% =====================================================================
%  R. Nielsen, GAMLE OG NYE PROPHETER.  Sex Forelæsninger.
%  Kjøbenhavn: Gyldendalske Boghandels Forlag (F. Hegel & Søn); trykt hos J. Jørgensen & Co., 1878.  112 pp.  Antiqua.
%
%  TRANSCRIPTION -- GENERATED; DO NOT EDIT BY HAND.  Rebuild from the repo root:
%    python3 ebook-method/gamle-nye-profeter/draft.py --emit --force \
%      && python3 ebook-method/collate/check.py gamle-nye-profeter --italics --apply \
%      && python3 ebook-method/gamle-nye-profeter/draft.py --post
%
%  STATUS: COMPLETE (2026-10-12): title page, Indhold and the whole text, pp. 1-112, with the one footnote (p. 105).
%
%  EDITION.  The text of record is the Royal Library's scan (KB alma99122156072605763; ~/bibliotek/Nielsen, Rasmus/
%  gamle-nye-profeter.pdf, SCANS.tsv), 250 ppi colour, with KB's OCR as text layer.  The only printing; no e-book and no
%  earlier transcription exists (checked 2026-10-12: the repo, the web, KB).
%
%  ACCURACY (2026-10-12).  The words are those of KB's layer (draft.py), collated word by word with an independent
%  tesseract reading of the images (collate/check.py, witness 'ocr'): 518 disagreements in 28,097 words.  35 were settled
%  by check.py's own rules and 119 by rules.py: accents tesseract adds, the printer's space before ;:?!, and six
%  candidate rules first measured on a BLIND PILOT of 74 crops (reader A, 16/16 controls): the print sets „...“ only
%  (tesseract's » and % are its „ and “) 12/12 and 12/12; tesseract's glyph faults where the layer's word is attested
%  11/11; the layer's equal-length glyph faults where tesseract's is attested 11/11; tesseract's aa read an/ag 5/5.  The
%  rest were read at the image BLIND (sheets.py: 352 word crops with two planted controls a sheet, 72 of 72 read right --
%  two read short where the crop cut the line; five Sonnet readers, .parts/reads/r-b1-B..F.tsv); 85 of those were
%  decided in the main session from the readers' lines and line crops (zoom.py: sheets za, zb).  LAYER_SPLITS opened five
%  words the layer split and the collation had passed („Ny rationali sme“, „F orsøg“).  The layer reads V as a backslash
%  (\iser, \idens): draft.py maps it.  LETTERSPACING (the print's emphasis, here \emph): candidates by box width
%  (spacedwidth.py: a per-letter width fitted over tesseract's boxes; spaced words run 1.2-1.3 times their predicted
%  width), 69 lines read blind (l-sheets), 55 runs set.  No italics (slant 99 % <= 1.05, no runs).
%  Sweeps after OPEN 0, each spot read at the image: check.py --moved (nothing real), the readers' whole lines against
%  the text (5 differences, all the readers'), paragraph starts against tesseract's indents (none), and every word in no
%  other transcription (542 words; zoom zc: two pairs of words run together, three line-end compounds whose hyphen is
%  kept, misprints below).  SAMPLED PAGE CHECK (pagecheck.py, seed 202610122): pp. 11, 14, 94 and the title page read
%  whole and blind, 793 words: 0 word, 0 punctuation and 0 markup errors (95 % upper bound 3.8 per 1000 words).
%  TRANSLATION as a last reading: the six translators listed 47 odd spots (ebook-method/gamle-nye-profeter/tr/odd-*.tsv),
%  all read at the image (zoom zo-01..03, zq): 21 mended by POST patch -- specks both OCRs read as stops, commas,
%  hyphens or apostrophes (pp. 4, 6, 25, 45, 52, 56, 79), „ read as a comma (p. 30), a question mark read as a colon
%  (p. 31), „Gud er“ run together (p. 84), the pronoun I read as 1 (p. 51), an elided compound „Phantasie- men“ (p. 29),
%  a lost paragraph break (p. 92), two damaged sorts (p. 15, 16) -- the rest are the print's (Atmosphære. der, Nebukadnezer,
%  Schwarz/Schwartz, Øjeblik, vanheldet, quotation marks the print leaves open or never opens).
%  What this cannot see: a misprint both OCRs normalise into a real word -- the corpus sweep finds only non-words.
%  PRINTED AS IS (each seen at the image; a % comment at the site): „prostesterer“ (p. 66), „Selstændighed“ (p. 56),
%  „Naturnødvendighen“ (p. 71), „alfor“ (p. 86); also, unmarked, the print's „Père Enfentin“, „Georger“ (p. 74),
%  „Schwarz“ (p. 108) against the note's „Schwartz“ (p. 105).
%  PRINTER'S DEFECTS (a sort that did not ink or broke; the letter supplied, a % comment at the site): „kunne“ (p. 16,
%  blank before „unne“), „Propheter“ (p. 16, „Prophe er“), „Forening“ (p. 15, „Fo: ening“).
%  The ø rule (Hans 2026-10-08) found nothing beyond what the collation decided.
%  Not transcribed: signature marks, library stamps, specks; the large initials of the lecture openings are set as plain
%  letters; the rule under each lecture heading is drawn by \lecture.  Quotation marks are as printed („ “).
%
%  PAGE MAP.  printed p = PDF page - 11 (p. 1 = PDF 12, p. 112 = PDF 123); half-title PDF 6, title PDF 8, Indhold PDF 10.
%  Each lecture begins on a new page: I p. 1, II p. 20, III p. 39, IV p. 58, V p. 78, VI p. 96.
%  The Indhold sets »Første Forelæsning« etc. letterspaced (here \emph, house convention) and has no page numbers; its
%  two section heads »Gamle Propheter.« and »Nye Propheter.« (bold sans) stand only there, not in the text.
%
%  TOOLS (ebook-method/gamle-nye-profeter/, copied from phantasiens-magt): draft.py, ocr.py, rules.py, sheets.py +
%  reads.py + READER-BRIEF.md, zoom.py, prescreen.py, pagecheck.py + PAGECHECK-BRIEF.md, patches.py (PATCHES, POST),
%  decisions.tsv, spacedwidth.py + spaced.tsv, and the sweeps splitsweep.py, linecheck.py, parcheck.py, unknown.py.
%
% =====================================================================
\documentclass[12pt,a4paper,openany]{book}
\usepackage[utf8]{inputenc}
\usepackage[T1]{fontenc}
\usepackage{libertinus}
\usepackage{libertinust1math}
\usepackage[danish]{babel}
\usepackage[protrusion=true,expansion=true]{microtype}
\usepackage{setspace}
\usepackage[margin=1.3in]{geometry}
\usepackage{fancyhdr}
\usepackage{hyperref}
\hypersetup{pdftitle={Gamle og nye Propheter}, pdfauthor={Rasmus Nielsen}, hidelinks}
\pagestyle{fancy}
\fancyhf{}
\fancyhead[LE,RO]{\thepage}
\fancyhead[C]{\textit{R. Nielsen: Gamle og nye Propheter (1878)}}
\setstretch{1.3}
\usepackage{marginnote}
\renewcommand{\thefootnote}{*)}
\newcommand{\tocline}[3]{\noindent\makebox[2.2em][r]{#1}\hspace{0.6em}#2\dotfill\ #3\par}
\newcommand{\opage}[1]{\marginnote{\footnotesize\textit{[#1]}}}
% a lecture opens a new page: heading, a short rule; the print's large initial is not reproduced
\newcommand{\lecture}[1]{\clearpage\begin{center}{\large #1}\\[0.4em]\rule{3em}{0.4pt}\end{center}\medskip\noindent}

\begin{document}

% --- front matter: PDF 8, title (read at the image, 2026-10-12); PDF 6 is a half-title »GAMLE OG NYE PROPHETER.«
\begin{titlepage}
\begin{center}
\vspace*{3em}
{\Huge GAMLE OG NYE PROPHETER.}\\[2.5em]
SEX FORELÆSNINGER\\[1em]
{\small AF}\\[1em]
{\large R. NIELSEN.}\\[2em]
[publisher's device]\\[4em]
KJØBENHAVN.\\[0.5em]
GYLDENDALSKE BOGHANDELS FORLAG (F. HEGEL \& SØN).\\[0.5em]
{\small TRYKT HOS J. JØRGENSEN \& CO.}\\[0.5em]
{\small 1878.}
\end{center}
\end{titlepage}

% --- PDF 10: Indhold (read at the image, 2026-10-12); no page numbers in the print
\begin{center}{\large Indhold.}\\[0.3em]\rule{3em}{0.4pt}\end{center}
\noindent\emph{Første Forelæsning:} Propheternes Kald og Stilling i Fortid\par
\noindent\hspace*{2em}og Nutid.\par
\begin{center}\textbf{Gamle Propheter.}\end{center}
\noindent\emph{Anden Forelæsning:} Moses, Samuel, Elias.\par
\noindent\emph{Tredie Forelæsning:} Prophetiske Fremsyn; Messiasbilleder.\par
\noindent\emph{Fjerde Forelæsning:} Spaadommenes Opfyldelse; Brydninger\par
\noindent\hspace*{2em}mellem det Gamle og det Nye.\par
\begin{center}\textbf{Nye Propheter.}\end{center}
\noindent\emph{Femte Forelæsning:} St. Simon, Fourier, Proudhon: For-\par
\noindent\hspace*{2em}negtelse af Gud og Forsyn; Nihilisme.\par
\noindent\emph{Sjette Forelæsning:} Opgjørelse; Nyrationalisme.\enspace Slutning.\par
\clearpage


% ===== TEXT (pp. 1--112) =====
% --- p. 1 ---
\setcounter{footnote}{0}%
\lecture{Første Forelæsning.}\opage{1}Som Menneskehedens store Lærere, Folkenes Førere
fremad mod Haabets forjættede Land, som Aandens
Fakkelbærere paa dunkle, ubanede Veie indtage Propheterne
en høi Rang. Men --- spørger man --- ere disse
Førere nu ogsaa paalidelige? De gamle Propheter ere jo
døde, og Synerne hos de nye komme stundom her og der
i betænkelig Forvirring. Er det virkelig en Historiens
Lov, at enhver aandelig bevæget Tidsalder maa have
sine Propheter? En kritisk Sammenligning mellem Oldtid
og Nutid kunde paa Forhaand friste til at antage
det Modsatte.

Den gamle Verden, fattig paa Erfaringer, men rig
paa Anelser og geniale Indskydelser, famlede i et udvortes
Halvmørke. I Mangel af sikkre Love og tilforladelige
Kjendsgjerninger maatte de indre Lys med
Syner og Aabenbaringer være dens ypperste Ledestjerner:
den gamle Verden trængte til Propheter. Men
den nye Verden med sit væbnede Øie, sine Kikkerter
og Mikroskoper gaaer frem ad Beregningens Vei og
kræver en anden Veiledning. I Politik, i Statshuusholdning,
i Industri og Handel adspørger man ikke
Orakler; man søger i en Pengekrise ikke Raad hos
Propheter, men hos kyndige Fagmænd.

% --- p. 2 ---
\setcounter{footnote}{0}%
\opage{2}Propheten er som en Seer paa Bjerget. Fjernt i
Horizonten dukke Fremtidsbilleder op af Mulighedernes
Taage og tegne sig i svævende Omrids paa Skyernes
luftige Grund; det er Straalebrydninger af det indre
Lys, der give disse Skikkelser Liv og Farve og fuld
Nærværelse. Her er intet tilforladeligt Afstandsmaal;
Propheten seer det Tilkommende som et Nærværende:
han regner ikke, han skuer. Den ræsonnerende Forsker
kommer derimod stykkeviis til sine Erkjendelser. For
ham er det Tilkommende en, ifølge Tingenes Sammenhæng,
lovbestemt Følge af det Forbigangne; han stoler
ikke paa en blot anskuende Phantasi; han tvivler og
prøver, han eftertænker og regner. Det Sidste forekommer
Nutidens Oplyste at være nok saa sikkert som
det Første. Vel sandt, vi kunne ikke forud beregne
en politisk Omvæltning som Astronomerne beregne en
Solformørkelse; uforudseelige Omstændigheder ville altid
spille en Rolle i de historiske Begivenheder; den omhyggeligst
beregnede Fremtid har gjerne en eller anden
Overraskelse i Baghaanden. Men kunne vi end ikke
nøiagtig forudsige, hvad der nødvendig maa skee, saa
kunne vi dog ved Beregning udfinde, hvad der efter al
Sandsynlighed vil skee. Sandsynlighedsregningen, der
i vore Dage har vundet et udstrakt Omraade, viser sin
Overlegenhed paa de allerforskjelligste Felter. Har
man om et Lands Befolkning et tilstrækkeligt Antal
statistiske Opgivelser for en længere Aarrække, saa
kan man, uden prophetisk Inspiration, med temmelig
Sikkerhed udregne, hvormange Dødsfald der ville indtræffe
i det kommende Aar, hvormange Selvmord, hvormange
der ville hænge sig, hvormange der ville drukne
sig, skyde sig osv. Hvad Ægteskaber angaaer, udregner
% --- p. 3 ---
\setcounter{footnote}{0}%
\opage{3}man ikke blot, hvormange der i Aarets Løb ville gifte
sig, men endog, hvormange med unge Piger, hvormange
med ældre Piger og hvormange med Enker. Saa vidt
er man naaet, at Postvæsenet kan beregne, hvormange
Breve der imellem Aar og Dag ville blive afleverede
med forfeilet Udskrift og hvormange uden Udskrift,
fordi man i sit Hastværk har glemt Udskriften. Sandsynlighedsregningen
er i stadig Fremgang; de statistiske
Tabeller voxe saa vel i Alsidighed som i Omfang og
ville, for de Dannedes Vedkommende, vel snart omtrent
kunne siges at have afløst Prophetierne.

Denne Betragtning er vildledende, thi den forrykker
Synspunktet.

I den udvortes Verden, ja selv i den indvortes, er
der Meget, der kan beregnes, og Enhver, der i slige
Ting vil gjætte sig frem eller stole paa Anelser eller
vælge iblinde, maa tage Følgerne af, at han af Magelighed
har forsømt at tænke og regne. Men med
Aandens høieste Anliggender er det noget Andet.
Skulde der komme en Tid, i hvilken det Prophetiske
ingen Betydning havde, da maatte det vel være en Tid,
hvori det i dybere Forstand Poetiske heller ingen Betydning
havde, en Tid, hvor Phantasie og Følelse, Hjertets
Rørelser og ideale Længsler ingen Betydning havde,
en Tid, hvori Sjælen, Samvittigheden med alt det Personlige
satte Kræfterne til paa Erhvervelse af endelige
Ting og Beregning af Nytte og Fordeel. I en saadan
Tid vilde ganske vist Religionen være en Indbildning,
Troen Overtro, Aandens Majestæt dethroniseret. Men
en saadan Tid er endnu ikke kommen og vil forhaabentlig
aldrig komme. 

% --- p. 4 ---
\setcounter{footnote}{0}%
\opage{4}Den forberedes dog, siger man. Længe nok har
den gamle østerlandske Fabel om en mirakelgjørende
Aands vilkaarlige Magt tyranniseret den menneskelige
Natur; det er paa Tide, at denne Natur faaer Lov at
røre sig frit og skabe sig selv sine egne Tilværelsesformer.
Aanden er ikke et Væsen for sig, ophøiet
over Naturen; det er Naturen selv, som i Mennesket
kommer til Bevidsthed og bliver til Aand. Alle menneskelige
Forhold, private saa vel som offentlige, bestemmes
ved naturlige Krav og have deres sidste Grund i uafviselige
Naturforhold. Mandens Kjærlighed til Kvinden,
Moderens Kjærlighed til Barnet, Folkets Kjærlighed til
Fædrelandet, hvad er det Andet end lutter Naturbestemmelser?
Det Samme gjælder om Skjønhed og
Ære, ja om Sandhed og Ret. Den, der ikke har en
oprindelig, d. e. en naturlig Sands for det Skjønne,
faaer aldrig at vide, hvad Skjønhed er. Dersom et
Menneske af Naturen ingen Æresfølelse har, hvad nytter
det saa at ville indprænte ham Ærens Grundsætninger?
Ved passende Anvendelse af Straf og Belønning kan
man maaskee tvinge Uvorenheden ind i en foreskreven
Sædeligheds Former; man kan jo afrette det gjenstridige
Individ, som man afretter den vilde Hest, ved at lade
ham tørste. Men fremmer det Sædeligheden selv? Om
man ogsaa med det uhyre Apparat af en overnaturlig
Vilkaarligheds Vrede og Naade kunde stille ni Helveder
paa den ene Side og lige saa mange Himmeriger paa
den anden i Udsigt, man vilde dog ikke ved udvortes
Tryk formaae at presse sand Erkjendelse af Sandhed,
Frihed og Ret ind i en eneste Menneskesjæl.

Der er et uopløseligt Baand mellem Natur og
Sandhed. At fornegte, hvad der følger af en Sags
% --- p. 5 ---
\setcounter{footnote}{0}%
\opage{5}Natur, er at fornegte Sandheden; at holde sig strængt
til Sagens Natur er at holde sig til Sandheden. Der
gjøres Væsen af Synden; man græder høitidelig over
den faldne Natur. Men hvad er det, der frembringer
Synden og fordærver Naturen? Det er Vilkaarligheden.
Naturen i al Naturlighed kan ikke fordærve sig selv;
det er just det Vilkaarlige, som dukker op, idet Naturlivet
gjennem Mennesket skal gaae over til Aandsliv,
der frembringer Synden. Idet Vilkaarligheden miskjender
Naturens Krav og modarbeider dens Stræben,
fremkommer det Falske, det Unaturlige. Vilkaarligheden
er det Ondes Princip, den virkelige Aarsag til
Splid og Uorden. Hvad er Uorden vel andet end
Unatur? Man vil ved vilkaarlige Forskrifter tvinge
Personer af ulige Sind, modstridende Naturer, til vedblivende
at leve sammen i Ægteskab, fordi de nu engang
ere i Ægteskab, leve som om de af Hjertet vare enige,
fordi Vilkaarligheden har befalet, at de skulle være
enige, leve saa det seer ud, som om de inderlig elskede
hinanden, skjøndt de vitterlig ere nær ved at hade hinanden
til Døden; denne Usandhed, denne sjælefortærende
Unatur skal --- saa lyder Vilkaarlighedens Bud --- være
det rette Tegn paa sædelig Orden, det ægte Beviis paa
Ægteskabets Hellighed. Hvad er det, som bringer
Samfundet i Uorden? Dette Unaturlige, at Noget, der
efter sit Væsen er af lavere Art, bliver sat som et
Høiere, medens det af Naturen Høiere holdes under et
Tryk, der stempler det til et Lavere; dette Unaturlige,
at Mange maae lide, for at nogle Faa kunne nyde,
at den Ene maa gjøre Arbeidet, for at den Anden kan
høste Lønnen; dette Unaturlige, at Folkefærd af een
Rod og eet Blod rives fra hinanden med vilkaarlig
% --- p. 6 ---
\setcounter{footnote}{0}%
\opage{6}Magt, medens Folk af forskjellig Afstamning med forskjelligt
Sprog, forskjellige Sæder og forskjellige Udviklingsbetingelser
skulle tvinges til at opgive, hvad
der udgjør deres Livs Naturgrund, og leve under en
Alt nivellerende Ensformighed i samme Stat under
samme Love, alene fordi den mægtige Vilkaarlighed
vil det.

Det er mærkeligt at see, hvorledes Vilkaarligheden,
den geistlige saa vel som den verdslige, sveder og arbeider
for at holde Styr paa Millionernes Ustyrlige og
sikkre det borgerlige Samfund Orden og Ro. Det vil
ikke lykkes: Synden, ak den slemme Synd! Synden
tager til. Troen tager af, den fordærvede Natur bliver
værre og værre. Ganske naturligt; saa længe man arbeider
med Vilkaarlighedens Princip, vilkaarlig Løn og
vilkaarlig Straf, overnaturlig Vrede og overnaturlig
Naade, skal Synden, det Unaturlige, nok trives. Gaaer
til Sagens Natur, skaffer os Lys over Lovenes Grund,
naturlige Forhold i Samfundets Indretning, og Uroen
vil høre op og Synden tage af; saa kommer den nye
Tid: saaledes som det nu er, kan det umulig blive.

Den eiendommelige Blanding af Sandt og Falsk,
som ligger til Grund for denne Udtalelse, skulle vi her
ikke nærmere drøfte. Kun det skal fremhæves, at Udtalelser
af denne Natur grunde sig paa Mere end Sandsynlighedsregning
og statistiske Tabeller. Det er prophetiske
Røster, vi høre; de nye Propheter stille sig
frem for at afløse de gamle.

Enhver Tid, hvori høiere Livskræfter røre sig og
store Aandsspørgsmaal ere i Bevægelse, maa, under en
eller anden Form, nødvendig have sine Propheter. Det,
som bestemmer en Tidsalders Karakteer, er den herskende
% --- p. 7 ---
\setcounter{footnote}{0}%
\opage{7}Verdensanskuelse. Vel kunne forskjelligartede Verdensanskuelser
være samtidige og brydes med hverandre;
men, med mindre Alt er i chaotisk Opløsning, altid
vil der være Grundtræk af en Anskuelse, som har
Fremtiden for sig. Vor kirkelige Forestilling om Himmel
og Jord tilhører Oldtiden, den videnskabelige Nutiden.
Forskjellen mellem Oldtid og Nutid er i denne Henseende
saa gjennemgaaende, at alle Tegn tyde hen paa
en radikal Omdannelse. Her er et Thema for stemningsrige
Variationer. Det hører til Prophetiens Væsen
at anslaae Tilværelsens Grundakkorder og røre de Strenge,
som dybest vibrere i Slægternes Øre. Idet Propheterne
staae paa fremragende Høider, tage de sig ud som eenlige
Skikkelser, men see vi nærmere til, ere de ingenlunde
eenlige; der er Tanker, der gjøre dem Selskab, de have
deres Liv i Kræfternes Midtpunkt. Hvad der gjærer i
Slægten, dens Trang og Længsel, dens Smerte og Haab,
bevæger sig kraftig i dem og klarer sig for deres Syn;
det Anede, det endnu Uforklarede, det hvorefter de Mange
famle, det i Sindet Bundne, som ikke har fundet det
løsende Ord, det hente de frem af Dybets Nat og løfte
det op i Daglyset.

Det Overnaturliges Forhold til det Naturlige, Frihedens
Forhold til Nødvendigheden, Verdensstyrerens
Villie og Begivenhedernes Sammenhæng, disse Aandens
udødelige Emner holde Bevægelsen i Gang og udgjøre
det stadige Mellemværende mellem gamle og nye Propheter.


Oversee vi det mythisk Forblommede, hvori Naturligt
og Overnaturligt sammenblandes --- Oraklerne. Sibyllerne,
Vølverne ---, og holde os til den skarpt udprægede
Modsætning af Aand og Natur, saa finde vi
% --- p. 8 ---
\setcounter{footnote}{0}%
\opage{8}det for Oldtiden Typiske hos Israels gamle Propheter.
Jehova, Aandens Gud, er Verdens Skaber, hans Villie
almægtig. Aandens Væren og Virken er altsaa ikke
afhængig af endelige Naturbetingelser; den almægtige
Aand virker, naar den vil, hvad den vil og som den
vil, med ubetinget Frihed. --- Med en saadan vilkaarlig
Brug af Almagtskræfterne kunne de nye Propheter paa
ingen Maade forsone sig. En mirakelgjørende Villie,
der skalter og valter med Naturen efter Behag og leger
olympisk Leg med Aarsager og Love, finde de i den
Grad stridende mod Erfaring og sund Fornuft, at de i
Forargelse herover ofte gaae til den modsatte Yderlighed
og fornegte Aandens oprindelige Frihed. Det er
da ikke Aandens frie Villie, men den evige Naturs
iboende Nødvendighed, der er det alt bestemmende
Første. Hvad Naturen fordrer, maa Aanden ville: den
guddommelige Villie er, ifølge denne Synsmaade, lige
saa lidt i Stand til at forandre Naturens Love, Naturens
Løb, som den menneskelige.

Denne Modsætning forstærkes end yderligere ved
den Maade, hvorpaa Aandsmeddelelsen antages at
foregaae.

De gamle Propheter modtoge Aanden umiddelbart
som Jehovas Sendebud; Aanden kom over dem med
Jehovas Ord, forkyndende Jehovas Villie. Kaldelsen
til Prophet var et Aandens Mirakel. Forholdet mellem
den guddommelige Aand og den menneskelige Bevidsthed
vedblev desaarsag at være mirakuløst. Propheten
veed at skjelne mellem sine egne Tanker og de Tanker,
der indgives ham af Aanden, mellem sine egne Ord til
Jehova og Jehovas Ord til ham. Hvor omstændeligt
og med hvilken Tydelighed Aanden kan forudsige Pro%
% --- p. 9 ---
\setcounter{footnote}{0}%
\opage{9}pheten Ting, som intet Menneske af sig selv kunde vide,
derpaa haves mangfoldige Exempler. Da Samuel (1 Sam.
c. 10) havde salvet Saul til Konge over Israel, kyssede
han ham og talede saaledes: „Naar du idag gaaer fra
mig, skal du støde paa to Mænd ved Rachels Grav paa
Grændsen af Benjamins Lod i Zelzach; de skulle sige
til dig: Aseninderne, som du gik ud at lede efter, ere
fundne; men nu har din Fader slaaet sig den Sag med
Aseninderne af Hovedet, han er bekymret for Eder og
siger: Hvad skal jeg gjøre for min Søn? Naar du derfra
gaaer længer frem og kommer til Egen Thabor,
saa skal der møde dig tre Mænd, som drage op mod
Bethel til Gud, den ene af dem bærer tre Gedekid,
den anden tre Brød, den tredie en Bimpel Viin. De
skulle byde dig Hilsen og give dig to Brød, dem skal
du modtage af deres Haand. Derefter skal du komme
til Guds Høi, hvor Philisterstøtten staaer; naar du
kommer der til Byen, skal du møde en Skare Propheter,
som komme ned fra Høien med Strængespil,
Pauker, Fløiter og Harper foran sig, og som prophetere.
Da skal Jehovas Aand komme heftig over dig; du skal
prophetere tillige med dem og blive en heel anden
Mand end før. Naar nu disse Tegn ere indtrufne, saa
gjør, hvad du føler Drift til, thi Gud er med dig“.
Og det skete nøiagtig saaledes. I det Øieblik Saul
„vendte Samuel Ryggen, vendte Gud hans Hjerte, saa
det blev et andet. Og alle disse Tegn indtraf paa
denne Dag. Da han naaede Høien, kom en Skare Propheter
ham imøde, og Guds Aand kom heftig over ham,
og han propheterede midt iblandt dem. Alle de, som
kjendte ham fra tidligere Tid og saae, at han propheterede
tilligemed Propheterne, sagde til hinanden ind%
% --- p. 10 ---
\setcounter{footnote}{0}%
\opage{10}byrdes, hvad gaaer der af Kis' Søn, er Saul ogsaa
iblandt Propheterne ?“

En saadan Aandsmeddelelse, en saa smaalig nøiagtig
Forudsigelse af mødende Tilfælde kan --- ifølge
de nye Propheter --- kun finde Tiltro hos dem, der
føle sig oplagte til al Slags Tro, ogsaa til Tro paa
Clairvoyance og Spiritisme. Ingen, som har noget
Begreb om Tilværelseslove, aandelige saa vel som
naturlige, vil kunne gjøre sig nogen Forestilling om,
hvorledes en Aand, som ikke er vor Aand, kan træde
talende op i vort Indre og underrette os om, hvormange
Mænd med Gedekid og Brød og Viinbimpler, vi skulle
møde paa Veien. --- Gjøre vi disse Viismænd opmærksomme
paa, at der dog vist er Mange, og det endog
blandt de Oplyste, som endnu oprigtig troe paa Samuels
prophetiske Fremsyn, saa svare de os, at det skeer ved
et fromt underfundigt Selvbedrag. Det gaaer med prophetiske
Fremtidsbilleder som med naturlige Lysbilleder:
en lille Dreining af Speilfladen kan forandre det Hele.
Samuels Sind er fyldt af Tro paa Jehova; Aanden
rører sig i ham ved Tanken om det afgjørende Skridt
at skulle vælge Folket en Konge. Saa taler han til
Saul med en Myndighed, som hans Alder og Stilling
og Aandsoverlegenhed giver ham. Saul skal ikke længer
lede om sin Faders Aseninder; han skal finde Propheterne
og blive en Aandens Mand; thi kun ved at blive
en heel anden Mand end før, „vil han som Konge kunne
styre Israel“. Det virker paa den begavede unge Mand,
Ordet slaaer ned i hans Sjæl som Lyn; saaledes finde
vi ham beaandet i de Beaandedes Kreds, --- Saul iblandt
Propheterne. Det er forstaaeligt, thi det er naturligt.
Idet man, halvt bevidst, halvt ubevidst, underskyder
% --- p. 11 ---
\setcounter{footnote}{0}%
\opage{11}en saadan Forklaring, medens man paa Overleveringens
Autoritet og fordi det Overnaturlige altid glimrer i
Phantasien holder sig bogstavelig til Fortællingen, eller,
som det saa koket indsmigrende hedder, annammer det
aabenbarede Ord i al Eenfoldighed, kan man jo let
bilde sig selv og Andre ind, at man endnu i Vantroens
Tid er En af de gammeldags Troende. Men en Taagesky
af Tvetydigheder ruger paa Bunden af saadanne
Sjæle. Dersom Selvbesindelsen vaktes, og Eftertankens
Lys fik Lov til at bryde igjennem, vilde de gammeldags
Troende, sikkerlig til Overraskelse for dem selv,
komme til at gjøre moderne Opdagelser. Den bogstavtroende
Opfattelse og den underskudte Forstaaelse tilhøre
to aldeles modstridende Anskuelser, Oldtidens og
Nutidens.

Efter den oprindelige Mening har Jehovas Aand
vilkaarlig, d. e. uden nogetsomhelst psychologisk Hensyn
til Naturbetingelser, indgivet Samuel hvert Ord i
den prophetiske Forudsigelse, hvad da ogsaa bekræftes
derved, at Tegnene punktlig indtræffe. Men ikke at
tale om, at en saa naturløs, vilkaarlig Aandsmekanik
staaer i skjærende Strid med de indre saa vel som
ydre Erfaringer, hvori ethvert Nutidsmenneske aander
og lever, saa er heri tillige en Modsigelse, som selv
den mest ubetingede Mirakeltro ikke formaaer at
fjerne. Naar Jehova er alvidende, og den Alvidende
nu en Gang i al Vilkaarlighed har besluttet at meddele
et Menneske, hvad der vil skee i den nærmeste Fremtid,
saa kan man under denne Forudsætning finde det
i sin Orden, at Forudsigelsen indtil de mindste Enkeltheder
gaaer i Opfyldelse. Men netop ved at holde
Forudsætningen fast drives vi af en ubøielig Logik
% --- p. 12 ---
\setcounter{footnote}{0}%
\opage{12}videre fremad til det Spørgsmaal: hvorledes er det muligt,
at den alvidende Aand, som her saa nøiagtig har truffet
alle Biting, har forfeilet Hovedtingen? Samuel er i
Aanden forvisset om, at Saul er den rette Mand, men
det er jo, hvad den følgende Historie viser, en stor
Feiltagelse. Saul er netop ikke den rette Mand. Er
det nu Jehova, der har taget feil, eller er det Samuel
eller maaskee dem begge? Hvordan man end vender
og dreier sig, kommer man ind i en Labyrinth af Modsigelser.
Vidste den Alvidende, at Saul virkelig var
den værdigste, da han valgtes til Konge, hvorfor vidste
den Alvidende da ikke ogsaa, at han senere vilde gjøre
sig uværdig? Siger man: den Alvidende vidste meget
godt, at Saul vilde svigte sit Kald, men valgte ham
alligevel, fordi han vilde vælge ham, saa have vi den
rene Vilkaarlighed. Siger man: Sauls Ulydighed var
en Yttring af hans frie Villie, og hvad der i en nær
Fremtid vil skee i Kraft af Menneskenes frie Villie,
kan end ikke den Alvidende selv forud vide; saa kommer
man i Forlegenhed med Gedekiddene, Brødene og
Viinbimplen; med mindre man vil antage, at Alt var
bestilt forud, og at de Mænd, Saul mødte, ikke mødte
ham med, hvad de bare, af fri Villie, men ifølge høiere
Ordre. Hvis det Sidste antages, hvad er saa Samuels
hele Forudsigelse? Et aftalt Spil, en Prophetiens
Komedie. Til at komme ud af Vilkaarlighedernes
Labyrinth er der kun een Udvei, at tage Naturbetingelserne
med og holde sig til den underskudte
Forklaring.

Ved Naturbetingelsernes Indtræden gaaer Vilkaarligheden
bort, og Aandsmeddelelsen viser sig i et andet
Lys. Alle Aandsgaver ere, oprindelig seet, Naturgaver;
% --- p. 13 ---
\setcounter{footnote}{0}%
\opage{13}det er kun i de Begavede, de af Naturen særlig Anlagte,
med eet Ord: Genierne, at de store Verdenstanker, vi
kalde Ideer, forklares til Aand. Ideerne komme ikke
ved overnaturlige mirakuløse Indskydelser, endnu mindre
falde de talende ind i den beaandede Sjæl! Geniet
modtager sine Indtryk fra Omverdenen som andre
Mennesker, Forskjellen er kun den, at han ifølge en
egen organisk Oplagthed bearbeider sine Indtryk paa
en eiendommelig Maade, omsmelter dem i Erindringen,
samler dem i Anskuelsen, former dem, klarer dem og
oversætter dem i sit eiendommelige, dog almeenforstaaelige
Sprog.

Den indre Omdannelse skeer ofte paa saa uvilkaarlig
Maade, at den Begavede kan falde i Forundring
over sig selv og spørge: hvad er det dog, der foregaaer
i mig? Det arbeider herinde uden at jeg arbeider, det
tænker, selv naar jeg ikke tænker; det groer som
Græsset groer, det voxer som Fosteret i den Frugtsommelige.
Hvad jeg bringer frem er ikke af mig selv;
min Virksomhed, om jeg end anstrænger mig til det
Yderste, er ganske underordnet, jeg er kun en Tjener,
en Fødselshjælper. Saaledes kan det vel forekomme
ham, at de nye Syner med de nye Tanker ere gaadefulde
Indskydelser. Men det er ikke Indskydelser af
en overnaturlig Aand, som taler ind i Sjælen; det er
Sjælekræfternes eget Spil, den geniale Natur, i hvis
hemmelige Værksted der arbeides efter Love. Den Begavede
kan umulig gjøre Skilsmisse mellem det, som
det Gaadefulde virker i ham, og det, han selv bevirker.
Ved at forøge en saadan Adskillelse og antage det
Naturlige for et Overnaturligt, vilde han synke ned
under en naturløs Aandsmekanisme og ophøre at være
% --- p. 14 ---
\setcounter{footnote}{0}%
\opage{14}Genie. Geniet er netop Eenhed af Natur og Bevidsthed,
den lykkelige Eenhed, hvori Naturvirksomhed og
Selvvirksomhed gaae sammen til Aand. Geniets Ord er
et Menneskeord; men dette Menneskeord, som Naturtanken
har født og Ideen holdt over Daaben, faaer ved
Indvielsen en magisk Kraft, der vidunderlig griber og
bevæger vort Sind. Dette Menneskeord er Aandens
Ord, et prophetisk Ord, med Afglands af det Guddommelige.


Altsaa: Oldtidsherolderne ere Propheter af overnaturlig
Aand, Nutidsherolderne ere Propheter af Genie;
her staae vi ved Knudepunktet. Skjøndt Meget af det,
de nye Propheter have at sige os om Naturlovene, vel
fortjener at lægges paa Sinde, saa er der dog Noget,
som de deels skjule, deels oversee. De hævde Lovenes
Gyldighed, og vi maae give dem Ret; hvor der ingen
Love er, eller hvor Lovene vilkaarlig brydes, er der
ingen Sammenhæng, ingen Mening eller Forstaaelse.
Men de tale, som om der ingen andre Love var til end
Naturlovene, som om Aanden ikke ogsaa havde sine
Love. Det seer bestandig ud, som skulde de physiske
Nødvendighedslove overalt træde i Aandslovenes Sted.
Her maae vi være paa vor Post; thi, dersom den antydede
Grundsætning skal gjennemføres, saa har det
meget omfattende Følger. Deraf følger nemlig ikke
blot, at den guddommelige Aand umulig kan være
noget selvstændigt, fra Naturen forskjelligt Væsen, men
deraf følger ogsaa, at vor Sjæl, vort Selv, vor Villie
heller ikke kan være noget fra vor legemlige Organisme
forskjelligt Væsen. Er det da virkelig blevet beviist,
at Aandslovenes Selvgyldighed umulig kan bestaae med
Naturlovene? Ingenlunde! Den hele Beviisførelse ind%
% --- p. 15 ---
\setcounter{footnote}{0}%
\opage{15}skrænker sig til en slet og ret Paastand. Det maa
indrømmes, at vor Indsigt i Aandslovene er for ufuldkommen
til, at vi i det Enkelte skulde kunne gjøre tilstrækkeligt
Regnskab for Vexelvirkningen mellem Natur
og Aand. Men, spørge vi, kunne Physik-Philosopherne,
vi mene hermed saadanne kritiske Halvtænkere, som
leve af at kopiere den forskende Physiks Resultater,
da virkelig fra deres Side i alle Enkeltheder --- ved
Hjælp af Naturlovene gjøre Rede for, hvorledes det
Uorganiske bliver til et Organisk, Organismen til et
levende Legeme, det levende Legeme til Noget, der
tager sig ud som en Sjæl, et Selvvæsen, en bevidst
selvbestemmende Villie? Kunne de nye Propheter maaskee
forklare Geniet? De ville sikkerlig komme til at
fatte sig kort og lade det beroe ved den kategoriske
Erklæring, at det Overordentlige i Geniet er Naturlovene
underlagt. Med samme Ret kunne vi da paastaae,
at det Overnaturlige i de gamle Propheters
Inspiration er Aandslovene underlagt. Eller er det
maaskee blevet beviist, at det Overnaturlige er det
Naturløse, det Unaturlige? Nei, her er intet Beviis,
her er kun en Paastand.

Dersom vi ikke kunde kjende Geniet paa dets Frembringelser
og heller ikke tilegne os disse Frembringelser,
glæde os over dem og benytte dem, før vi tilbunds
havde udgrundet, hvad Geniet er, og til hvilken mekanisk
eller kemisk Forening af Stoffer det bliver at henføre, % printer's defect: the r broken („Fo: ening“)
saa maatte vi jo indtil videre give Afkald paa
Kunst og Poesi; vi maatte betragte Civilisationens Heroer
som mistænkelige, gaadefulde Størrelser, ja, hvad mere
er, vi vilde slet ikke være i Stand til at forstaae de nye
Propheter. Men saa sandt vi kunne skatte og beundre
% --- p. 16 ---
\setcounter{footnote}{0}%
\opage{16}Geniets Frembringelser uden at kunne begribe den
geniale Inspiration, hvi skulde vi da ikke anerkjende
de vidunderlige Frembringelser af Aabenbaringens Aand,
skjøndt vi ikke formaae at forklare den overnaturlige
Inspiration. Det forholder sig ikke ganske, som den
moderne Bevidstheds ivrigste Prædikanter høirøstet forsikkre
paa Gader og Torve --- især paa Physik-Philosophernes
Halmtorv ---, at Aandens Selvstændighed er
lutter Vilkaarlighed, at denne Vilkaarlighed er Skyld i
Dualismen, og Dualismen Skyld i vor aandelige Ulykke.
Det er neppe sikkert, at vi blive Livsgaaderne kvit,
fordi vi blive Vilkaarligheden kvit, thi Aandens ubetingede
Frihed er dog noget Andet end lovløs Vilkaarlighed;
vi tør ikke uden videre stole paa, at Døden kvitterer
Frihedens Regnskab, naar vi blot fornegte Friheden
og sove rolig ind i Nødvendigheden. Nødvendigheden er
ikke mere forstaaelig end Friheden. Enhver, som føler
Kald til at anstille Undersøgelser, vil overbevise sig
om, at Nødvendighedsproblemet er uløseligt, netop fordi
Frihedsproblemet er uløseligt, og omvendt; med andre
Ord: at Geniet i sidste Instans er lige saa uudgrundeligt
som Aanden.

Nødvendigheden er Naturens Uudgrundelige; Aandens
Uudgrundelige er Friheden. Den oprindelige Eenhed
af Frihed og Nødvendighed er skjult i den evige
Viisdom. Hvad enten vi nu nærmest føle os tiltalte af
Aandens eller af Geniets Propheter, maae vi stole paa % printer's defect: the t did not ink („Prophe er“)
Viisdommen --- paa Daarskaben kunne vi jo ikke stole; % printer's defect: the k did not ink (a blank before „unne“)
vi maae see paa Styrelsens Gang i Historien. Menneskeslægtens
Udvikling er en naturbetinget Aandsudvikling,
en ofte hæmmet, af mangfoldige Hindringer gjennemkrydset,
men i det Hele fremadskridende Frihedsudvikling.
% --- p. 17 ---
\setcounter{footnote}{0}%
\opage{17}Vilkaarlighed og Naturskranker høre til Udviklingens
elementære Forudsætninger. Ved de mødende Skranker
opstaaer en Følelse af Nødvendighed, ved Vilkaarlighedens
Brug en Følelse af Frihed. Livet kan lige saa
lidt undvære Nødvendigheden som Friheden, begge
hæve sig for det udvidede Syn som ideale Magter i
uophørlig Strid. Skal Friheden tjene Nødvendigheden
eller Nødvendigheden Friheden?

Tilstandene vexle. Bag Mythens glimrende Slør
begynder Nødvendigheden en storartet Leg, der sætter
Slægtens ungdommelige Kræfter i Bevægelse. Den mægtige
Natur tager Menneskeskikkelse paa; det er som
en Spøg og dog frygtelig Alvor. Paa Bjergenes Toppe,
paa Havets Bund, i Skovenes Tykning, i Floder og
Træer, overalt er der menneskelignende Væsener, der
leve, føle og ville. Elementernes Kamp er en Frihedskamp,
og Friheden seirer: Uhyrerne synke, Olympen
befolkes af salige Guder, og fra Afgrunden høres Giganternes
Suk. Det er en Heroernes Tid, en Geniernes
Tid med spiredygtige Tanker og store Handlinger, et
Civilisationens Foraar, da Evnerne skyde i Veiret, Naturen
indvies til Aand og Ideerne blomstre. Alt er
Skjønhed og Harmoni, thi Friheden er nødvendig og
Nødvendigheden fri. Men Herligheden varer kun kort;
den legende Nødvendighed tager Masken af, og hvad
viser sig? --- Skjæbnens forstenende Aasyn! Tilværelsen
mørknes, Guderne blegne, Frihedsmodet synker; Forstanden
staver paa Tingenes Sammenhæng, og jo længere
den kommer frem i Texten, desto tydeligere viser det
sig, at Friheden kun var en Drøm: det virkelige Liv
er Trældom.

% --- p. 18 ---
\setcounter{footnote}{0}%
\opage{18}Tilstandene vexle. For at Skjæbnenødvendigheden
ikke skal kvæle Friheden, har den evige Viisdom føiet
det saa, at Aanden kan komme til sin Ret. Det begynder
med Friheden forfra, men ikke, som før, i Tvetydighed.
Naturfriheden var en Tvetydighed og derfor
en Skuffelse. Tvetydigheden hæves, idet Aanden adskiller
sig fra Naturen og forkynder sig selv som Alskabningens
Herre. Forkyndelsen er i Miraklet. Det Naturlige lider,
Miraklet slaaer igjennem; for Miraklets Slag maa Nødvendighedstraadene
briste og Skjæbnen synke. Det er
Befrielsens Værk: Aandens almægtige Villie aander
paa de afmægtige Villier, stiller dem Frihedsmaalet for
Øie og vækker dem til fornyet Kamp. Det er Aabenbaringens
Tid, en Martyrernes Tid: det Timelige opoffres
for det Evige, det Overnaturlige faaer Magt med
det Naturlige.

Tilstandene vexle. Det overnaturlige Gjennembrud
virkede saa stærkt paa den menneskelige Phantasie, at
den væmmedes ved det Naturlige og sværmede kun for
det Overnaturlige. Naturen falmede; ved at betragtes
i det overnaturlige Lys tog den sig ud som en gammel
Skjøge, af hvis Ungdomsskjønhed der neppe var Spor
tilbage. Istedenfor at forstaae den udtalte Modsætning
af Aand og Natur som et Sandhedens Udtryk for dybere
Eenhed, søgte man Befrielsen i et vilkaarligt Brud paa
Naturens Love. Den virkelige Tingenes Orden gik
under i lutter Mirakler, Legender og hellige Eventyr.
Friheden blev miskjendt, fordi Nødvendigheden blev
miskjendt. Men Tilstandene vexle. Den evige Viisdom
har føiet det saa, at Naturen paa Ny kan komme
til sin Ret. Opdagelsernes Tid, Naturvidenskabernes
% --- p. 19 ---
\setcounter{footnote}{0}%
\opage{19}Tid, de frie Tankers og de frie Forskningers Tid er
atter Geniernes Tid, men --- paa høiere Vilkaar.

De modsatte Tilstande gaae, hver efter sin Retning,
til saa stærke, eensidige Yderligheder, at det af den
Grund er vanskeligt for Mennesker, der leve i den ene
Tilstand, ret at forstaae, hvad der oprindeligt er levet
i den anden. Men vi maae stole paa Viisdommen.
Kun den styrende Viisdom kan give os en Forstaaelsens
Nøgle, en Gaadens Løsning, en Maalestok til Bedømmelse
af \emph{gamle} og \emph{nye Propheter}.

% --- p. 20 ---
\setcounter{footnote}{0}%
\lecture{Anden Forelæsning.}\opage{20}Det er en fremmed, længst forsvunden Verden, Israels
gamle Rige, vi skulle gjøre Bekjendtskab med. Men,
indvender man, dette Rige kan da ikke siges at være
os fremmed. Det israelitiske Folks Historie læres jo
saa godt som i alle vore Skoler og kjendes selv af den
menige Mand. Gives der noget Oldtidsfolk, om hvis
Liv og Skjæbne, man tør forudsætte almindelig udbredt
Kundskab, saa er det nok det jødiske. --- Og dog er
hiint Rige os fremmed, mere fremmed end Grækernes
og Romernes Kultur, mere fremmed end Urtidens Tilstande
i Steenalderen, Broncealderen og Jernalderen,
hvis fattige Levninger udgraves af Jorden. Og hvorfor
saa fremmed? Fordi det, som skal karakterisere dets
inderste Væsen, er os uforstaaeligt.

Der maa være foregaaet en stor Forandring, ja en
Forvandling, med Himmel og Jord, siden hiint gamle
Vidunderrige forsvandt. En Forvandling? Den virkelige
Natur er dog i store Træk den samme. Det gamle
Ægypten med Nilens Flod, der synker og stiger, er til
endnu. Endnu staaer Sinai Bjerg, en Kæmpe af Granit,
urokkeligt paa sin brede Fod, skjøndt Aartusinders
Storme have hvirvlet om dets graa Hoved. Og Kanaans
Land, Jordans Flod og det døde Hav, Morja Bjerg og
% --- p. 21 ---
\setcounter{footnote}{0}%
\opage{21}Zions Klippe, Libanons Høie, Esdrelons Dale, ja Roserne
paa Sarons Enge --- Alt er her endnu som i Oldtids Dage.

Ja, det Naturlige er her endnu, men hvor er det
Overnaturlige? En Torden kan rulle fra Sinais Top og
Lynene hvisle som Slanger; men hvor er Basunernes
Lyd? Hvor er Aandens myndige Ord? Jehovas Røst
er forstummet. Stormene vaagne, Uveiret farer endnu
ad de gamle Veie, men hvor er Cherubernes Vogn, og
hvor er Han, der fordum „agede paa Stormens Vinger?“
Hvor er den Almægtiges Arm? De forfærdelige Gjerninger,
de frelsende Gjerninger, de store Tegn, naar
han fra Skyen udrakte sin Haand? hvor er det altsammen,
og hvad er det blevet til? --- Hvad det er
blevet til, det skal jeg sige Dig, svarer den moderne
Kritik, hvis Du endelig vil vide det. Det er blevet til,
hvad det var og endnu er og altid vil være, religiøs
Phantasie, Frembringelser af en Aand, hvis inderste
Væsen er Blændværk.

Dersom Kritiken nu har Ret, saa maae vi jo kunne
prøve det. Er det et Blændværk, da maa det som
ethvert andet Blændværk forsvinde, naar vi i Dagens
Lys kalde Forstanden til Hjælp og anlægge Virkelighedens
Maalestok. Men paa hvem skal Maalestokken
anlægges? Ikke paa Israels Folk i Almindelighed; thi
dette Folk var med al sin Stivsindethed dog vægelsindet;
det skiftede Guder, det skiftede Tro, det bevægedes hid
og did som Sivet og Røret. Vi maae anlægge Maalestokken
paa de fremragende Karakterer, Folkets Førere,
Aandens Udkaarede; vi maae tage deres Liv og Gjerninger
skarpt i Øiesyn, for, om muligt, at opdage
Blændværket.

% --- p. 22 ---
\setcounter{footnote}{0}%
\opage{22}Hvis en Mand i vore Dage traadte frem for det
menige Folk og talede saaledes: Jeg veed, hvem der
har frastjaalet Eder Friheden og lagt Trældommens
Aag paa Eders Skuldre. Men frygter ikke, fatter Mod!
stoler paa Aandens Almagt og stoler paa mig; jeg skal
bryde det tunge Aag, jeg skal trods Hindringer og
Farer føre Eder sikkert til Maalet. Jeg behøver ikke
verdslige Midler; jeg bruger ikke kjødelige Vaaben; thi
den Villie, som virker i mig, er ikke min Villie, men
hans, som bærer Himmel og Jord\dots{}, hvis en Mand,
som sagt, optraadte saaledes, hvad vilde vi da dømme
om ham? Han er hildet i et beklageligt Blændværk.
Muligt, at han selv troer paa sin Sendelse og mener
det ærligt. Men han er en Sværmer, en Vanvittig,
mildest talt en Phantast. Lad ham forsøge sig og
prøve sine Kræfter; vi spaae ham en sørgelig Udgang;
vor Sandsynlighedsregning og vore statistiske Tabeller
bevise, at han vil lide Skibbrud med sit Forsøg og
strande paa Virkelighedens Klipper.

Men er det dog ikke besynderligt; i hiin gamle
Vidunderverden ere de store Førere netop traadte frem
paa en saadan Maade og have talt til Folket i et saadant
Sprog. Hvad skulle vi nu dømme? Vare hine
Mænd da Sværmere og Phantaster? Saa meget er vist,
at de ikke lede Skibbrud paa det Sandsynlige, men
førte deres Værk til Ende. Man dømme om det Overnaturlige,
som man vil, saa er her dog et Uforklarligt.
Hvis Troen paa det Overnaturlige kan bevise sin Sandhed
derved, at den giver Kraft til at løfte det Naturlige
og gjøre det Ophøiede virkeligt, saa indsees ikke, hvorledes
det Overnaturlige kan være det Usande, thi Troen
paa det Usande, en usand Tro, kan umulig formaae at
% --- p. 23 ---
\setcounter{footnote}{0}%
\opage{23}forvandle et Blændværk til Sandhed. Hvis Troen paa
det Overnaturlige er og forbliver en Indbildning, der i
det Høieste kun har psychologisk Sandhed, saa maa
der i hiin Vidunderverden have været en Tid, da den
psychologiske Sandhed var stærkere end den virkelige
Sandhed; men saa er det umuligt at forstaae, hvi den
stærkere Sandhed ikke, i det mindste for hiin Tids Vedkommende,
tillige skal gjælde for den høiere Sandhed.

Det er Karakterer af stor psychologisk Sandhed,
ophøiede Skikkelser, vi have for os. Lad os bede dem
træde frem, at vi ret kunne betragte dem; vi ville nævne
tre: \emph{Moses}, \emph{Samuel}, \emph{Elias}.

Da Moses er en verdenshistorisk Karakteer, det
jødiske Folks Lovgiver og Grundlægger af dets Statsforfatning,
er det naturligt, at man har sammenlignet
ham med andre af Oldtidens store Lovgivere, med Lykurg
i Sparta, Solon i Athen o. A. Men ved slige Sammenligninger
overseer man sædvanlig den karakteristiske
Forskjel og gjør Moses til Viismand i samme Forstand
som Solon. I sin Afhandling: „Die Sendung Moses“
har Schiller behandlet Sagen fra et rationalistisk Standpunkt.


Indviet i de ægyptiske Præsters Viisdom, landflygtig
fra Pharaos Hof, har Moses i Midians Ørk omsider
udgrundet, hvorledes han bedst skulde udfrie sit
Folk, indprente det Patriarchernes Tro paa den eneste
sande Gud og sikkre dets Fremtid som Guds Folk ved
en hensigtssvarende Lovgivning. Uden at have modtaget
nogen overnaturlig Aabenbaring, har han dog for
at skaffe sig Tiltro hos det uvidende Folk, maattet
nedlade sig til et fromt Bedrag, idet han har beraabt
sig paa en mirakuløs Kaldelse, har foregivet at staae
% --- p. 24 ---
\setcounter{footnote}{0}%
\opage{24}i et eget personligt Samkvem med Jehova selv, og er
saaledes traadt talende og handlende frem i Jehovas
Navn.

Det Uhistoriske i en saadan Opfattelse kan ikke
godtgjøres ved ligefrem Henviisning til det gamle Testamentes
Vidnesbyrd. Disse Vidnesbyrd kunne jo blot
bevise, enten at det fromme Bedrag er blevet saa
mesterlig gjennemført, at selv de Ypperste i Folket
have antaget de foregivne Mirakler, de udvortes saa vel
som Aandsmiraklet, for fuld historisk Sandhed, eller at
Præsterne, snildt anende Bedraget, have hjulpet til med
at inspirere Folkephantasien og i den derefter, maaskee
langt senere, optegnede Historie, ligeledes ved et fromt
Bedrag, givet de svævende Forestillinger om det Skete
en fast dogmatisk-mirakuløs Form og derved sikkret
den Fremstillingsmaade, der er typisk for Theokratiet.
Her er Noget, der aabent og uforbeholdent maa indrømmes
Kritiken. Det kan ikke godtgjøres, at Mosebøgerne
ere skrevne af Moses selv eller af et Øienvidne.
Det kan ikke skjules, at de Mellemled af Betingelser,
som skulde gjøre os en Forstaaelse mulig, ere paa mange
og høist vigtige Punkter udeladte. Det maa indrømmes,
at den Mirakeltroende ved at tage Alt efter Bogstaven,
hvad Indsigt i Sagens indre historiske Sammenhæng
angaaer, ikke kommer et Haarsbred videre end den
Ikke-Troende, som, med Forkastelse af det Overnaturlige,
holder sig til Sagn og Myther. Begge gjøre sig,
om end paa modsat Maade, skyldige i samme Feil: at
paastaae en Viden om det, som ikke kan vides. Vor
eneste Ledestjerne, det Eneste, vi have at holde os til,
er den \emph{psychologiske Sandhed}.

% --- p. 25 ---
\setcounter{footnote}{0}%
\opage{25}At den psychologiske Sandhed ikke kan bestaae
med nogetsomhelst vitterligt Bedrag, være sig fromt
eller ufromt, er dog vel uimodsigeligt. Men skulde den
ikke kunne bestaae med et ubevidst Selvbedrag? Selv
om vi antage dette, havde vi dog vundet Noget, som
er vist, den Erkjendelse nemlig, at Moses, hvad Aabenbaringsbevidsthed
angaaer, staaer paa samme, væsentlig
paa selvsamme Trin som Israels Folk med alle dets
Præster og Propheter. Det psychologiske Spørgsmaal
bliver altsaa dette: af hvilken Livskilde have Aandens
Mænd i det Gamle Testamente selv troet at øse det
Mod, den Viisdom og Kraft, som lod dem overvinde
fast overmenneskelige Hindringer og satte dem i Stand
til at udføre deres Mission?

Det er et eget Syn at see Moses, en vaabenløs
Mand, ja en Rømningsmand, staaende over for Majestæten
Pharao, bydende ham i Jehovas Navn og med
Jehovas Myndighed: Lad mit Folk fare! Pharao føler
sig som Magthaver, og en Magthaver slipper ikke saa
let sit Bytte. Hvad stoler han da paa, den vaabenløse
Mand? I Ørknens Eensomhed ved Bjerget Horeb, den
Gang han grundede paa Mulighederne, var han tvivlraadig
og bange; han følte sin Afmagt, og Sukket lød:
Herre, send en Anden. Nu da han midt i Begivenhedernes
Strøm seer Virkeligheden, som den er, seer
den mægtige Pharao selv i Øinene, nu er han slet ikke
bange. Han føler ikke sin Afmagt, tvertimod, der aander
en voxende Selvfølelse, en stedse stigende Følelse af
Magt i de gjentagne Ord: Lad mit Folk fare! Det er
en stor Forandring, der er foregaaet med Moses. Skulde
denne Forandring være skeet ved et Selvbedrag, da
maa der være Selvbedrag af høist forskjellig Orden;
% --- p. 26 ---
\setcounter{footnote}{0}%
\opage{26}thi Kjendsgjerningerne vise, at ogsaa Pharao befinder
sig i et Selvbedrag.

Mod Indbildninger og Blændværk har den virkelige
Verden stærke Midler; det gjælder en Styrkeprøve
mellem Moses og Pharao. Gik det an at stille Moses
i Række med Genierne, saa kunde man i de Sammentræf
af Omstændigheder, hvorunder de naturlige Landeplager
komme ham til Hjælp, og den Snildhed, hvormed
han forstaaer at benytte og tyde dem til sin Fordeel,
maaskee finde en Bekræftelse paa det bekjendte Ord:
Geniet har Lykken med sig. Men at opfatte Moses som
Geni er at forfeile den psychologiske Sandhed. Geniets
Bevidsthed er mythisk i sin Grund; Geniet i høiere
Potens er en Heros, avlet af en guddommelig Fader
med en menneskelig Moder, en Halvgud, i hvem Villien
er Naturkraft, og Naturkraften Villie: i Geniet har Illusionen
Kjød og Blod. Moses er, væsentlig forstaaet,
intet Geni; det mythisk Forblommede er hans Væsen
saa fjernt som muligt. De sandselige Tilsyneladelsers
Æsthetik er ham ganske ubekjendt. Intet af det Dæmringstrylleri,
som laaner Naturen Sjæl og gjør Klippeblokke
til Aabenbarelser af mægtige Aander, er at spore
hos ham. For hans Bevidsthed er Jorden Jord, Skyerne
Skyer og Havet Hav. Han er i sin Betragtning af
Virkelighedens Phænomener saa ædruelig som en Naturforsker
i det nittende Aarhundrede, i sin Bedømmelse
af de borgerlige Samfund saa praktisk forstandig, at
han vist gjerne havde brugt de statistiske Tabeller,
hvis han blot havde havt dem.

Moses, siger man, er et praktisk Geni, et Villiesgeni,
hans Vidunderkraft er hans Villie. Men dersom
den psychologiske Sandhed skal søges i dette Punkt,
% --- p. 27 ---
\setcounter{footnote}{0}%
\opage{27}hvor komme vi saa hen? Hvis det er muligt, at Den,
der virkelig maa føle, at han har al Styrke og Hjælp
at søge i sin egen Villie, kan hilde sig i den Indbildning,
at det er en af ham uafhængig Villie, en Villie
uden for ham, hvorfra Hjælpen kommer, da maae vi
tilstaae, at et saadant Selvbedrag nærmer sig stærkt
til det Mirakuløse. Hvorledes det er psychologisk forklarligt,
at et praktisk Geni med en prosaisk forstandig
Naturbetragtning kan --- naar det da ikke skeer ved
et vitterligt Bedrag, for at blænde Folket --- føle sig
inderlig forvisset om, at det ikke er Omstændighederne
eller i Tilfældets Spil Naturkræfterne, men Jehova selv,
som i Farens Stund kommer sin lydige Tjener til Hjælp,
formaae vi ikke at indsee.

Med Fjenderne bag efter sig og Havets Vande
foran sig er Moses og hans Folk i Fare. Hvad gjør
nu „det praktiske Geni“? Han rører Vandene med sin
Stav. „Han udrakte sin Haand over Havet“, siger den
gamle Skrift, „og Jehova lod Havet bortfare med en
stærk Østenvind hele Natten, og han gjorde Havet til
det Tørre, og Vandene skiltes ad“. Tage vi Jehova
bort, og beholde Østenvinden, saa er Tilfældet selv
mirakuløst; lade vi det praktiske Geni digte sin Villie
om til en almægtig Jehova, saa er Selvbedraget mirakuløst.
Kun i den Forudsætning, at Moses i sit Indre
oprindelig har følt sin Villie bevæget og styrket af
Jehovas Villie og desaarsag i det stærke Veir, som
skiller Vandene ad, gjenkjender den samme almægtige
Villie, er der psychologisk Sandhed. Denne psychologiske
Sandhed er Eet med Aabenbaringens Under.

At forklare Aandsmiraklet, det psychologiske Under,
er umuligt. Vor Tidsalders religiøse Psychologie giver
% --- p. 28 ---
\setcounter{footnote}{0}%
\opage{28}kun blege, farveløse Efterbilleder af det oprindelige
Gudsforhold. Sandheden af de indre Oplevelser er en
Troens Hemmelighed, som Enhver maa beholde for sig.
Den Troeskritik, der skal kunne anvendes paa Historien
om Moses, maa støtte sig paa følgende to Grundsætninger:
1) alle Beretninger om udvortes Mirakler maae
hente deres Lys fra Aandsmiraklet, 2) Aandsmiraklets
psychologiske Sandhed maa erkjendes ved en Bedømmelse
af Moses' Karakteer og det Værk, som var
hans Mission.

Nutidens Oplyste ere virkelig i Forlegenhed med
det Historiske i de mange Mirakler. Der er dem, der
ræsonnere saaledes: Alle Mirakler ere lige mirakuløse;
man maa desaarsag enten antage dem alle eller forkaste
dem alle. Antage dem kunne vi ikke, uden at
begaae et aandeligt Selvmord; vi maae altsaa forkaste
dem, d. v. s. anvise dem Plads i det store Musæum for
Oldtidens Sagn og Myther. Denne Synsmaade er kun
conseqvent, naar man med det Samme forkaster Aandsmiraklet
og gjør Moses' Tro til en Indbildning. Men
at holde paa Aandsmiraklet og stille Moses paa Aabenbaringens
Trin mellem Troens Forbilleder, skjøndt man
lader hans andre Mirakler staae hen i al Tvetydighed,
gaaer paa ingen Maade an. Er den Forestilling, Moses
selv har om et Mirakel, tvetydig, saa er hans Forestilling
om Jehova jo ogsaa tvetydig, og hans Tro tvetydig.
Moses er da en Mand af en tvetydig Karakteer.
Staaer det derimod rigtigt til med Moses' egen Tro; er
der ingen Tvetydighed, men fuld Klarhed i hans Indre,
naar han taler om sin Sendelse og om alt det, som
Jehova selv aabenbarer ham, altsaa ingen Tvetydighed
i hans Forestilling om Aandsmiraklet, da kan der heller
% --- p. 29 ---
\setcounter{footnote}{0}%
\opage{29}ikke være nogen Tvetydighed i hans Forestilling om
Begivenhedsmiraklerne i det Udvortes. Er Moses berettiget
til at optræde som Jehovas Udsendte, da er
han ogsaa berettiget til at betragte enhver overordentlig
Begivenhed, der tjener hans Formaal, som en Haandsrækning
af Jehova, et Tegn fra den Almægtige.

Her er slet ikke Spørgsmaal, hvorvidt Moses har
været i Stand til at drage nogen objectiv Grændse mellem
det, Jehova virker igjennem naturlige Aarsager, og det,
der skeer paa overnaturlig Maade. At Vandene i det
røde Hav adskiltes efter Jehovas Villie, derom er Moses
overbeviist, men det physisk-kritisk-metaphysiske Forhold
mellem Østenvinden og den guddommelige Villie
overlader han den moderne Bevidsthed at opdage. --- Saa
kan da, indvender man, Forskjellen mellem Mirakler
og naturlige Tildragelser fra dette Synspunkt baade
sættes vilkaarlig og udslettes vilkaarlig, den religiøse
Phantasie har altsaa lige saa frit et Spillerum i Aabenbaringslæren
som i Mythologien. Hvor er nu Troeskritiken,
og hvad faaer man ud af Aabenbaringsmiraklet?
Ikke at tale om, hvor naturligt det falder
Oldtidens Folk at indføre Guderne talende og handlende,
saa er det jo vitterligt, at Sværmere og Fanatikere have
foregivet indre Aabenbaringer, uden at man deraf kan
slutte at de vare Bedragere. Overnaturlige Syner kunne
jo, hvad Physiologien viser, forklares som Hallucinationer,
naturlige Udbrud af et religiøst Delirium.

Vi gjøre opmærksom paa, at det ikke er en Phantasie-
men en Villiesfordobling, en indre Modsætning af guddommelig
og menneskelig Villie, om hvis personlige
Karakteer Undersøgelsen dreier sig. Paa Frugterne
skal Træet kjendes, paa Handlingerne skal Villien kjendes;
% --- p. 30 ---
\setcounter{footnote}{0}%
\opage{30}paa den klare Opfattelse af Livets Modsætninger skal
Forstanden kjendes. Vi spørge: er der i Karakteren
Moses et eneste Træk, som tyder hen paa mythisk Forblanding
af Guddommeligt og Menneskeligt? Forraader
hans Inspiration Noget, som kunde vække Mistanke om
ekstatisk Forvirring? Bære Aandslovene fra Sinai maaskee
Spor af en syg Hjerne, et betænkeligt Præg af
Hallucinationer?

Kan Physiologien gjøre Rede for, hvorledes en
Mand, hvis inderste Væsen er Indbildning og Vildfarelse,
formaaer at udføre et Kæmpeværk som det: at frelse
et Folk fra Trældom? Har man nogen Psychologie, som
forklarer, hvorledes det gaaer til, at et Menneske først
omsætter sin egen Villie i Form af en almægtig Gudsvillie,
og derpaa, uden at ane det uhyre Selvbedrag,
bøier sig under denne Gudsvillie som en lydig Tjener,
arbeider med et utaknemmeligt Folk et langt, langt Liv
igjennem, uden at svækkes, uden at sløves, har man
virkelig nogen Psychologie, som forklarer os dette?

„Og Moses steg op fra Moabs Sletter til Bjerget
Nebo, til Toppen af Pisga, som er lige overfor Jericho,
og Jehova lod ham see Landet Gilead indtil Dan; og
hele Naphtali og Landet Ephraim og Manasse og hele
Landet Juda indtil Havet bagved det; og Sydlandet og
Sletten, Dalen i Jericho, Palmestaden indtil Zoar. Og
Jehova sagde: dette er det Land, som jeg har tilsvoret
Abraham og Isak og Jakob,\dots{} men du skal ikke komme
derind\dots{} Og Moses var eet hundrede og tyve Aar gammel,
da han døde; ikke blev hans Øine svækkede og ikke veg
hans Kraft fra ham. Men aldrig opstod der en Prophet
ydermere i Israel som Moses, thi med ham havde Jehova
staaet Ansigt til Ansigt.“ (5te Mosebog K. 34.).

% --- p. 31 ---
\setcounter{footnote}{0}%
\opage{31}Den næste Prophet, vi nævnede, var Samuel, den
sidste blandt Israels Dommere. Hans Tid er en anden,
hans mere begrændsede Virkekreds en anden, men
Aanden er væsentlig den samme. Tidlig, allerede som
Barn, blev Samuel kaldet til Prophet. Der hviler over
hans Kaldelse et ubeskriveligt Skjær af idyllisk Fred
og dog aandig, spændende, anelsesfuld Forventning. Viet
af sin fromme Moder til Jehovas Tjeneste, sendes han
som lille Dreng til Silo, til Ypperstepræsten Eli. Eli
er svag og hans Sønner ryggesløse; det varsler ilde for
Elis Huus. Samuel sover nær Helligdommen, hvor
Jehovas Lampe brænder. Han sover trygt i den stille
Nat, men vaagner op, thi en Stemme kalder. Han
troer, at det er den Gamle og gaaer ind til ham og
spørger: Herre kaldte du paa mig? Da det gjentager
sig, forstaaer Eli, at det er Jehova, der kalder paa
Drengen. „Derfor sagde Eli til Samuel: Gaa hen og
læg dig! naar nu Jehova kalder paa dig, skal du sige:
„Tal, Jehova! thi din Tjener hører.“ Saa kommer Aabenbaringen.“
Jehova sagde til Samuel: Nu vil jeg udføre
Noget i Israel, der skal runge i begge Øren paa den,
som hører det\dots{} jeg vil dømme Elis Huus for bestandig
paa Grund af den Misgjerning, at han vidste, hvorledes
hans Sønner bragte Forbandelse over sig, og dog revsede
han dem ikke.“ Hvad skal nu Samuel gjøre? Et Barn
og dog en Prophet, er han bange for at tale til Eli om
Synet. Dog Eli kaldte paa Samuel og sagde: „Hvad
var det, han talte til dig om? Nu maa du ikke dølge
det for mig. Gud straffe dig i sin Vrede, hvis Du
skjuler Noget for mig af alt, hvad han talede til dig.
Da sagde Samuel ham det altsammen og dulgte ikke
Noget for ham.“ --- Forunderlig er den tilsyneladende
% --- p. 32 ---
\setcounter{footnote}{0}%
\opage{32}Sindsro, hvormed den gamle Ypperstepræst hører sin
Dom. Han angrer ikke, thi hans Sjæl er mat; han
handler ikke, thi hans Villie er lam, han beder ikke,
thi hans Haab er udslukt. Med en næsten apathisk
Resignation svarer han: „Han er Jehova og kan gjøre,
hvad der tykkes ham godt.“

Og Varslerne gik i Opfyldelse. Samuel voxer op
til Mand og æres af Folket; han holder som Dommer
Jehovas Lov i Hævd og faaer i sin Alderdom det særegne
Kald at vælge og salve de to første Konger i
Israel.

I Samuels Historie er det Overnaturlige saa eiendommelig
indsprængt i det Naturlige, at Forstaaelsen
synes at ligge nær og dog pludselig undviger, naar vi
mene at have den i Hænde. Det er det Psychologiske,
der tiltrækker og frastøder. Samuel befinder sig i et
kritisk Vendepunkt, thi det mosaiske Theokratie er i
Færd med at skifte Form. Dommernes Tid var en
Urolighedernes og Lovløshedens Tid, „hver Mand gjorde,
hvad han fandt for godt“. Intet Under at Folket fandt
Tilstanden utaalelig og begjærede en Konge. Skjøndt
det var et Forlangende, der ingenlunde stred imod
Loven, er det dog forstaaeligt, at det maatte vække
Bekymring hos Samuel. Ikke at tale om, at Dommerværdigheden
derved gik tabt for hans Familie, saa var
her et religiøst Moment af særlig Vægt og Betydning.
Var Folkets Fordring udsprunget af Tro paa Jehova,
den evige Konge; mon ikke snarere af verdslig Attraa
og selvkloge Hensyn? Hvor fandtes den Mand, der
beklædt med jordisk Kongemagt, ubetinget vilde være
Jehovas Tjener? Det er forstaaeligt, at den gamle
Prophet, deels af hierarchisk Skinsyge, men især af Iver
% --- p. 33 ---
\setcounter{footnote}{0}%
\opage{33}for Aandens Ret imod Overgreb fra Kongens Side, kan
have fundet det nødvendigt at holde skarpt Øie med
Saul og ved given Anledning sætte ham paa Prøve.
Philisterne rykkede i Marken med en frygtelig Hær.
„Saul var endnu i Gilgal, og hele Folket, som var hos
ham, blev forfærdet. Han tøvede i syv Dage, indtil den
Tid, Samuel havde sagt, men Samuel kom ikke til
Gilgal, og Folket adspredtes\dots{} Da sagde Saul: Bringer
Brænd- og Takofferne hid til mig, og han forrettede
Brændofferet. Som han lige var færdig med Brændofferet,
kom Samuel, og Saul gik ham imøde for at
hilse ham. Samuel sagde, hvad har du gjort? Og Saul
sagde: Da jeg saa, at Folket adspredte sig fra mig,
og du ikke kom efter det fastsatte Antal Dage, og
Philisterne samledes i Michmas, saa tænkte jeg: Nu
komme Philisterne ned til Gilgal over mig, uden at
jeg har forsonet Jehovas Ansigt; derfor fordristede jeg
mig til at forrette Brændofferet. Og Samuel sagde til
Saul: Du har handlet daarligt,\dots{} nu skal dit Kongedømme
ikke bestaae. Jehova vil lede om en Mand efter
sit Hjerte og beskikke ham til Fyrste over sit Folk,
fordi du ikke har overholdt, hvad Jehova bød dig.“
Det er endelig af almeenpsychologiske Grunde forstaaeligt,
at Saul, som er forkastet af Propheten, og som
meer og meer faaer Mistanke om, at det er David, der
er udseet til at være hans Efterfølger, bliver bitter og
mørk, og seer i sine nærmeste Omgivelser Forræderi
og Rænker.

Ingen skal sige, at Aanden i denne Historie er
gaaet Naturen for nær. Saul og David ere ikke skyggeagtige
Væsener, theokratiske Abstractioner; det er
Individualiteter i Kjød og Blod med store Kræfter og
% --- p. 34 ---
\setcounter{footnote}{0}%
\opage{34}store Lidenskaber; hvilke Modsætninger af Lys og Skygge!
Sauls Dæmon og Davids Harpespil, de vilde Udbrud af
Sauls Had og Davids vidunderlige Lykke: Spydet udslynges,
men det rammer ham ikke. Kampe og Eventyr,
Venskab, Fjendskab, Elskov og Poesi, alt forener sig
om at gjøre denne første Strid om Thronen i Israel
saa tiltrækkende, som muligt. I Gibea under en Tamarisk
sidder Saul paa en Høi med sit Spyd; mørkt er
hans Aasyn, men hans Aand er endnu mørkere. Hans
Tjenere, der staae omkring ham, maae høre af hans
Mund, at de ere Forrædere tilhobe. Hvilken Smerte,
at Jonathan selv er i Forbund med Isai Søn! Achimelek
i Nobe kaldes til Regnskab, og Edomiten Doegs Sværd
nedhugger Jehovas Præster. Men Saul er viet til Undergang;
hans Kraft er brudt. Tilsidst er han kun en
Sjæleruin, en udbrændt Vulkan, en fortvivlet Mand, der
i Nattens Skjul søger Raad hos den Døde. Han tyer
til Kvinden i Endor. Hun spænder hans Phantasie og
maner Skygger frem af Graven. Den Døde taler: „Hvorfor
vil du spørge mig? Jehova er vegen fra dig og bleven
din Fjende.“ --- Prophetens Stemme! Saul er dømt og
synker i Afmagt for Samuels Skygge.

Men just som Nutidsbevidstheden ad denne Vei er
i Begreb med at tilegne sig Historiens Indhold, viser
sig pludselig det Frastødende: Jehovavillien med den
theokratiske Mechanisme afbryder Forstaaelsen. „Hvilken
Vilkaarlighed, at gjøre det til en Forbrydelse af Saul,
at han, i Farens yderste Øieblik, som Konge vil sørge
for sit Lands Forsvar, istedenfor at spilde Tiden med
at vente paa en Mand, der nøler og ikke holder sit
Ord“. Dog fra Aabenbaringens d. e. fra Aandsmiraklets
Synspunkt har Samuel Ret. Det, der fordres af Saul,
% --- p. 35 ---
\setcounter{footnote}{0}%
\opage{35}er ubetinget Lydighed mod Jehovas Lov, er med andre
Ord: Aandsvillie, ikke selvraadig Driftvillie. Ved at
mistroe Propheten har Saul selv undergravet den Grund,
hvorpaa han skulde staae, thi kun i Kraft af Prophetens
Valg er han Konge. Vi skulle tilstaae, at det mirakuløse
Gudsforhold, hvorpaa den profetiske Myndighed hviler,
ikke kan blive gjennemsigtigt for os; men vi skulle
tillige erkjende, at dette Gudsforhold, hverken er tilspeculeret
eller tilræsonneret, men oprindelig sat af
Aanden som Aand; vi skulle erkjende, at Israels Folk
har havt sit aandelige Liv i dette Forhold med lige
saa fuldgyldig Sandhed, som noget Nutidsfolk kan have
sit verdslige Liv i en Kulturidee; vi skulle erkjende,
at det, vi kalde det Haarde, det Vilkaarlige, det
Mechaniske, netop er det utvetydige Udtryk for den
Adskillelse af Gudsvillie og Menneskevillie, der betinger
en virkelig Aabenbarelse af Aandsvillien.

Vi gaae nu længere frem, for at møde Thisbiten
Elias. Ogsaa om denne Prophet gjælder det, at endskjøndt
Tidsvilkaarene ere forandrede, saa er dog Aanden
den samme. Efter at have løsrevet sig fra Juda Rige
og faaet Kalvedyrkelsen i Dan og Bethel indført, er
Israel i fuld Gang med at forandre Kalvedyrkelsen
til den groveste Afgudsdyrkelse. Det kommer da til
en hæftig Kamp imellem Jehova og Afguderne. Omris
Søn Achab, der har taget Jesabel, en Datter af Kong
Ethbaal i Sidon, til Ægte og bygget Baal et Bedehuus
i Samaria, er en ivrig Afgudsdyrker. Der tales om fire
og halvtredssindstyve Baalspropheter og om fire hundrede
Astartepropheter, som spiste ved Jesabels Bord.
Mod denne Folket ødelæggende Vederstyggelighed optræder
Propheten Elias. Hans Villie er en hellig Liden%
% --- p. 36 ---
\setcounter{footnote}{0}%
\opage{36}skab, Aanden i ham er Ild og Flamme; i sin Nidkjærhed
vover han Alt, han sætter Himmel og Jord i Bevægelse.
Han taler til Achab: „Saa sandt Jehova, Israels Gud
lever, hvem jeg tjener, saa skal der i disse Aar hverken
falde Dug eller Regn uden paa mit Ord.“ Saa flygter
han, ikke af Feighed, men paa Jehovas Bud; han drikker
af Bækken Krith, og Ravnene bringe ham Føde. Da
Bækken er udtørret, gaaer han til Enken i Sarepta;
Melet i Krukken slipper ikke op, og Olien i Kruset
mangler ikke: han opvækker hendes Søn fra de Døde.
Da Landeplagen havde varet paa tredie Aar, og Hungersnøden
var steget høit i Samaria, fik Elias Befaling at
vise sig for Achab. Achab møder ham med disse Ord:
„Er det dig, som forvirrer Israel?“ Propheten svarer:
„det er ikke mig, som forvirrer Israel, men dig er det
og din Faders Huus, fordi I svigte Jehovas Bud og du
følger Baalerne.“ Det kommer nu til en afgjørende
Prøve. Hele Israel med Baals og Astartes Propheter
forsamles, efter Achabs Bud, om Jehovas Prophet paa
Bjerget Karmel. „Da traadte Elias hen for hele Folket
og sagde: Hvor længe ville I halte til begge Sider?
Er Jehova Gud, saa følger ham; er Baal det, saa følger
ham! Men Folket svarede ham ikke et Ord.“ Først,
da Offerprøven var endt, og det havde viist sig, at den
døve Afgud ikke gav en Lyd, hvor meget hans Præster
end raabte og skare sig til Blods med Knive, men at
Jehova gav Elias Svar ved den Ild, som faldt ned og
fortærede hans Brændoffer: da kastede Folket sig paa
sit Ansigt og raabte: „Jehova er Gud, Jehova er Gud!“
Baalspropheterne dræbes, og Landeplagen hører op.
Men Jesabel raser. Elias maa flygte til Ørknen; mismodig
og træt, ønsker han sig Døden; „Jehova, tag nu
% --- p. 37 ---
\setcounter{footnote}{0}%
\opage{37}mit Liv, thi jeg er ikke mere værd end mine Fædre.“
Han styrkes og husvales og kommer til Gudsbjerget
Horeb; Nidkjærheden vil fortære ham, han klager og
beder: „Hærskarernes Gud! Israels Børn have svigtet
din Pagt, nedbrudt dine Altere og dræbt dine Propheter
med Sværd; jeg er alene tilbage, men ogsaa
mig tragte de efter at skille ved Livet.“ Men de svale
Vinde kjøle ham, et Pust af hellig Hengivelse bringer
Ro i hans Sind og dæmper hans Lidenskab. Jehova
selv vil komme. „En stærk og vældig Storm, som løsrev
Bjerge og knuste Klipper, gik foran, men Jehova
var ikke i Stormen. Efter Stormen kom der Jordskjælv,
men Jehova var ikke i Jordskjælvet. Efter
Jordskjælvet kom der Lynild, men Jehova var ikke i
Lynilden. Efter Lynilden kom der en stille, sagte Lyd;
da hyllede Propheten sit Ansigt i sin Kappe“, og fornam
i al Sagtmodighed Jehovas Nærværelse.

Kun endnu nogle Skridt, og han er ved Maalet.
Efter at have kaldet Elisa til Discipel, og forudsagt
Achab og Jesabel en blodig, forsmædelig Død, ender
han sit stormfulde Liv i et Uveir. „Han drog bort
fra Gilgal, og sagde til Elisa: Bliv kun her, thi Jehova
sender mig til Bethel; men Elisa svarede: Saa vist
som Jehova lever og din Sjæl, jeg forlader dig ikke.
Saa gik det i Fart fra Gilgal til Bethel, fra Bethel til
Jericho, fra Jericho til Jordan, hvor Elias svøbte sin
Kappe sammen og slog paa Vandet, saa det skiltes ad.
Idet de gik over, sagde Elias til Elisa: „Forlang Noget,
som jeg skal gjøre imod dig“, og Elisa bad: „Maatte
der kun være to Dele af din Aand over mig!“ Svaret
lød: „Du har ønsket dristig. Hvis du seer mig blive
% --- p. 38 ---
\setcounter{footnote}{0}%
\opage{38}borttagen fra dig, skal du faae det; hvis ikke, da faaer
du Intet.“

Hvor er det os fremmed og fjernt, hiint Israels
gamle Vidunderrige! Hvad nytter det vel at lade, som
om man bogstavelig troer Historien om Elias, og saa,
naar det kommer til en Forklaring, ikke at have Andet
at støtte sig til end nogle Elementer af almeenreligiøs
Psychologie med Talemaader og Forsikkringer og alskens
Sandsynligheder. Kun Aanden selv kan erkjende, hvad
Aanden vidunderlig frembringer. Havde Elisa blot seet
med sandselige Øjne, da havde han vel seet Lyn og
Uveirsskyer, men ikke Elias som han blev borttagen
fra ham; havde han kun seet det Virkelige i et naturligt
Lys, da havde han ikke seet Himmelfarten. Men
at han midt i det Sandselige virkelig saae det Oversandselige,
og inderlig i Jehovas Aand oplevede Miraklet,
høre vi af hans Udbrud: „Min Fader, min Fader, Israels
Vogn og hans Ryttere!“

Aabenbaringsvirkelighed er Aandsvirkelighed, Virkelighed
for Aanden.

% --- p. 39 ---
\setcounter{footnote}{0}%
\lecture{Tredie Forelæsning.}\opage{39}Forfatterne paa det gammeltestamentlige Trin staae
saa umiddelbart under Indtryk af det Overnaturlige, at
de i den efterdannende Virksomhed, hvorunder Beretningerne
nedskreves, typisk lade det Mirakuløse gribe
ind i bratte Overgange og overraskende Enkeltheder.
Begivenhederne ere ikke photographerede; de have
gjennemgaaet en Efterdannelse i religiøs Phantasie.
Man kan, hvad Principet angaaer, meget vel holde sig
overbeviist om Aandsmiraklets psychologiske Gyldighed
uden dog derfra at slutte til den historiske Overleverings
bogstavelige Ufeilbarhed. Fortolkerne paa Nutidens
Trin, kunne kun have et middelbart og afledet Indtryk
af det Overnaturlige; de ere henviste til en Opfattelse,
der stiller de sjælelige Naturbetingelser i Forgrunden;
de maae lade sig nøie med en Afglands af det
Mirakuløse, et Efterskjær, der sætter sit guddommelige
Præg paa de store Karakterers Tænken, Villen og
Handlen. I dyb Anerkjendelse af de prophetiske
Karakterers vidunderlige Styrke og sjælekraftige Sandhed
har Troespsychologien til alle Tider et religiøst Element,
hvori den kan mødes med Aabenbaringspsychologien.


% --- p. 40 ---
\setcounter{footnote}{0}%
\opage{40}Naar den prophetiske Aand særlig skal bestemmes
som en Spaadommens Aand, da vogte man sig vel for
at opfatte Spaadommen aandløst. Nogle Exempler.

\emph{Den blinde Achia} og \emph{Jeroboams Hustru}
(1 Kong. K. 14). Den samme Prophet, som havde forudsagt
Jeroboam, at han skulde blive Konge over de ti
Stammer, kommer i sin Alderdom til at forudsige ham,
at hans Undergang er nær, at hans Huus skal falde.
De nærmere Omstændigheder ere bekjendte: Jeroboams
Søn Abia blev syg; Kongen, som vidste med sig selv,
hvor dybt han havde krænket Propheten, lader sin
Hustru forklæde sig og gaae til Silo, for at høre en
Spaadom om Sønnens Liv. „Saaledes gjorde Jeroboams
Hustru, hun stod op og gik til Silo og kom ind i
Achias Huus. Denne kunde ikke see, thi hans Øines
Lys var udslukket af Alderdom.“ Holde vi nu det
Overnaturlige i Baggrundens Dæmring, da er det ikke
uforklarligt, at den blinde Seer kan have gjenkjendt
den forklædte Dronning, at han kan have forudsagt
Død og Fordærvelse over det hele afgudiske Huus og
at det Billede af Fald og Ødelæggelse, den gamle
Blinde saae i Aanden, kan i sine Enkeltheder have
været desto forfærdeligere, jo stærkere hans Inspiration
i det givne Øjeblik var. Dette er da ogsaa Fortællingens
almeenpsychologiske Grundkarakteer. Men Aabenbaringspsychologien
fordrer en slaaende, mirakuløs Indvirken
af det Overnaturlige, og Beretningen lyder: „Jehova
sagde til Achia: See, Jeroboams Hustru kommer ind
for at udfritte dig angaaende hendes Søn, thi han er
syg. Saa og saa skal du tale til hende, thi naar hun
kommer, vil hun lade, som om hun er fremmed. Saa
snart Achia hørte hendes Skridt nærme sig til Døren,
% --- p. 41 ---
\setcounter{footnote}{0}%
\opage{41}sagde han: Kom ind Jeroboams Hustru, hvorfor lader
du, som om du var fremmed; jeg er jo sendt til dig
med Sørgebudskab. Gaa hen og sig til Jeroboam: Saa
siger Jehova, Israels Gud: Fordi jeg ophøiede dig over
det Folk, hvortil du hørte, og satte dig til Fyrste over
mit Folk Israel, rev Kongeriget fra Davids Huus og
gav dig det; og fordi du ikke var som min Tjener
David, men handlede ugudeligere end alle....., see
derfor vil jeg bringe Ulykke over dit Huus..... Hvem
der døer hos Jeroboam, skulle Hundene i Byen æde,
hvem der døer paa Marken, skulle Himlens Fugle fortære;
det har Jehova sagt. Reis dig nu og gaa hjem,
naar din Fod træder paa Byens Gade, skal Drengen døe.“

Det vil indsees, at det hverken kan være den ordrette
Fastholden af alle Enkeltheder, ei heller den
naturalistiske Bortforklaring af det Overnaturlige, endnu
mindre den hinkende Mægling mellem Sandsynligt og
Usandsynligt, hvori en Fortolkning her kan søges.
Enhver Uhildet maa kunne see, at det netop er Straalebrydningen
af overnaturligt og naturligt Lys i Aabenbaringspsychologiens
Prisme, der giver Billedet Karakteer
og Farve.

\emph{Nebukadnezars Drøm} (Daniel K. 2). I Aabenbaringspsychologien
er Blikket umiddelbart rettet mod
det Mirakuløse, hvad enten det nærmest kommer indenfra
eller udenfra. Hvorpaa skal en sand Prophet vel
kjendes fra de falske? Paa den Kjendsgjerning, at den
Alvidende aabenbarer ham Ting, hvorom intet Menneske
af sig selv kan have nogen Viden. Det er dette Beviis,
der bliver ført for Nebukadnezar, Kongen af Babylon,
Jehovatemplets første Ødelægger. I sit andet Regjeringsaar
drømte Nebukadnezar en Drøm; „og hans Aand
% --- p. 42 ---
\setcounter{footnote}{0}%
\opage{42}gyste, og med Søvn over ham var det forbi. Og Kongen
befalede at kalde de Skriftlærde og Besværgerne og
Troldmændene og Chaldæerne, for at forkynde ham
hans Drøm; og de kom og stode for Kongens Ansigt.
Og Kongen sagde til dem: jeg har drømt en Drøm, og
min Aand gyste, saa jeg vil vide Drømmen. Og Chaldæerne
sagde til Kongen paa Aramaisk: Kongen leve
evindelig! siig dine Tjenere Drømmen, og vi ville angive
Udtydningen. Kongen svarede og sagde til Chaldæerne:
Befalingen er udgaaet fra mig; dersom I ikke
tilkjendegive mig Drømmen og dens Udtydning, skulle
I vorde huggede i Stykker og Eders Huse gjøres til en
Skarndynge.“ Nebukadnezar er paa een Gang en klog
Mand og en gal Mand, et praktisk Geni og en Phantast.
Chaldæerne gjøre Omsvøb; men Kongen staaer
fast paa, at saa sandt de kunne \emph{udtyde} ham Drømmen,
maae de ogsaa kunne \emph{sige} ham Drømmen. Kunne de
ikke det Sidste, saa er det Løgn og Bedrageri med det
Første. En ubøielig og sund Logik! Paa denne Logik
bleve de chaldæiske Viismænd henrettede. Daniel
kaldes; han stedes for Kongen og siger: „Hemmeligheden,
som Kongen begjærer at vide, kunne de Vise,
Besværgerne, de Skriftlærde, Stjernetyderne ikke kundgjøre
Kongen. Men der er en Gud i Himmelen, som
aabenbarer Hemmeligheder, og han har kundgjort
Nebukadnezar, hvad der skal skee i de sidste Dage.“ I
Hemmelighedernes Aabenbarelse synes Aandsmiraklet
med Eet at blive saa ubetinget, at ethvert Glimt af
naturlig Forstaaelse maa forsvinde; Nutidens religiøse
Psychologie er nær ved at fortvivle. Og dog er her
Antydninger af almeenpsychologisk Natur. Overfor den
vældige Erobrer, den af sin Storhed og Magt berusede
% --- p. 43 ---
\setcounter{footnote}{0}%
\opage{43}Nebukadnezer, staaer Daniel med en Sindsro, en Overskuelsens
Klarhed og Aandsoverlegenhed, som minder
om Moses, da han stod over for Pharao. --- Hvad
drømmer vel en Nebukadnezar om? Sin egen og sit
Riges Storhed! Hvad kan det vel være for Drømmebilleder,
som i Syner over hans Hoved omkring hans
Leie kunne bringe hans Aand til at gyse? Billeder af
Storheden og Storhedens Fald! „Du Konge, du saae,
og see! et Billede, samme Billede var stort og Glandsen
deraf overordentlig; det stod foran dig, og dets Udseende
var forfærdeligt.“ Strax ved denne Begyndelse
kan Propheten læse i Kongens Aasyn, at han har truffet
det Rette; han fortsætter: „Hovedet paa dette Billede
var af reent Guld, Brystet og Armene af Sølv, Underlivet
og Hofterne af Kobber; Benene derpaa vare af
Jern, Fødderne deels af Jern og deels af Leer.“ Det
psychologiske Eftertryk ligger paa, at Drømmens Indhold
ved Daniels Ord paa Ny bliver til for Nebukadnezars
indre Syn, og vistnok i klarere, bestemtere,
mere betydningsfulde Træk. „Du saae, indtil en Steen
blev hugget løs, og det ikke ved Hænder; og denne
slog Billedet ved dets Fødder, der vare af Jern og Leer,
og knuste dem. Derpaa knustes tilsammen Jernet,
Leret, Kobberet, Sølvet og Guldet, og de bleve som
Avner fra en Tærskeplads om Sommeren, Vinden borttog
dem, og deres Sted blev ikke fundet mere; men
Stenen, som sønderslog Billedet, blev til en stor Klippe
og opfyldte hele Jorden.“ Nebukadnezar lytter og
grunder; han drømmer sin Drøm om, idet han hører
Prophetens Udtydning; og Udtydningen er i Sandhed
storartet.

% --- p. 44 ---
\setcounter{footnote}{0}%
\opage{44}Opfattelsen kan nu svinge mellem to Poler, den
mirakuløse og den almeenpsychologiske, men stadig
inden for Aandens Kreds. Er der Sandhed i Drømmen
og dens Udtydning, saa er der ogsaa Virkelighed; men
det er ikke Politik-Virkelighed med statistiske Tabeller
til Ledetraad i Statshistorien, det er med al sin historiske
Vægtfylde Aandsvirkelighed, Virkelighed for
Aanden.

\emph{Messianske Spaadomme}. Saa længe Aandsmiraklet
er i fuld psychologisk Virksomhed, og Aabenbaringstroen
lever i sit oprindelige Gudsforhold, kunne
Spaadommene opfattes enkeltviis, ja sætningsviis, uden
at Opfattelsen derfor bliver aandløs: et Navn, et Billede,
et flertydigt Udtryk er ofte nok til at sætte en Strøm
af aandeligt Liv i Bevægelse. Men naar Indtrykket af
det Oprindelige er afmattet, naar det Nærværende er
blevet saa fattigt paa Aand, at det kun lever i skyggeagtige
Efterdannelser af det Forbigangne og maa trøste
sig med Haabet om et Tilkommende, da er det en betænkelig
Sag med sætningsvise Udtydninger af gamle
Spaadomme.

Propheterne have spaaet om en Messias, som
skulde komme. Kan man virkelig paavise, at disse
Spaadomme passe paa Jesus af Nazareth og paa ingen
Anden? De Kristne paastaae det, Jøderne benægte det,
begge Parter med lige Ret, saa længe man eensidig
holder sig til Ord og Sætninger. Det kan fra dette
Synspunkt ikke nægtes, at Rabbinernes Indvendinger
ofte have glimret ved skarp Kritik og fordunklet vore
kirkelige Dogmatikeres Forsvar. Vi vælge et eneste
Sted (Jes. 7, 14—16): Propheten taler til Judas Konge,
som, truet af overmægtige Fjender, i sit Mismod ikke
% --- p. 45 ---
\setcounter{footnote}{0}%
\opage{45}tør begjære et Frelsens Tegn fra Jehova: „Derfor skal
Jehova give Eder et Tegn: See, Jomfruen vorder frugtsommelig
og føder en Søn, og du skal kalde hans Navn
Immanuel; Mælk og Honning skal han spise, indtil han
lærer at forkaste det Onde og vælge det Gode. Thi
førend Drengen lærer at forkaste det Onde og vælge
det Gode, skal det Land vorde øde, for hvis tvende
Konger“ (Rezin, Konge i Syrien, og Pheka, Konge i Israel)
„du føler Gru.“ Dette skal nu, ifølge den orthodoxe
Dogmatik, absolut være et messiansk Skriftsted,
ja et Hovedsted. „Den Jomfrufødte, Immanuel,
Gudmennesket, et Frelsens Tegn! --- et Skriftsted, som
ovenikjøbet er citeret af Evangelisten Matthæus (K. 1,
v. 22—23): den, som ikke kan see, hvorledes Propheten
her har \emph{skuet} i \emph{Aanden} det store frelsende Under,
at Jesus i Tidens Fylde skulde fødes af Maria, den rene
Jomfru, maa ikke have theologiske Øine i Hovedet.“
Der behøves imidlertid hverken rabbinsk Lærdom eller
overtheologisk Skarpsindighed for at opdage, at den
orthodoxe Paastand hviler paa svage Grunde. Om en
mirakuløs Jomfrufødsel er her slet ikke Tale. Propheten
taler, ifølge Grundtextens Ord, ikke om den
rene Jomfru, men snarere om sin unge Hustru, hvis
Sønner (jvnfr. K. 8. v. 3) han giver symbolske Navne.
Den Frelse, hvorom han varsler, er saa langt fra at
henlægges til en fjern Fremtid, at det tvertimod ved
Ordene: „førend Drengen lærer at forkaste det Onde
og vælge det Gode“, udtrykkelig tilkjendegives, hvor
nær Frelsen er. Naar Landet skal være frelst, førend
Immanuel har lært at skjelne mellem Godt og Ondt,
hvorledes kan Immanuel saa være Frelseren? Naar
derimod Evangelisten, benyttende den alexandrinske
% --- p. 46 ---
\setcounter{footnote}{0}%
\opage{46}Oversættelse, seer tilbage paa det jesajanske Sted, betoner
Ordet „Jomfruen“, og saaledes finder Spaadommen
anden Gang opfyldt, ja i høiere Betydning fuldkommen
opfyldt, da maae vi i Aabenbaringens Navn give ham
Ret; men hans Ret er i Aanden.

Betragtet i Aandens Lys, er hele det gamle Testament
en Spaadom om det Nye, saa vist som den opvoxende
Yngling er en Spaadom om Manden, den sig
udfoldende Blomst en Spaadom om Frugten. Loven er
Evangeliets Grundlag, og Evangeliet er Lovens Fylde.
Men Overgangen skeer ikke rolig organisk som i Plantelivet;
Aabenbaringens høiere Trin forholder sig ikke
ligefrem til det lavere som Manddommens Daad til
Ungdommens Drømme.

Hvad der giver Moseloven Præg af guddommelig
Aabenbaring, nemlig den ubetingede Forudsætning, at
ethvert af dens Bud er umiddelbart dikteret af Jehova
selv, sætter Aandsvillien i Stand til at hævde sin Høihedsret
over Menneskelivets Naturside i alle borgerlig-sædelige
Forhold og grundlægge en Stat, der ikke tilsteder
det politiske Liv nogen særlig Selvstændighed
over for det religiøse. En Følge heraf var, at Propheterne,
skjøndt ikke beklædte med verdslig Magt,
kunde udøve en mægtig Indflydelse paa Styrelsens
Gang, og deres myndige Ord være afgjørende for Krig
og Forbund med nærmere og fjernere Folkeslag. Det
berettes (Jes. K. 7), at da Rygtet om Rezins og Phekas
truende Fremrykken udbredtes i Juda, „da bævede
Kongens Hjerte og hans Folks Hjerte som Træerne i
Skoven bæve for Vinden“. Paa Jehovas Bud gaaer
Jesaias Achas imøde og siger: „Tag dig iagt, hold dig
rolig, ei frygte du og ei vorde dit Hjerte forsagt\dots{}
% --- p. 47 ---
\setcounter{footnote}{0}%
\opage{47}for disse to rygende Brande, for Rezins og Syriens og
Remaljas Søns fnysende Vrede. Syrien har oplagt onde
Raad imod dig med Ephraim og Remaljas Søn, de sige:
Lader os drage op til Juda, trænge det, tage det til
os og sætte Tabeel til Konge over det. Men saa siger
Herren Jehova: Det skal ei stadfæstes, det skal ikke
skee.“ Achas og Folket see Faren der, hvor den ikke
er, og see den ikke der, hvor den er. De ville slutte
Forbund med Assyrien, maaskee endog søge Hjælp hos
Ægypten, og tænke ikke paa, at naar disse store Hedningefolk
faae Indpas i Landet, ville de træde Alt under
Fødder og ødelægge Theokratiet. Ægyptere og Assyrere
komme ikke som Hjælpere; de komme først paa Vredens
Dag som Jehovas Redskaber til at fuldbyrde en Straffedom.
„Paa den Dag skal Jehova fløite ad Fluen, som
er ved Ægyptens Floders Grændser, og ad Bien, som
er i Assurs Land; da skulle de komme og leire sig tilsammen
i de afsluttede Dale, i Klipperevner og paa
alle Tornehække og paa alle Enge \dots{}“ Paa hvert
Sted, hvor der var tusinde Viinstokke til tusinde Sølvsekler“,
der græsse nu Øxne og Faar og vilde Dyr.
„Med Pile og med Bue skal man gaae derhen, thi til
Torne og Tidsler vorder hele Landet.“

Menneskelig seet, er Propheternes Liv uophørlig i
Fare. Ordets Myndighed og Kongernes Magt, Evighedens
Ret og Øieblikkets Krav, Lovens Fordringer og
Folkets Lidenskaber: her er ulige Vaaben! Israels
Folk, som efter sin Grundbestemmelse er et Aandens
Folk, Jehovas Folk, er efter sin elementære Naturbestemmelse
tillige et sandseligt, selvisk og haardhjertet
Folk, et gudfrygtigt og dog ugudeligt Folk, trofast og
dog troløst, lydigt og dog trodsigt, modigt og dog
% --- p. 48 ---
\setcounter{footnote}{0}%
\opage{48}forsagt, opoffrende og dog egoistisk, et stridigt og gjenstridigt
Folk, i Strid med sin Gud og i Strid med
sig selv.

Det tydeligste Vidnesbyrd om Aandens Magt i den
prophetiske Inspiration er den Redelighed, den Villiesreenhed
og det urokkelige Mod, hvormed Jehovas
Sendebud udrette hans Ærinde. Propheterne ere folkelige
Mænd i Ordets sande Betydning; de leve med
Folket og for Folket, de trøste det i Prøvelsens Tid,
de styrke dets Haab, de vække dets Kraft og aabne
det lyse Udsigter; men altid paa een Betingelse, den
ufravigelige Betingelse, at Folket ærligt vil holde sig
til Lovens Ord og Aand. Aldrig høre vi dem sige:
dette Folk er et Guds Folk, thi dets Hjerte er som
Guds Hjerte; hvad et saadant Guds Folk med et Guds
Hjerte nu saadan har faaet Syn for, om det saa var
for lidt Kalvedyrkelse, maa dog være det Rette: Folket
kjender det, Folket forstaaer det. Nei, naar Folket
synder, og Aanden vredes, saa lyder Ordet strengt; vi
høre Jesajas (K. 1) raabe: „Vee et syndigt Folk! ---
En Oxe kjender sin Eier, et Asen sin Herres Krybbe;
men Israel kjender det ikke, mit Folk forstaaer det
ikke.“

Lovaabenbarelsens Særkjende maa netop søges
deri, at Aandsvillien staaer utilsløret, reen og klar over
for Driftvillien; det sandselige Selv skal bøie sig
ubetinget for Aandens Bud. Det var det Haarde ved
Israels Udvælgelse, at den guddommelige Aand kom
det naturlige Menneske saa umiddelbart nær. Mod
Vreden i Prophetens Ord blev Folkets Hjerte forhærdet,
thi Sandheden saarer; naar Aanden lyste som Himlens
Sol, blev Øjet blændet af de stærke Straaler. Sand%
% --- p. 49 ---
\setcounter{footnote}{0}%
\opage{49}heden og Forhærdelsen, Aabenbaringen og Forblindelsen
svare psychologisk til hinanden. Det forblindede Folk
slog sine Propheter ihjel; Forholdet til Aand gjorde
Israels Børn til Aandsforbrydere: hvilken forunderlig
Ening af Velsignelse og Forbandelse! Det var Frihedens
Vilkaar, og dog en Skjæbne saa haard, at Kjærligheden
selv maatte græde derover. „Jerusalem, Jerusalem, du
som ihjelslaaer Propheter og stener dem, som ere sendte
til dig, hvor ofte vilde jeg forsamlet dine Børn, ligerviis
som en Høne forsamler sine Kyllinger under sine
Vinger; men I vilde ikke.“

At Jehovas Lov maatte opfyldes, derom vare Propheter
og Præster ubetinget enige. Men Opfyldelsen
har to Sider, en udvortes og en indvortes. Der gives
en Efterlevelse af de enkelte Bud i det Udvortes, som
til en vis Grad er betinget af Lydighed, men Lydigheden
er kun et Skin, thi Sindelaget svarer ikke til
Gjerningen. Den Opfyldelse, som Aandsloven kræver,
er væsentlig indvortes, en Lydighed i Sind og Hjerte,
en Efterlevelse i Gjerning i Sandhed. Ved at kaste et
Blik paa Præsternes og Propheternes ulige Vilkaar i
det jødiske Theokratie vil det blive os forstaaeligt, at
hine nærmest holdt sig til det Udvortes, disse frem
for Alt til det Indvortes. Præsterne, den hellige Overleverings
Vogtere, det Bestaaendes Støtter og Bærere,
havde det Kald at varetage Offertjenesten i Templet,
oplæse Loven, sørge for Renselser, Sabbather og Høitider,
og nøie paasee, at de religiøse Ceremonier foregik
efter den foreskrevne Orden. Deres Virksomhed
var, hvad man kalder officiel, og saaledes ogsaa deres
Stilling. Levi Stamme med faste Stæder og faste Indtægter,
Tiender, foreskrevne Dragter o. s. v. \emph{afgav}
% --- p. 50 ---
\setcounter{footnote}{0}%
\opage{50}Theokratiets Embedsmænd. Propheterne vare ikke
Embedsmænd, havde hverken Tiender eller Ornater
eller faste Boliger; de levede frit som Fuglene under
Himlen, og toge til Takke med det Uvisse: vilde Menneskene
ikke sørge for dem, saa maatte Ravnene gjøre
det. --- Det var med høitidelige Ceremonier, at Moses
(3. Moseb. K. 8) indviede Aron. „Han gav ham Underkjortelen
paa, og han ombandt den med Bæltet og
iførte ham Overkjortelen; og han lagde Livkjortelen
derover, og han ombandt ham med Livkjortelens broderede
Bælte. Og han lagde Brystsmykket over ham,
og paa Brystsmykket satte han Urim og Thumim. Og
han satte en høi Hue paa hans Hoved og paa Huen
foran en Guldplade, den hellige Prydelse.“ Saa var
Ypperstepræsten Aron i fuldt Ornat. Paa lignende
Maade bleve Præsterne, Arons Sønner, begavede med
Underkjorteler og Bælter og høie Huer. Her var noget
Udvortes, Noget for den forsamlede Menighed at see
paa. Med Propheternes Kaldelse forholdt det sig
anderledes. „Elias traf Elisa, Saphats Søn, som pløiede
med tolv Spand Oxer for sig (han trak selv med det
tolvte); idet Elias gik forbi kastede han sin Kappe
over ham (1 Kong. K. 19); dermed var Elisa kaldet.
Aanden gjør ikke mange Ceremonier; den kommer, og
saa er det nok: Amos, Hyrden fra Thekoa, blev kaldet
fra Faarehjorden.

Det er forstaaeligt, at Propheterne gjorde en ganske
anden Fordring til det Indvortes end Præsterne. Fordringen,
Nødvendigheden af en indvortes Lovopfyldelse
har Jesajas (K. 1) udtalt med en Røst, der lyder
igjennem Tiderne, og med et Sandhedens Eftertryk,
der kan bringe Udvortesheden selv til at ryste: „Hører
% --- p. 51 ---
\setcounter{footnote}{0}%
\opage{51}Jehovas Ord, I Sodomas Fyrster; mærk Guds Ord, du
Gomorrhas Folk! Hvad skal jeg med eders mange
Offere, siger Jehova; jeg er mæt af Vædderes Brændoffer
og af Slagtekvægets Fedme; til Studes og Lams
og Bukkes Blod har jeg ikke Lyst. Naar I komme at
sees for mit Aasyn, hvo har forlangt dette af eders
Haand? I søndertræde mine Forgaarde\dots{} Jeg
skjøtter ei om Uret og festlig Dag. Eders Maaneder
og eders bestemte Høitider hader min Sjæl; de ere mig
en Byrde, jeg er kjed af at bære dem. Naar I udstrække
Eders Hænder, skjuler jeg mine Øine for Eder.
Om I end vil gjøre eders Bønner mange, jeg vil dog
ikke høre Eder: eders Hænder ere fulde af Blod. ---
Lader af at gjøre ilde, lærer at gjøre godt;\dots{} skaffer
den Faderløse Ret, fører Enkens Sag; kommer saa,
og lader os gaae i Rette med hinanden, siger Jehova,
og om saa eders Synder end ere røde som Purpur,
som Snee skulle de vorde hvide“ (jvfr. Jesajas Kap. 48;
Kap. 58).

Det gaaer i den aandelige Verden som i den
legemlige, at en Bevægelse, der skal overvinde Hindringer,
forudsætter et Tryk og en Drivkraft. De forfærdelige
Trusler i Forening med de store Forjættelser
afgave, hvad Lovopfyldelsen angaaer, Trykket og Drivkraften.
Havde Forjættelserne udelukkende været
knyttede til en udvortes Opfyldelse af de udvortes Bud
efter Bogstaven, da vilde Drivkraften snart været opbrugt,
thi Erfaringen lærte, at en retfærdig Gjengjældelse
ikke fandt Sted i det Udvortes; men det prophetiske
Ord holdt Tvivlen nede ved at sætte Skylden
i Sindelaget, pege paa de i Sandhed Retfærdige
(Ezechiel 14, 14) og henvise til Fremtiden. Havde
% --- p. 52 ---
\setcounter{footnote}{0}%
\opage{52}Forjættelserne udelukkende været knyttede til den indvortes
Opfyldelse, da vilde det faste mosaiske Grundlag
have løst sig op i en ubestemt Aandelighed, en
formløs Inderlighed, der stred mod Aabenbaringens Formaal: at opdrage Folket til virkelig Erkjendelse
af Jehovas Villie, til ubetinget Lydighed mod Aandens
Love.

For at forstaae de gamle Propheter, maae vi ret
mærke os dette Eiendommelige, at Lovens Opfyldelse
lige saa meget peger udefter som indefter, og at det
Samme gjælder om Trusler og Forjættelser, medens det
Udvortes dog er og forbliver saa ueensartet med det
Indvortes, at det Mirakuløse overalt skinner igjennem.
Den psychologiske Spænding, som herved fremkommer
og sætter Troesphantasien i Røre, kaster Forstaaelsens
Skriftlys paa de Brændinger og Bølgegange, hvorunder
de prophetiske Taler idelig bevæges mellem Kjendsgjerning
og Billede, Virkelighed og Symbolik. --- „Vee
dem, der staae aarle op om Morgenen, fordi de hige
efter Drik. De sidde langt ud paa Natten, og Vinen
hidser dem. Harper, Cithare, Pauker, Fløiter og Viin
er ved deres Gjæstebud, men Jehovas Gjerning ville de
ikke see. Vee dem, som sige om det Onde: det er
godt, og om det Gode, det er ondt. De gjøre den
Ugudelige til retfærdig for en Løns Skyld og fjerne
den Retfærdiges Ret fra ham. Derfor, som Ildens
Flamme fortærer Halm, og tørt Græs synker sammen
i Luen, saa skal deres Rod være som Muggenhed og
deres Blomster fare hen som Støvet; thi de foragtede
Jehova Zebaoths Lov, de lode haant om Israels Helliges
Ord. Derfor blusser Jehovas Vrede op mod hans Folk;
han udstrækker sin Haand over det, han slaaer det
% --- p. 53 ---
\setcounter{footnote}{0}%
\opage{53}saa Bjergene skjælve: Ligene ligge som Skarnet paa
Gaden.“ (Jes. K. 5.) Dette Vexelskin af Billede og
Kjendsgjerning, Symbol og Virkelighed er karakteristisk
for de messianske Spaadomme.

\emph{Messiasbilleder}. I Loven er Gudsforholdet
ubetinget sat som et Villiesforhold; af Lovopfyldelsens
strænge Fordringer udspringe de billedrige Skildringer
af den forventede Frelser. Først seer det ud, som
kunde Folket virkelig opfylde Loven, naar det kun alvorlig
vilde; men efterhaanden bliver det indlysende,
at den samme Gud, som har givet sit Folk Loven,
ogsaa ved sin Aand maa forhjælpe det til at faae Loven
opfyldt. Forholdet mellem Jehovas overnaturlige Villie
og den menneskelige Frihed bliver da nærmere bestemt
saaledes: at Loven overtrædes er Menneskenes egen
Skyld, og de guddommelige Straffe retfærdige; men at
Loven opfyldes, er betinget af Guds Naade, Guds
Aand, thi af sig selv har Mennesket ikke Kraft til at
gjøre Loven Fyldest. Den virkelige Tilstand er nu den,
at Folket paa Grund af Overtrædelserne, befinder sig i
Nød og Fare, trykket og truet af overmægtige Fjender.
Frelsens Vilkaar er, at Folket erkjender sin Skyld og
sin Afmagt, omvender sig til Jehova og stoler ubetinget
paa hans Forjættelser. Naar da Propheten skildrer
Synden og Straffen, holder han sig nær til det Virkelige,
om end Trækkene vise sig forstørrede i Aabenbaringens
Speil, men taler han om Naaden og Frelsen,
saa gaaer den i Haabet glødende Phantasie langt ud
over Virkelighedens naturlige Grændser; Lyset fra det
Overnaturlige bryder frem med farverige Straaler og
forherliger hele Naturen. „Det Folk, som vandrede i
Mørke, skuer et stort Lys, et Lys er opgaaet over dem,
% --- p. 54 ---
\setcounter{footnote}{0}%
\opage{54}som boe i Dødens Skyggers Land\dots{} Thi et Barn
er født os, en Søn er givet os. Fyrstendømmet skal
være paa hans Skuldre; hans Navn skal kaldes Under,
Raadgiver, stærke Gud, Evighedens Fader, Fredens
Fyrste. Paa Fyrstendømmets Storhed og paa Freden
er der ingen Ende“. --- Davids Throne og hans Kongerige
skal befæstes ved Ret og Retfærdighed (Jes. Kap. 9).
Spørge vi, hvem denne Fredsfyrste er, og hvor fra han
skal komme, saa svarer Propheten (Jes. Kap. 11):
„Der udgaaer en Kvist af Isai Stub, og et Skud af
dens Rødder skal blomstre frem. Paa ham hviler Jehovas
Aand, Viisdoms og Forstands Aand, Raads og
Krafts Aand, Kundskabs og Jehovas Frygts Aand\dots{}
Han slaaer Jorden med sin Munds Stav og med sine
Læbers Aande skal han dræbe de Ugudelige\dots{} Da
skal en Ulv dvæle ved et Lam, og et Pantherdyr ligge
ved et Kid; en Kalv og en Løveunge skulle være
sammen: et Barn kan føre dem.“

Det er det aandelige Israel, af hvis Rod den Vidunderlige
fremskyder. Grundtrækkene af hans Væsen,
hans Virken og timelige Vilkaar kan sees i det udvalgte
Folks Historie; thi det er dette Folk, og dette
alene, som har maattet taale og lide for Aandsvilliens
Seier, for Sandhedens og Retfærdighedens hellige Sag.
„Han opvoxede som en Kvist for Jehovas Aasyn og
som et Skud fra en tør Jordbund. Han havde ei
Skikkelse eller Herlighed, at vi skulde see paa ham,
ikke Anseelse, at vi kunde have Lyst til ham. Han
var foragtet og forladt af Mennesker, en Mand af
Smerter, erfaren i Sygdom, som En, for hvem man
skjuler sit Ansigt; saa foragtet, at vi regnede ham for
Intet. Visseligen, han bar vore Sygdomme og vore
% --- p. 55 ---
\setcounter{footnote}{0}%
\opage{55}Smerter lagde han paa sig; vi agtede ham som den,
der var slagen, ramt af Gud og haardt plaget. Og
han blev gjennemboret for vore Overtrædelsers Skyld
og slagen for vore Synder. Straffen laa paa ham, at
vi skulle have Fred, og ved hans Saar ere vi lægede.“
(Jes. K. 53, 2—5). Det er det aandelige Israel, som
for Prophetens Syn antager personlig Skikkelse, det
er „Jehovas Tjener“, som i Taalmod lider for det
kjødelige Israel og tager dets Overtrædelser paa sig;
det er ham, som i frivillig Lydighed opfylder Loven
og bringer det store Offer, der ene formaaer at udslette
Skylden.

Jehovas Tjener er Jehovas Søn; men den Søn, som
skulde gjenoprette Israels Rige og reise Davids ludende
Hytte,“ (Amos 9, 11) maatte ikke blot have en Tjeners
Taalmod, men ogsaa en kongelig, ja overnaturlig Magt.
„Jeg vil give dig Folkene som fnyse, og Jordens vide
Kreds til din Eiendom; du skal knuse Hedningene med
Jernstav, som Pottemagerkar skal Du slaae dem i
Stykker.“ (Psalm. 2). Under den babyloniske Landflygtighed,
da der slet ingen Udsigter var til en Gjenoprettelse
af Riget ved verdslige Vaaben, maatte Troen
paa det Oververdslige dobbelt forstærkes. Hvad der
ikke er at see paa Jorden, seer da Propheten i Himmelens
Skyer. „Jeg saae i Natten Syner, og see: med
Himmelens Skyer kom En, som en Menneske-Søn, og
kom hen til den Gamle af Dage og stilledes frem for
ham. Og ham blev givet Magt og Ære og Rige, at
alle Folk, Stammer og Tungemaal skulde tilbede ham.
Hans Herredømme er et evigt Herredømme, som
ikke skifter, hans Kongerige gaaer ikke tilgrunde.“
(Daniel 7, 13—14).

% --- p. 56 ---
\setcounter{footnote}{0}%
\opage{56}Hvor ulige disse Messiasbilleder --- Tjeneren i
Ringhed og den mægtige Konge, den sagtmodige Fredens
Fyrste og den overnaturlig udrustede Kriger, den
paa Jorden Lidende og Sønnen af den Gamle i Skyerne
--- hvor ulige disse Billeder end indbyrdes ere, saa ere
Familietrækkene i dem alle dog saa væsentlige, at de
falde sammen i eet Grundbillede, Billedet af ham, der
opfylder Loven.

Men paa det gammeltestamentlige Standpunkt kan
Loven ikke opfyldes. At opfylde Loven er Eet med at
ophæve Loven; det er en Aabenbaringens Hemmelighed,
som Ingen af de gamle Propheter har kunnet gjennemskue.
Lovens Opfyldelse er betinget af Jødedommens
Forvandling, en Forvandling, hvori den gaaer til Grunde.
Saa længe det Overnaturlige og det Naturlige, det guddommelige og
det menneskelige Væsen, vedbleve at staae
saa skarpt mod hinanden som de, til Aandsvilliens
utvetydige Aabenbarelse, netop maatte staae paa Jødedommens
Trin, saa længe maatte Lovens Opfyldelse i
Aand og Sandhed forblive en uopfyldt Fordring. Efter
Udlændighedstiden, da Zion atter saaes i Jehovas Lys
og det andet Tempel var reist, hørte Afgudsdyrkelsen
op og det grove Frafald. Loven oplæstes, Folket
lyttede og forstod, at dets Liv og Selstændighed ufravigelig % PRINTED AS IS: „Selstændighed“
beroede paa en punktlig Lovopfyldelse. Men
det var egentlig kun en Opfyldelse i det Udvortes, en
ceremoniel Præsteopfyldelse, hvortil man naaede:
Aanden var borte, de prophetiske Røster vare forstummede.
Saa levede man hen paa en Overlevering.
Pharisæisme og Sadducæisme stredes om Skriftfortolkningen,
og politiske Idealer fra Makkabæernes Heltetid
blandede sig med de messianske Forventninger. For
% --- p. 57 ---
\setcounter{footnote}{0}%
\opage{57}ikke at løses op i græsk og romersk Hedenskab, holdt
man sig, under Tryk af den verdslige Overmagt udenfra,
til Lovens Bogstav, til Gudstjenestens Former, og Formerne
krystalliserede. Saaledes var den Tilstand, der
faldt umiddelbart sammen med „Tidens Fylde“. Messias
kom. Det aandelige Israel, den snævre Kreds af
de Stille i Landet, hilsede ham i Tro og Haab med
Forventningens ydmyge Glæde; men Præsterne og
Pharisæerne forargedes paa ham, frygtede ham, hadede
ham og saae i ham en Folkeforfører, en Gudsbespotter,
en Forbryder mod Loven. Han, der opfyldte Loven
blev dømt efter Loven og korsfæset som Misdæder: en
sjælerystende Modsigelse, dybt begrundet i Jødedommens
Væsen! Loven opfyldes; men Opfyldelsen er omgiven af
Rædsler. Messias døer: Jødedommens Sol formørkes,
dens Grundvolde skjælve, dens Grave aabne sig, Forhænget
i Templet revner; men --- Loven blev opfyldt,
Spaadommene bleve opfyldte.

% --- p. 58 ---
\setcounter{footnote}{0}%
\lecture{Fjerde Forelæsning.}\opage{58}Forsoningen af det Overnaturlige med det Naturlige,
af Aandens Villie med Menneskets Væsen blev --- paa
overnaturligt Grundlag --- fuldbyrdet af Gudmennesket.
At Gud var i Kristus og forligede Verden med sig selv,
vil sige, at han forligede Menneskets Bestemmelse som
Aand med dets oprindelige Naturbestemmelse. Det er
ikke en Deel af Mennesket, dets Hjerte alene eller dets
Hoved, dets Følelse eller dets Forstand, dets Phantasie
eller dets Villie, men det hele Menneske, Mennesket
som det udgaaer fra Naturens Haand, der skal omdannes.
Hele Naturen, og da i særlig Forstand
Menneskets Natur, er Aandens Værksted, det Grundlag,
hvorpaa den bygger, det Redskab, hvormed den
arbeider, det Materiale, hvortil den trænger, det umiddelbare
Stof, den former og danner. Men, vi gjentage det,
Forsoningen grunder i det Overnaturlige. Kristi Fødsel
er et Under, hans Opstandelse er et Under, Gjenfødelsen,
den nye Fødsel af Vand og Aand, er et Under, som
ingen skriftklog Nikodemus begriber.

Men at Underet i det Nye Testamente viser hen
til en langt inderligere Forening af det Naturlige med
det Overnaturlige end i det Gamle, sees blandt Andet
deraf, at det Prophetiske forandrer sin Karakteer. De
% --- p. 59 ---
\setcounter{footnote}{0}%
\opage{59}gamle Propheter stode under Aanden som almægtig
Villie; den kom over dem, naar den vilde, og drog sig
igjen tilbage, naar den vilde. Propheten fra Nazareth
derimod lever i Aanden, den er „blivende over ham“,
den er som Personlighedens Aand Eet med hans hellige
Villie, som Kjærlighedens Aand inderlig Eet med hans
rene Natur. De Gamle talede i Aandens Navn: \emph{saa}
\emph{siger Jehova}; Faderens Udsendte taler i sit eget
Navn: \emph{Sandelig}, \emph{sandelig siger jeg Eder}.

Paa Lovens Trin er Aandens Rige et overjordisk
Rige, som den almægtige Villie --- vistnok til Menneskenes
Bedste, men dog i egentlig Forstand for sit
eget Herredømmes Skyld --- vil oprette paa Jorden.
Fra dette Synspunkt kan Jehova lignes ved en Konge,
der fører Strid med jordiske Magthavere, og som, endskjøndt
han med eet Slag kunde knuse dem alle, dog
for sit Navns og sin Æres Skyld nedlader sig til at
betjene sig af menneskelige Kræfter. Han hverver da
Tropper, lover de Trofaste, der ville lyde hans Bud,
en høi Løn, og truer de Feige og Trodsige, der overtræde
Disciplinen, med haard Straf. Det er gode Vilkaar,
han byder; Krigerne maae ansee det for en sand
Lykke at stride under hans Banner. Ved at tjene en
saadan Konge tjene de sig selv; thi han har Land
at uddele, og det Land, han har lovet dem, er skjønt:
det flyder med Mælk og Honning. De troe Tjenere
kunne være forvissede om, at han, naar Seiren er
vundet og Landet erobret, vil holde sit Ord og lade
dem leve længe og lykkeligt i Landet. Men Riget,
forstaaer sig, er og bliver hans eget Rige: hans Stridsmænd
ere kun Tjenere. Det er altsaa ikke egentlig
for Rigets Skyld, men for den høie Løn og de gode
% --- p. 60 ---
\setcounter{footnote}{0}%
\opage{60}Vilkaar i Riget, at hans Tjenere stride. Hvis der var
en anden Konge, som bød dem en høiere Løn og et
rigere Land, saa foretrak de naturligviis, hvad da
ogsaa sees deraf, at Mange, der mene ikke at have
faaet deres Lønning og Brød i rette Tid, blive Overløbere,
at tjene en anden Konge. Her er saa vel paa
de Lydiges som paa de Ulydiges Side en endelig Beregning
af Tab og Vinding, et udvortes Forlig mellem
Opoffrelse og Egennytte.

Paa Evangeliets Trin er Menneskenaturen væsentlig
forsonet med Aanden; her passer et andet Billede.
Den Aandernes Fader, der fører Krig med de aandsfjendske
Verdensmagter, har en Søn, han elsker, og
denne Søn er et Menneske, en virkelig Menneskesøn i
Kjød og Blod, og Menneskesønnen er Rigets Arving.
Idet nu Sønnen strider for Udbredelsen af Faderens
Rige, strider han med det Samme for Udbredelsen af
sit eget Rige, og samler Stridsmænd. Men hvad er
det for Stridsmænd? Ikke Leietropper, hvervede Soldater,
der tvinges til Lydighed ved streng Disciplin og
kæmpe for udvortes Løn. Nei, Lønnen er indvortes,
Lydigheden frivillig: Rigets Maal er Menneskenes Frelse.
Saa kalder han dem paa Frelsens og Frihedens Vilkaar.
„\emph{Kommer til mig}!“ Lader os stride for Aandens
Rige, thi dette min himmelske Faders Rige er ogsaa
mit Rige; jeg, Menneskesønnen, er Rigets Arving.
„\emph{Kommer til mig}!“ I, mine Brødre og Søstre, ere
mine Medarvinger; Aandens og Sandhedens og Frihedens
Rige er ogsaa eders Rige, et sandt og virkeligt
Menneskerige. \emph{Kommer til mig}! Lader os stride
for vort eget Rige, at det kan seire, og vi med det. ---
Her er ikke Tale om Lydighed af Frygt eller om
% --- p. 61 ---
\setcounter{footnote}{0}%
\opage{61}udvortes Løn, Egennytten er udelukket; Livet sættes
ind, for at Livet kan vindes. Lønnen er, at indgaae i
Riget; Straffen er, at udelukke sig selv fra Riget.
Saaledes er Aandens Maal udtrykkelig sat som Menneskets
Maal. Menneskets Aandsudvikling gaaer i Eet
med Menneskenaturens Forklarelse i Kraft af det Overnaturlige,
en Forklarelse til Sjæl og Legeme, en Seier
over det Forkrænkelige, en Opstandelse til Uforkrænkelighed
og evigt Liv.

For Prophetierne om dette Riges Udbredelse,
Kampe og Seire afgive Talerne (Mathæus Kap. 24,
Lukas Kap. 21) om Jerusalems Undergang og Verdens
Ende det typisk Afsluttende. Jødedommen gaaer under,
de hedenske Magter oprive hverandre, Løgnens Afgrund
aabner sig og sluger sit Bytte. Men naar Alt, som
ikke har Sandhedens opholdende Grund, styrter og
falder sammen, staaer Sandhedens Aabenbarer forklaret
som Ordet selv, det Livets og Personlighedens, Retfærdighedens
og Kjærlighedens Ord, der skal dømme
Verden. I Johannes' Aabenbaring er denne Verdensomvæltning
videre udmalet i Syner og Billeder og
gammeltestamentlige Symboler. --- Vor Karakteristik
af Aabenbaringspropheterne er hermed sluttet.

Med Aposteltidens Ophør seer det ud, som om
Aabenbaringens Strøm standsede, og det Overnaturlige
med Eet drog sig tilbage; thi vi kunne ikke troe paa
den mirakuløse Aandsmeddelelse, hvorpaa den katholske
Kirke vil begrunde sin Ufeilbarhed. De montanistiske
Propheter har den tidlig forkastet; Bernhard af Clairvaux
var en mystisk Hierark, Joachim af Floras „evige
Evangelium“ fandt ingen Indgang; Birgitta, Nordens
Sibylla, staaer som en Undtagelse. Det er en Præste%
% --- p. 62 ---
\setcounter{footnote}{0}%
\opage{62}kirke, som med sit hierarkiske Inspirationsdogme maa
kvæle al prophetisk Inspiration ved at kvæle Aandens
frie Rørelse. Ikke blot Vantroens frække, udfordrende
Ord, ikke blot Mystikerens og Sværmerens vilde, phantastiske
Ord, men selv det i Sandhed straffende, tugtende
Ord, der vender sig mod dens egen Fordærvelse, maa
den bringe til Taushed: den har brændt en Arnold,
den har brændt en Johan Huss, ja hvad mere er, den
har fordømt og brændt en Savonarola.

Imidlertid kunne Katholicismens Ordførere her
møde os med en vægtig Indsigelse. „Ved at bryde
med den katholske Kirke have I Protestanter, uden at
vide deraf, brudt med Kristendommen og fornegtet
Aabenbaringen, thi I have fornegtet det Overnaturliges
fortsatte Virksomhed. Det er denne Fornegtelse, som
Eders nye Propheter, Propheterne af Geni, nu endelig
gjøre vitterlig. Den Aandsfornegtelse, hvormed de overraske Eder, den Vantro, hvormed de forarge Eder, er
i Grunden kun en renere, tydeligere Udtalelse af Eders
egen Vantro, Eders eget Frafald, Eders egen Aandsfornegtelse
fra Begyndelsen. De have Ret imod Eder
og ville efterhaanden rive Eder alle med sig. Vantroens
Flod svulmer, den bliver stedse bredere og
stærkere; den har Tilløb fra saa mange baade aabenbare
og skjulte Kilder, at Oversvømmelsen snart vil
stige til en sand Syndflod.

I ville gjerne troe, at det Overnaturlige en Gang
i fordums Dage har viist sig og virket paa Jorden;
men i Nutiden tør det slet ikke vise sig, slet ikke virke.
I antage med os, at Aposteltiden var en Miraklernes
Tid, at Aandens Udgydelse paa Pintsefesten var et
kirkestiftende Mirakel; men at den Kirke, som er stiftet
% --- p. 63 ---
\setcounter{footnote}{0}%
\opage{63}ved et Mirakel, ogsaa maa opholdes og styres ved et
Mirakel, ville I ikke antage. Sund Logik maatte dog
kunne lære Eder Eet af To: enten er Aandens Udgydelse
paa Pintsefesten et virkeligt Mirakel, og da er
vor Lære om Aandens Bliven i Kirken og dens overnaturlige
Nærværelse paa de oekumeniske Kirkeforsamlinger
uimodsigelig; eller, hvis vor Kirkes Dogme om
den sig fortsættende Inspiration drages i Tvivl, da maa
Inspirationen paa den første Pintsefest ogsaa drages i
Tvivl. Enten ere Beretningerne om Apostelmiraklerne
troværdige, og da maae de Mirakler, som i Tidernes
Løb ere skete ved de store Vidner i Kirkens Skjød,
ved og med dens Helgene og Helgeninder, ogsaa være
troværdige; eller: dersom de Mirakler, ved hvilke Kirken
endnu den Dag idag forherliges af sine Hellige,
drages i Tvivl, da maae Apostelen Peters Mirakler
drages i Tvivl, Kristi Mirakler drages i Tvivl. Kristi
Himmelfart, Kristi Opstandelse, Kristi Fødsel af Maria,
den rene Jomfru, alt dette maae I Protestanter med
Eders Kultur og moderne Bevidsthed drage i Tvivl, og
hvor ere I saa? I kritiserede Kirken og troede paa
Skriften uden Kritik; I gik videre: I kritiserede Skriften,
for at troe paa Fornuften; I gik videre: I kritiserede
Fornuften, den rene Fornuft, for at troe paa?
--- ja, det er et Spørgsmaal. Eders lovpriste Fremskridt
i fri Oplysning er i Korthed dette: I begyndte
med at skyde Kirken til Side, for at troe paa Evangeliet,
I ende med at skyde Evangeliet til Side, for at
troe paa Kritiken“.

Vi have anført disse Indvendinger, ikke for at
forsvare os imod Katholikerne, men for at komme til
Klarhed over Protestantismens Væsen og dens principielle
% --- p. 64 ---
\setcounter{footnote}{0}%
\opage{64}Stilling til Kultur og Videnskab. Katholikernes Paastand,
at den protestantiske Kirke nødvendig maa ende
i Selvopløsning, finder stærk Medhold hos den moderne
Bevidstheds Ordførere. „Saa længe Protestantismen
endnu ikke har givet Slip paa den sidste Rest af
gammel Tro paa Kristendommen som overnaturlig
Aabenbaring, saa længe vil den som en middelalderlig
Levning vedblive at være i Strid med Fornuft og Kultur.“
--- Den Verdensanskuelse, som svarer til Aabenbaringstroen,
er, siger man, \emph{dualistisk}, Nutidens
fornuftige Verdensanskuelse er derimod \emph{monistisk}.
Slagordene ere: \emph{Dualisme} og \emph{Monisme}.

For at den dannede Almeenhed ikke skal have
Bryderi med at udfinde Betydningen af disse dunkle
Ord, giver man en kort og fattelig Forklaring. Ordet
\emph{Dualisme} betyder alt det Dumme, det Uklare, Ulogiske
og Obscurantistiske, man kan forestille sig; Ordet \emph{Monisme},
der lyser som Stjerne paa Physik-Philosophernes
Stormhat, betyder det Kloge, det Klare, det logisk
Sammenhængende, det, som alene kan tilfredsstille den
frie Tanke. Efter denne Forklaring er det intet Under,
at de Dannede, der ville holde Skridt med Tiden,
vende Dualismen Ryggen for i rask Trav at følge den
paa Stormhatten lysende Monisme.

„Hvis Monismen ikke var Sandheden, da saae det
sort ud med Fremtiden. At aflede Naturen af Aanden,
hævde Aandsvillien og sætte Friheden øverst, er dualistisk
og gammeldags; at aflede Aanden af Naturen,
forgude Driftvillien og sætte Nødvendigheden øverst,
er monistisk og ægte moderne. Hvis den moderne
Bevidsthed ikke var Eet med den moderne Natur, hvis
det lod sig tænke, at Stormhattene kunde blæse af for
% --- p. 65 ---
\setcounter{footnote}{0}%
\opage{65}de physik-philosophiske Genier, som de laadne Huer
paa Marken blæse af „Melkebøtterne“, hvad skulde det
saa blive til med Frihed og Fremskridt? Skulde Aandsvillien
virkelig være høiere end Driftvillien, og Nødvendigheden
magtes af Friheden, saa var Menneskehedens
Sag forloren, Alderdommen uden Trøst, Manddommen
uden Mod, Ungdommen uden Haab. Vover
man at tvivle om Monismen, saa vover man at tvivle
om Darwins storartede Hypothese, det eneste videnskabelige
Dogme, om hvis ubetingede Gyldighed den
frie Tanke endnu ikke har vovet at tvivle. Kunde
man, var det muligt, kunde man tvivle om Menneskenaturens
monistiske Eenhed med Abenaturen, Menneskeformens
monistiske Afstamning fra Abeformen --- skrækkelige
Tanke! saa maatte man jo fortvivle om Livets
Maal; thi, naar Mennesket ikke kjender sin Afstamning,
hvorpaa skulde det da kjende sit Maal, naar det
ikke kjender sin Fødsel, hvorledes skulde det kjende
sin Fødsels Adel? Dersom Nødvendigheden ikke var
den seende Friheds blinde Moder, dersom Menneskevilliens
gigantiske Anstrængelser med Dyder og Lyder,
historiske Bedrifter og poetiske Forbrydelser ikke monistisk
lod sig tilbageføre til den monistiske Naturs
monistiske Moderskjød: hvad skulde det da blive til
med de naturlige Menneskerettigheder, med Livets
Harmonie i den endeløse Kamp for Tilværelsen!“ ---
Dette er i Grundtræk den moderne Bevidstheds monistiske
Program.

Den for Dualisme udskregne Verdensanskuelse er
for Resten ikke Dualisme, nemlig ikke Dualisme i den
Forstand, at Tilværelsen skulde falde fra sig selv i to
Stykker, et overnaturligt og et naturligt; Aabenbarings%
% --- p. 66 ---
\setcounter{footnote}{0}%
\opage{66}troen er, efter sit Princip, ikke dualistisk i den Forstand,
at Naturlovenes System skulde falde uden for
Aandslovenes, og desaarsag nu og da vilkaarlig brydes og
vilkaarlig heles af en overnaturlig Almagtsvillie; en
saadan Fremstilling er en Forvrængning, der kun skal
tjene til at anbefale og forherlige Monismen. Den
Dualisme, hvori den protestantiske Aabenbaringstro er
begrundet, har Eenheden af Aand og Natur, Frihed og
Nødvendighed, Gud og Verden til Forudsætning. Denne
oprindelige, altomfattende Eenhed er en Eenhed af
stærke Modsætninger; men netop derved en stærk og
dyb og indholdsrig oprindelig Eenhed. Eenheden af
Aand og Natur, af Gud og Verden, er en Frihedens
Eenhed, en almægtig Villieseenhed. Herved bestemmes
Forholdet mellem Viden og Tro.

Protestantismen anerkjender den frie Tankes Ret
i videste Udstrækning, d. v. s., den anerkjender Alt,
hvad grundig Undersøgelse formaaer at stadfæste; den
prostesterer mod Alt, hvad der beviislig er huult og % PRINTED AS IS: „prostesterer“
grundløst. --- „Er der en levende Gud, der faderlig sørger
for sine Skabninger, særlig for Mennesket, eller er det
Alt frembringende Væsen i uendelig Nyden af sine
Frembringelser uendelig ligegyldigt mod de levende
Individers Vee og Vel?“ Intet Menneske veed det,
ingen videnskabelig Undersøgelse formaaer at afgjøre
det. „Ere Sjælene udødelige eller forsvinde de med
Legemets Opløsning som Lys, der udslukkes?“ Ingen
veed det, intet videnskabeligt Beviis kan afgjøre det.
Men naar Videnskaben ikke kan afgjøre, hvilken af de
modsatte Antagelser der er den rette, naar den ene af
dem, ifølge Videnskabens Dom, lige saa vel er mulig
som den anden, saa er det jo Daarskab at fordre
% --- p. 67 ---
\setcounter{footnote}{0}%
\opage{67}videnskabelig Afgjørelse af det, som Videnskaben ikke
kan afgjøre. --- Men dersom saa Nutidens Dannede i
Følge heraf vilde være enige om, at det, som ikke kan
vides, heller ikke kan troes, saa vilde Protestantismen
protestere mod alle Nutidens Dannede. At tvivle om
Alt er Fortvivlelse, at leve paa løse Meninger er Løsagtighed;
Livets Alvor kræver Besvarelse af Livets
Spørgsmaal, en Afgjørelse af dets høieste Anliggender.
„Men hvem skal besvare de store Spørgsmaal, hvem
skal klare de høieste Anliggender?“ Mennesket selv og
ingen Anden. „Hvor mange Feiltagelser og Vildfarelser
ere da her ikke mulige?“ Her er kun een Feiltagelse
mulig, een eneste, som bestaaer i, at Mennesket tager
Feil af sig selv; her er kun een Vildfarelse mulig, een
eneste, og den bestaaer i, at Mennesket farer vild i
sig selv.

Ved at gjøre Menneskelivets høieste Spørgsmaal
til personlige Spørgsmaal og indrømme Mennesket Evne
og Myndighed til selv at besvare disse Spørgsmaal, har
Protestantismen gjort Troen til en i høieste Forstand
personlig Sag og lagt Tilværelsens Tyngdepunkt ind i
Personligheden. Men, vil man spørge, hvad har det
med den positive Religion at gjøre? Her svares: Personlighed
er betinget af Selvforstaaelse, men Selvforstaaelse
er betinget af Underviisning, Opdragelse og
Vækkelsesmidler, som den positive Religion tilbyder.
Først synes det, som hvilede Troen paa dens Autoritet
alene, men efterhaanden som det gaaer frem med den
personlige Udvikling, vendes Forholdet om, saa at Religionens
Autoritet i sidste Instans kommer til at hvile
paa den personlige Selvforstaaelse. Troesprincipet
falder da sammen med Selvforstaaelsens Princip:
% --- p. 68 ---
\setcounter{footnote}{0}%
\opage{68}Autoriteten er kun en Støtte, et Rækværk i Høiden, at
den Opadstigende ikke skal svimle. Vil Nogen mene,
at en saadan Opfattelse af personlig Tro umulig kan
bestaae med den Tro paa en historisk Aabenbaring,
som den protestantiske Kirke har vedkjendt sig og
endnu bestandig maa vedkjende sig for at være Kirke:
da lægge man vel Mærke til, at det ikke er Selvforstaaelsen,
der frembringer Religionen, men omvendt.
Der er en himmelviid Forskjel imellem at aabenbare
den personlige Sandhed og at anerkjende Sandheden i
den sig aabenbarende Personlighed. Aandsvillien er
Personlighedens Centrum; er Aandsvillien i sit Væsen
kun relativ, en forgjængelig Størrelse, der opstaaer i
Tiden og forgaaer i Tiden, da har det Personlige kun
et relativt Værd, og Friheden gaaer under i Nødvendigheden.
Skal Friheden seire og Personligheden gjælde
ubetinget, da maa Aandsvillien selv være en Villie, der
bærer den hele Tilværelse, en almægtig Villie, der udtaler
sig saa vel i Naturens som i Aandens Love.
Kun det Menneske, hvis personlige Selvforstaaelse gaaer
saa dybt, at han erkjender sin afmægtige Villie beaandet
og bevæget af den almægtige Villie, naaer til Personlighedens
Ankergrund.

Det vil heraf forstaaes, med hvilken Varsomhed de
indre, individuelle Erfaringer maae behandles, hvor omhyggelig
de maae sammenlignes med de ydre, hvilken
Udholdenhed i psychologisk Iagttagelse og Selvkritik
der hører til, for at den protestantiske Aabenbaringstro
skal kunne siges at være en fuldt bevidst Tro, en Tro,
der kan gjøre Rede for sin Oprindelse og sin Karakteer.
Den protestantiske Bevidsthed paa sit nuværende Kulturtrin
finder sig ingenlunde i umiddelbar og sympathetisk
% --- p. 69 ---
\setcounter{footnote}{0}%
\opage{69}Forstaaelse med de Aandsphænomener, den møder i
Aabenbaringen. Tvertimod; jo mere udvortes historisk
disse Phænomener opfattes, jo mere frastødende ere de.
Først idet det Udvortes føres tilbage til sit Indvortes,
begynder Forstaaelsen. Det er Aandsvilliens Forhold
til Menneskets Natur, der afgiver Forstaaelsens Grundlag.

Den protestantiske Kulturbevidsthed er et ideelt
Produkt af mangfoldige Omdannelser, langvarige, forudgaaende
Arbeider; i disse Arbeider har den sin Historie;
i denne Historie maa den søge Forudsætningerne for
en Forstaaelse af sig selv. Det Samme gjælder om den
protestantiske Troesbevidsthed. Den har i den aabenbarede
Religion og de paafølgende Opfattelser af dens
Karakteer og Væsen sin Historie og i denne Historie
Forudsætningerne for Forstaaelsen af sig selv. Religionen
er et Gudsforhold. At dette Gudsforhold ikke
kommer i Stand derved, at Mennesket selv søger og
finder en Gud, det kan Kulturbevidstheden bevidne.
Den er jo netop forvisset om, at man kan søge Gud
indefter og udefter, paa Jorden og i Himlen og hvor
man vil, og finder ham dog ingensteds. Kritiken gjør det
aabenbart, at det Gudsforhold, der vilkaarlig sættes af
Mennesket selv, er subjektivt og grundløst, at den Gud,
der frembringes af Menneskefølelse, Menneskephantasie
og Mennesketanke, lige saa vist maa blive en Afgud
som de Guder, der frembringes af Menneskehænder.

Saa sandt nu Gudsforholdet er Personlighedens
Grundforhold, og dette Grundforhold ikke kan have
Realitet, med mindre det oprindelig er sat af Gud, saa
vil den protestantiske Personlighed, der er vis paa sig
selv, holde sig forvisset om, at Gudsvillien, den almægtige
Aandsvillie, vedblivende staaer i et virksomt Forhold
% --- p. 70 ---
\setcounter{footnote}{0}%
\opage{70}til Menneskevillien: dette virksomme Forhold er Underet.
Underet, saaledes forstaaet, er ikke det Samme som
Miraklet. Al Væren er Under (alles Sein ist Wunder),
dette Ord af en berømt Physiker anvende vi her i den
Mening, at en guddommelig Tænken, en guddommelig
Villen gjennemtrænger den hele materielle Verden, at
Materiens Kræfter netop virke som naturlige Aarsager,
idet de virke efter deres naturlige Anlæg, at disse naturlige
Anlæg maae være at henføre til en guddommelig
Villen, at Tingene, med andre Ord, virkeliggjøre Skaberens
Villie ved at virke efter deres Natur. Under
Forudsætning af, at Aandsvillien er virksom i den hele
Natur, maa den da ogsaa være virksom i Mennesket,
og dens Virksomhed sees i Menneskenaturens eiendommelige
Anlæg.

Fra dette Synspunkt kan den protestantiske Personlighed
til religiøs-psychologisk Selvforstaaelse, uden
at komme i Modsigelse med Videnskaben, klare sig de
historiske Forudsætninger, paa hvilke hans Tro hviler.
At der i Menneskets Natur maa ligge en oprindelig
Trang til Religion, sees allerede af Mytherne. I det
mythiske Gudsforhold er det Guddommelige lige saa
vel virksomt som det Menneskelige. Den Tid er forbi,
da man ansaae de mythiske Guder for vilkaarlige Opfindelser.
Men under Mythens Form er det Guddommelige
skjult i det Naturlige; først paa Aabenbaringens
Trin skiller Aanden sig fra Naturen, som den
seende Frihed fra den blinde Nødvendighed; Gudsforholdet
sættes af Aanden selv som et Villies Forhold,
og Bevidstheden om det Personlige bliver til for
Mennesket.

% --- p. 71 ---
\setcounter{footnote}{0}%
\opage{71}Hvorledes det nu egentlig er gaaet til, at den guddommelige
Aand i Tidernes Løb, under Udviklingens
Gang, har forandret sin Villen til Villie, hvorledes den
evige Frihed er traadt op imod Naturnødvendighen og % PRINTED AS IS: „Naturnødvendighen“
har sat sin Selvaabenbaring i Mythens Sted, det er
Spørgsmaal, hvis Besvarelse ligger udenfor Videns
Grændse. Den Indvending, at ingen Forandring kan
foregaae med den guddommelige Virkemaade, eftersom
Gud jo er og skal være uforanderlig, er et Ræsonnement
i Luften; det ligger netop i Villiens Natur at
kunne forandre sin Virken, uden at forandre sit Væsen.
Mysteriet er det uudgrundelige Forhold mellem Frihed
og Nødvendighed; vi maae holde os til det Religiøs-Psychologiske.


Idet den protestantiske Personlighed føler sig frastødt
af Aandsmiraklet, opdager han i det umiddelbart
Frastødende tillige noget væsentlig Tiltrækkende. Det
er den talende Gud, Gudsvilliens overnaturlige Tilkjendegivelse,
der frastøder; og hvorfor? Fordi en saadan
Tilkjendegivelse falder udenfor enhver Erfaring i
Nutiden og forudsætter en sjælelig Tilstand, der længst
er forsvunden. Men i dette Fremmede er der ham dog
noget Bekjendt, noget med hans Tilstand inderlig Beslægtet.
Paa sit eget Bevidsthedstrin kan Personligheden
ikke i og ved sig selv opdage den overnaturlige
Villies Indvirken; han er i sine indre Oplevelser henviist
til en almindelig Forudsætten af guddommelig
Tænken og guddommelig Villen; han kommer ikke og
kan ikke komme videre, thi hans Gudsforhold er reflecteret.
Hvad opdager han saa? At dette reflecterede
Gudsforhold er en Speiling af hiint oprindelige, og
den religiøse Beaandelse en psychologisk Eftervirkning
% --- p. 72 ---
\setcounter{footnote}{0}%
\opage{72}af Aandsmiraklet. Han troer ikke paa Aandsmiraklet,
fordi han troer paa en Overleverings Autoritet; men
han troer paa Overleveringens Autoritet, fordi han i
Selvforstaaelse erkjender, at det perennerende Under,
hvori hans Gudsforhold grunder, den Villiens Tilvæxt,
hvorved hans Selvstændighed voxer, den Lysæther,
hvori hans Frihed aander, indeslutter Aandsmiraklet
som sin Tilblivelses Mulighed, sin oprindelige Forudsætning.
Saaledes forstaaet, fører det Afledede til fortsat
Fordybelse i det Oprindelige.

Den protestantiske Aabenbaringstro er væsentlig
forskjellig fra den katholske. Forskjelligheden kommer
kun deelviis tilsyne, saa længe der fornemmelig strides
om, hvad det er, \emph{hvorpaa} der skal troes, \emph{hvormeget} der skal troes o. s. v. Grundforskjelligheden
viser sig først, naar Spørgsmaalet er, \emph{hvorledes} der
skal troes, og hvad Troen er i sit Princip.

Den katholske Tro er umiddelbar Autoritetstro,
dens Mirakeltro en aandløs Tro, og den aandløse Tro
Overtro. Den katholske Kirke har ikke blot efterlignet
Jødedommen i Udvorteshed; den er med sit uhyre
Overleveringsapparat, sine Læreformer, sin Kultus, sin
Kirketugt og frem for Alt ved sin forviklede Verdenspolitik
gaaet langt videre. De som elske de sætningsvise
Spaadomme og lade Propheten „skue i Aanden“
det Fremtidige, naar han taler om det Datidige, kunne
med Rette paastaae, at det er den katholske Udvorteshed,
Jesaias har skuet i Aanden, idet de med Eftertryk,
ja med ti Gange forstærket Eftertryk her kunne
gjentage Jehovas Ord: „Jeg er kjed af Eders mange
Offere\dots{} Eders Maaneder og Eders bestemte Høitider
hader min Sjæl; de ere mig en Byrde.

% --- p. 73 ---
\setcounter{footnote}{0}%
\opage{73}Omsider, da Misbrugen var paa sit Høieste, kom
det til et Udslag fra Grunden, og Aandskræfterne rørte
sig. Men Gjennembrudet var ikke noget mirakuløst
Udslag; det var ingen ny Aabenbaring, det var en Tilbagevenden
til det Oprindelige, det var Reformationen.
Luther er en Prophet, men ikke i Kraft af et Aandsmirakel.


I Luthers Beaandelse træder det Overnaturlige ikke
særlig frem; det er tvertimod skjult under lutter Naturbetingelser.
Ved Bedømmelsen af Luthers storartede
Karakteer gjælder det at opfatte ham paa een Gang
som en Aandens Mand og en Mand af Geni. Hos en
Prophet som Jesaias, en Evangelist som Johannes, en
Apostel som Paulus, er Naturbegavelsens Storhed umiskjendelig:
men derfor vil man dog ikke opfatte Jesaias,
Johannes og Paulus som Genier. Og hvorfor ikke?
Fordi det Naturlige hos dem er saa ganske i Aandsvilliens
Tjeneste, at det ikke lyser med sit eget Lys,
men nærer Aandens Flamme og kun giver den et
vexlende Skjær. Hos Luther er Forholdet omvendt,
Aandsvillien er her og rører sig kraftig. Luther har
Mod, fordi han har Tro, og han har Tro, fordi han
dybt er grebet af Aanden. Man læse f. Ex. hans Brev
til Churfyrsten, da han forlod Wartburg: „Ed. Ch.
Naade veed, eller veed De det ikke, saa gjør jeg det
Eder her med vitterligt, at jeg har Evangeliet, ikke fra
Mennesker, men alene fra Himlen ved vor Herre Jesus
Christus\dots{} Men at jeg har tilbudt mig at ville stille
mig til Forhør og Dom, det er skeet, ikke fordi jeg
tvivlede, men af Ydmyghed for at lokke de Andre.
Men da jeg nu seer, at min altfor store Ydmyghed vil
geraade Evangeliet til Fornedrelse, og at Djævelen vil
% --- p. 74 ---
\setcounter{footnote}{0}%
\opage{74}indtage Pladsen aldeles, hvis jeg kun indrømmer ham
en Haandsbred, saa maa jeg, tvungen af min Samvittighed,
gribe Sagen anderledes fat. Jeg har viist Ed.
Ch. Naade Lydighed nok derved, at jeg, for at tjene
Ed. Ch. Naade, dette Aar har holdt mig tilbage\dots{}
Det veed jeg med mig selv, at om denne Sag stod saaledes
i Leipzig, som den staaer i Wittenberg, vilde jeg
alligevel ride derind, om der endog (Ed. Ch. Naade
tilgive mig min naragtige Tale) i ni Dage regnede lutter
Hertug Georger, og hver af dem var ni Gange saa
rasende som denne\dots{} Saadant være skrevet til Ed.
Ch. Naade i den Hensigt, at Ed. Ch. Naade maa vide,
at jeg kommer til Wittenberg under en langt høiere
Beskyttelse end Churfyrstens. Jeg har heller ikke i
Sinde at begjære Ed. Ch. Naades Beskyttelse, ja jeg
holder for, at jeg snarere skulde beskytte Ed. Ch.
Naade, end De kunde beskytte mig. Hertil kommer,
at dersom jeg vidste, at Ed. Ch. Naade kunde og vilde
beskytte mig, saa vilde jeg ikke komme. Denne Sag
skal og kan ikke hjælpes og fremmes ved Sværdet;
Gud maa her alene raade uden al menneskelig Omsorg
og Bistand. Derfor er den, som er stærkest i Troen,
her den bedste Beskytter; og da jeg nu sporer, at Ed.
Ch. Naade er heelt svag i Troen, kan jeg ingenlunde
holde Ed. Naade for den Mand, der skulde kunne beskytte
eller redde mig“.

Naar man betænker, at Luther, da han skrev
dette Brev, var en fredløs Mand, og at Kurfyrst Frederik
den Vise var den Eneste, som, menneskelig talt, kunde
tage ham under sin Beskyttelse, saa faaer man Indtryk
af en Tro og et Mod, der minder om de gamle
Propheter.

% --- p. 75 ---
\setcounter{footnote}{0}%
\opage{75}Ikke desto mindre er Aandsvirksomheden eiendommelig
betinget af hans individuelle Natur. Hans
Sjælekriser, hans Tidsalders Vilkaar, Folkestammens
Gemyt, hans storartede Lidenskab, de stærke Lyn fra
det aandelige Uveir, hvori han kommer, alt i Forening
tilhører en kulturhistorisk Atmosphære. der ombølger
hans Skikkelse og viser os ham som en Ideens Mand,
en Mand af Geni. Hvor uafhængig Luther end var af
alle Forstandsberegninger og endelige Klogskabshensyn,
saa formaaede han dog med al sin Troeskraft og sin
religiøse Karakteers Redelighed ikke at fastholde sin
evangeliske Frihedstanke i dens oprindelige Reenhed;
han maatte omsider bøie af for Verdensforholdene og
finde sig i Skranker, han ikke var i Stand til at
gjennembryde.

Mærkeligt er det at see, hvorledes det præstelig-theologiske
Element i Luther lidt efter lidt faaer Overvægt
over det prophetiske. Zwickauer-Propheter, Bondeoprør,
Kjættere og Sværmere af alle Farver viste i
tydelige Tegn, at de vilde Aander vare løsladte. „Da
vi“, siger Luther, „ved Hjælp af Evangeliet begyndte
at lære, at udvortes Ting ikke kunde gjøre salig,\dots{}
saa gik disse Folk, ogsaa store og lærde Folk, saa vidt
at Daaben som et udvortes Vand, Ordet som en udvortes,
menneskelig Tale, Skriften som et udvortes
Bogstav, gjort af Blæk, Viin og Brødet som baget af
Bageren, skulde være aldeles Intet, thi de vare udvortes,
forgjængelige Ting; saaledes faldt de paa Skriget:
Aand! Aand! Aanden maa gjøre det, Bogstaven ihjelslaaer!
--- Men „Paulus“ (1. Cor. 15, 3—4) anseer
Skriften for sit stærkeste Vidnesbyrd og viser, at der
ikke er noget Middel til at holde vor Lære og Tro, uden
% --- p. 76 ---
\setcounter{footnote}{0}%
\opage{76}det legemlige eller skriftlige Ord, forfattet med Bogstaver
og mundtlig prædiket ved ham eller Andre.
Thi her staaer klart: Skrift, Skrift. Men Skrift er
ikke lutter Aand, hvorom de sluddre, at Aanden maa
gjøre det alene, at Skriften er et dødt Bogstav og ikke
kan give Livet. Gjør ikke meget Væsen af Aanden,
naar Du ikke har det udvortes aabenbare Ord, thi det
vil visselig ikke være en god Aand, men den lede Satan
fra Helvede. Thi den hellige Aand har jo indesluttet
sin Sandhed, sine Raad og alle Hemmeligheder i Ordet
og aabenbaret dem i Skriften, saa at Ingen kan undskylde
sig eller har at søge eller forske efter noget
Andet“.

Saa gjerne vi end indrømme Gyldigheden af disse
Ord, naar de ret blive forstaaede, saa meget vi end
paaskjønne Luthers og Reformatorernes Bestræbelser
med at sikkre Kirken en Troesbekjendelse og bygge
Lærebegrebet paa et fast Grundlag, saa maae vi dog
udtrykkelig fremhæve, at det netop var paa et Feilgreb
i disse Bestræbelser, at Gjennemførelsen af Protestantismens
Princip strandede. Hvad der levede i
Luther selv, men aldrig gik op for Luthers Bevidsthed,
det store Problem, hvorledes det egentlig forholder sig
med Ordet og Aanden, med Skriftens Autoritet og
Sandhedens frie Tilegnelse, det kunde endnu mindre
blive klart for hans ivrige Tilhængere. Man higede
mod en ny Tilstand og arbeidede for et Bestaaende;
man trængte til Orden og Tugt i det Udvortes, og af
Skriftens Ord og Skriftens Lære, hvis Elementer den
reformatoriske Hede havde smeltet, og de vilde Aander
blandet med urene Bestanddele, udkrystalliserede de
lutherske Dogmer til juridisk Brug og dannede et fast
% --- p. 77 ---
\setcounter{footnote}{0}%
\opage{77}Grundlag for den lutherske Statskirke. Atter Autoritetstro,
atter Udvorteshed!

Den protestantiske Udvorteshed var i det syttende
og attende Aarhundrede lige saa aandløs som den
katholske, og langt mere oprørende, fordi den stod i
skrigende Modsigelse med Principet. Ligesom de
glødende Masser og hede Dampe i Jordens Indre under
Tryk af den faste Skorpe tiltage i Spændkraft, indtil
Skorpen revner, og der skeer et vulkansk Udbrud, saaledes
vandt den forskende Aand under det tunge Tryk
af den udvortes Autoritet omsider Spændkraft til at
bryde Skorpen og skaffe sig Luft. Efter enkelte, spredte
Eruptioner slaae Frihedens Flammer ud i en Kjæde af
Vulkaner, Rækkevulkaner, hvis Ild umulig kan slukkes;
thi naar man tilstopper Krateret i den ene, ere de
andre i fuld Virksomhed: i Lueskjær af disse Vulkaner
møde vi \emph{de nye Propheter}.

% --- p. 78 ---
\setcounter{footnote}{0}%
\lecture{Femte Forelæsning.}\opage{78}Den franske Revolution af 1789 var ikke et Oprør
iblandt mange; det var en Grundvoldenes Revolution,
et verdenshistorisk Jordskjælv. Rystelserne strakte sig
vidt og sporedes overalt, ikke alene i det Udvortes,
men ogsaa, og det fornemmelig, i det Indvortes. At
møre Støtter faldt, at Throner væltede med gamle Dynastier,
fornam man med Skræk i den synlige Verden;
men ogsaa i den usynlige, Sindenes, Tankernes, Ideernes
Verden var der Dynastier som faldt: Autoriteten faldt,
Legitimiteten faldt, Pieteten faldt, Kristendommens ærværdige
Majestæt blev traadt under Fødder. Det var
Driftvillien fra den blinde Nødvendigheds Afgrund,
der løftede Frihedens Banner og reiste sig med sin
hele Hær af selvisk flammende Lidenskaber. Der foer
en Ildstrøm gjennem de vilde Sjæle, og Menneskehoveder
meiedes ned som Tidselhoveder. Revolutionens Dæmon
rystede de blodige Vinger og svor, ikke at raste, før
alle Frihedens Skranker vare gjennembrudte, Forrettighedernes
Bolværker sløifede, og Tyranniets Taarne
jævnede med Jorden.

Da Uveiret omsider havde udraset, da de gamle
Altere med den gamle Tro, de gamle Throner med de
gamle Myndigheder igjen havde reist sig, var der vist%
% --- p. 79 ---
\setcounter{footnote}{0}%
\opage{79}nok Mange, som meente, at en Gjentagelse af slige rædselfulde
Katastropher kunde forebygges, og Ødelæggelsernes
Spor udslettes, naar de legitime Magthavere
toge Guds Naade til Hjælp, og Kirke og Politi forenede
sig om at anvende større Aarvaagenhed ; men Erfaringen
har hidtil ikke stadfæstet deres Forventning. Vi skulle
her ikke omhandle de politiske Forviklinger, men udelukkende
fæste vor Opmærksomhed paa „Kulturkampen“,
den Kamp, der føres mellem det Gamle og det Nye
med ideelle Vaaben. Det Gamle har ingen levende
Propheter, alle dets Propheter ere døde; men --- de nye
Propheter ere levende.

Hvilke ere da disse nye Propheter? Hvorledes see
de ud? Hvorfra komme de? Hvad ville de? Hvor er
det Maal, hvortil de pege hen? Det er vanskeligt at
besvare disse Spørgsmaal nøiagtig. Mængden er
broget, dens Tal er neppe til at opregne. Forskjellige
Physionomier, Masker, Kostumer: hvor ere de dog
forskjellige i Begavelse, Dannelse, Optræden og Holdning!
De komme jo fra alle Livets Kredse, fra Arbeidernes
Værksteder, fra Politikens Leire, fra Agitationernes
Ildsteder og fra Videnskabens stille Rige; vi
finde dem i Literaturen, vi høre dem i Folkeforsamlingerne;
de glimre i praktiske Eventyr, vi møde dem
i Romaner og Skuespil. En Byron og en Darwin, en
St. Simon og en George Sand, en August Comte og en
Hartmann, en Louis Blanc og en Alfred de Musset, en
Proudhon og en Stuart Mill, en Fourier og en Lasalle,
hvor høist ulige, ja hvor indbyrdes uenige! Ikke desto
mindre høre de alle under samme Kategori; de staae alle
paa Revolutionens Grund. De ere trods deres indbyrdes
Afvigelser enige i Eet, og dette Ene er afgjørende: det
% --- p. 80 ---
\setcounter{footnote}{0}%
\opage{80}Overnaturlige maa bort, kun det Naturlige gjælder;
Aandsvillien er en Skygge, kun Driftvillien har Kjød
og Blod.

Uden at indlade os paa Forskjellighederne i det
Enkelte, skulle vi her indskrænke os til, hvad der i
denne Sammenhæng vedkommer vor Opgave, en Karakteristik
af den fælles Hovedtanke, og til nærmere Belysning
anføre nogle Exempler.

\emph{En ny Religion}. Som den Første i Rækken
maae vi nævne St. Simon, Stifteren af den Religion,
der kaldes St. Simonisme, et karakteristisk Produkt
fra Revolutionens Arnested. Naar her tales om Religion,
maa det vel mærkes, at det er Religion i moderne Forstand,
en Religion uden Gud. --- Claude Henri, Greve
af St. Simon (født 1760, død 1825) hørte til den høieste
Adel i Frankrig --- han vilde nedstamme fra Karl den
Store --- og førte et eventyrligt Ungdomsliv; var Lafayettes
Vaabenfælle i den amerikanske Frihedskrig o. s. v.
Afvexlende rig og fattig, arbeidsom og ørkesløs, ærgjærrig
og udsvævende, satte han sig som Maal: at frelse den
faldne Menneskeslægt og at gjøre det bedre end Kristus.
Kristus havde kun sørget for det Aandelige og overseet
det Kjødelige, hvilket var en stor Eensidighed; St. Simon
vilde sørge baade for det Aandelige og det Legemlige
ved at indsætte Kjødet i dets Rettigheder (Nouveau
Christianisme). Istedenfor et hiinsidigt uvirkeligt Himmeriges
Rige, vilde han stifte et jordisk virkeligt Arbeiderrige,
helliggjøre Industrien, bearbeide Guddommen
ved at bearbeide Stoffet og forbedre de Fattiges Kaar.
Han sagde saaledes: „Hvis den Dag idag Døden bortrev
alle Prindser af Blodet, alle Kronens høie Embedsmænd,
Statsministrene, Præsidenter, Biskopper og 10,000
% --- p. 81 ---
\setcounter{footnote}{0}%
\opage{81}af Frankrigs største Godseiere, vilde det vække Sorg,
fordi de ere høitstaaende Personer; men Staten vilde
ikke rokkes et Haarsbred, og Dagen efter vilde der
være bødet paa disse 30,000 Piller, efterdi Tusinder af
Mennesker ere i Stand til at udrette det, som Prindser
af Blodet, Ministre, Rigmænd og høie Prælater have at
gjøre i Verden. Men hvis de, der saaledes døde, vare
de vigtigste Haandværksmestre, de største Arbeidsgivere,
Landets Kemikere, Physikere, Kunstnere og Digtere,
saa vilde Tabet af disse maaskee kun 3000 Personer
slet ikke kunne oprettes.“ --- St. Simon blev af sine
Disciple sat høit over Kristus. Han levede for sin
Idee, men døde dog ikke som Kristus, thi hans Død
var en Følge af et mislykket Selvmord. Til Disciplene
der stode ved hans Leie, talede han disse prophetiske
Ord: „Pæren er modnet, I ville plukke den.“

Og de plukkede „Pæren“, men bleve snart uenige.
Det var fornemmelig Ægteskabet og Kvindens Emancipation
der fremkaldte Striden. En af Ypperstepræsterne
i det St. Simonistiske Hierarkie, Père Enfentin, havde
nemlig opdaget, at der i Menneskeslægten gives to
Slags mandlige Individer. Der er nemlig Individer med
saa bestandige Følelser, at de, naar de engang ere
blevne forelskede, med en fast haardnakket Udholdenhed
hænge ved een og samme Person; det er „Othello-Naturer“,
eiendommelig anlagte paa livsvarigt Ægteskab
(le mariage définitif). I Modsætning til dem er der
andre Individualiteter, livlige bevægelige Naturer, der
sværme fra Skjønhed til Skjønhed --- ridderlige Sommerfugle,
der ikke blive hængende ved en enkelt Blomst,
for ei at gjøre de andre Blomster Uret! --- det er
„Don-Juan-Naturer“, som kun ere oplagte til forbi%
% --- p. 82 ---
\setcounter{footnote}{0}%
\opage{82}gaaende Ægteskab (mariage provisoire), Père Enfentin
tilstod aabenhjertig, at han hørte til de sidste. Forholdt
det sig nu saaledes med de mandlige Individer,
saa var der jo grundet Formodning om, at det maatte
forholde sig paa lignende Maade med de kvindelige;
kun at disse ikke, paa Grund af langvarig Underkuelse,
turde give deres medfødte Natureiendommelighed saa
ligefrem tilkjende. Père Enfentin fandt desaarsag, at
det var hans Pligt som Ypperstepræst at undersøge
Sagen lemfældig og føle sig frem, for at opdage „den
frie Kvinde“.

\emph{Nysædelige Bevæggrunde}. Den gamle Moral
pleier gjerne, naar Talen er om Dyd og Pligt, at henvise
ethvert Menneske til dets egen Samvittighed og
minde det om, at betvinge de farlige Lidenskaber. Men
fra det moderne Standpunkt maa denne Moral forkastes.
Drifterne, Tilbøielighederne, Lidenskaberne skulle ikke
betvinges, de skulle tvertimod løslades til fri og alsidig
Vexelvirkning. Ved at undertvinge Drifterne, lænkebinde
Lidenskaberne og afstumpe Tilbøielighederne
berøver man Individet dets medfødte Motiver, de kraftigste
Bevæggrunde til sædelig Virksomhed. Mennesket
kan ikke gjøre Andet end hvad det har Lyst til, naar
det skal gjøres til Gavns. Den iblandt de nye Propheter,
som med størst Energie har udtalt dette, er Charl.
Fourier (død 1837). I sin Lære om de fire Bevægelser
(le mouvement social, le mouvement animal, le mouvement
organique et le mouvement materiel), hvis anden
Deel udtrykkelig anbefaler sig til de Nydelsessyge (les
voluptueux), fremsætter han den Grundtanke, at ligesom
den physiske Verdensorden bestaaer ved den materielle
Tiltrækning efter Newtons Lov, saaledes vil den nye
% --- p. 83 ---
\setcounter{footnote}{0}%
\opage{83}Samfundsorden komme til at beroe paa en gjensidig
Tiltrækning af Individernes Tilbøieligheder (attraction
passionnée); enhver Hindring i denne Henseende maa
især være piinlig for Kvinden. Universet er en stadig
skabende Association, i hvilken alle Sole og Planeter
ere Medlemmer og Medarbeidere efter deres forskjellige
Egenskaber og Tilbøieligheder. Paa Jorden er Menneskeslægten
det skabende Element. Det er Samfundets
Opgave at gjenoprette Harmonien mellem Menneskenes
mangfoldige Drifter og Lidenskaber ved at gruppere
Individerne til forskjellige Arbeider efter deres Tilbøieligheder
og Interesser.

Naar man ret har indseet, at de sædelige Bevæggrunde
udelukkende maae søges i Drifter og Tilbøieligheder,
saa vil man kunne forstaae det Hensigtssvarende
i de provisoriske Ægteskaber og trænge ind i Erotikens
Mysterier. I en af sine Romaner lader George Sand
Elskeren sige: „Jeg har aldrig bearbeidet min Indbildningskraft
for paa Ny at tænde eller gjenoplive en
Følelse i min Sjæl, som ikke mere fandtes der, jeg har
aldrig paalagt mig selv Elskov som en Pligt, Bestandighed
som en Rolle. Naar jeg har følt Elskoven slukkes
i min Sjæl, har jeg sagt det, uden at skamme mig
derover og uden Samvittighedsnag.“ Elskerinden i en
anden Roman fra samme Haand, tilstaaer uforbeholdent:
„Af alle disse Forelskelser, som jeg barnligt og blindt
havde givet efter for, syntes ingen og ingen Forbindelse
mig saa dadelværdig som den, jeg til Trods for mig
selv forsøgte at lade vare ud over dens Tid.“

\emph{Fornægtelse} af \emph{Gud} og \emph{Forsyn}. Naar Dualismen
er utaalelig fordi den ikke blot forudsætter en Verdensnatur,
men ogsaa en Gud, saa maa Gud for Monismens
% --- p. 84 ---
\setcounter{footnote}{0}%
\opage{84}Skyld nødvendigviis afskaffes, thi den virkelige Natur
kan man da ikke afskaffe. Men naar Gud falder fra,
staaer Forsynet paa svage Fødder. Den Prophet og
Verdensfrelser, som ved at erklære Gud og Forsynet
aabenbar Krig især har gjort sig navnkundig, er Pierre
Jos. Proudhon (f. 1809, d. 1865). En sand Ildaand, ---
en Elias iblandt de moderne Propheter --- hæver
Proudhon sit Aasyn mod Himlen, tiltaler det høieste
Væsen og siger: „Nu er du dethroniseret og knust.
Dit Navn, der saalænge har været den Vises sidste
Ord, en Indvielse for Dommeren, Fyrstens Styrke, de
Fattiges Haab, de Angrendes Tilflugt, dette umedgjørlige
Navn vil for Fremtiden være viet til Foragt, til Forbandelse
og blive spottet af Menneskene.“ Hvad er vel
Gud? spørger han. „Gud er Daarskab og Feighed, Gud
er Hykleri og Løgn, Gud er Tyranni og Elendighed,
Gud er det Onde.“ Denne pathetiske Expectoration
finder et jævnt og roligt Udtryk i Positivisternes Grundsætning,
at det Absolute slet ikke er, at Alt er relativt.

At der ikke kan være noget Forsyn, derom er
Proudhon paa det Rene med sig selv, ikke paa Grund
af de naturlige Onder, Sygdomme, Alderdomsskrøbeligheder
o. desl., thi de ere jo Følger af Individernes
Endelighed; heller ikke paa Grund af de moralske
Onder, thi de kunne jo tilskrives Menneskenes Misbrug
af den frie Villie; men --- paa Grund af de oekonomiske
Onder. „Havde der været et guddommeligt Forsyn,
saa maatte det jo fra Evighed af have vidst, hvad vi
stakkels Dødelige efter 6000 Aars smertelige Erfaringer
nu først faae at vide, at Frihed, Rigdom, Kundskab
ikke kan fremmes og realiseres i Samfund som vore,
uden at styrte os i en Afgrund af Elendighed. Hvor%
% --- p. 85 ---
\setcounter{footnote}{0}%
\opage{85}for har Forsynet ikke advaret os, og ledet os paa bedre
Veie, men overladt os til vor egen ufuldkomne Logik?
Fordi der intet Forsyn er.“

Den moderlige Natur, der nu maa træde i Forsynets
Sted, opammer sine Børn ved Nødvendighedens
Bryst og gjør Experimenter med deres Opdragelse,
Experimenter, der baade ere fine og kemiske. „Enhver
Nation“ --- siger Taine --- „ligner et stort af Naturen
foretaget Experiment. Ethvert Land er en Smeltedigel,
hvori forskjellige Substanser ere kastede i ulige Forhold
og under særegne Betingelser. Hine Substanser
ere Temperamenterne og Karaktererne. Disse Betingelser
ere Klimaterne og Samfundsklassernes oprindelige
Situation. Blandingen gjærer efter faste Love, umærkeligt
gjennem Aarhundreder, og frembringer i det ene
Tilfælde et fast Stof, i et andet en Sammensætning, som
exploderer. Det er en Fornøielse at iagttage det dumpe
Arbeide, der bringer disse gigantiske Masser til langsomt
og uophørlig at koge. Man lærer de utallige
Kræfter at kjende, der knuse eller sprede eller sammendrage
den Mængde af levende Partikler, der ere
dem underlagte. Man føler det regelmæssige Fremskridt,
der gjennem en Række af Trin, man kan tælle,
og af Omdannelser, man kan forudsee, bringer dem til
den endelige og forud betegnede Tilstand. Man glæder
sig sympathetisk over Naturens Almagt, og man smiler,
naar man seer den evige Kemiker ved en ringe Forandring
i Forhold, i Betingelser eller i Substanser
frembringe Revolutioner, forfærdige Skjæbner, fremkalde
Storhed eller Forfald og forud bestemme hvert
Folk hvilke Værker det skal udføre, og hvilken Elendighed,
det skal døie. Det er et ædelt Skue at see det
% --- p. 86 ---
\setcounter{footnote}{0}%
\opage{86}uendelige Laboratorium udstrakt i Tiden som i Rummet,
hvor saa mange forskjellige Kar, nogle kolde og
fyldte med den golde Aske, andre kogende og røde
af frugtbare Flammer, aabenbare det bølgende Livs
Forskjellighed og de udødelige Loves Eensformighed.
Skjøndt indespærrede i en Krog af Rummet og af Tiden,
skjøndt Døgnvæsener, hvis Liv maaskee imorgen overskjæres
ved Virkningerne af en Explosion eller ved Tilfældigheden
i en Blanding, kunne vi dog opdage flere
af hine Love og danne os et Begreb om Livet i det
sammenhængende Hele.“

\emph{Verdens Forbedring}. Efter at have frigjort
sig fra de gamle Fordomme og kastet den tunge
supranaturalistiske Oppakning fra sig, føler det moderne
Menneske sig saa let, saa let; Brystet udvider sig, Panden
er løftet, Hovedet klart og Musklerne rede: nu skal
der tages fat paa Arbeidet. Men det er et svært Arbeide.
Verden skal forbedres, Samfundet omdannes fra
Grunden af, og den forkvaklede Menneskenatur saa
godt som gjøres heelt om. Hvad er det, der gjør, at
det ikke rigtigt vil gaae med de provisoriske Ægteskaber?
Den forkvaklede Menneskenatur. „Ægteskabet“, siger
George Sand igjennem sin kjære Jacques, „er nu og
altid efter min Mening en af de mest odiøse Institutioner.
Jeg tvivler ikke paa at det vil blive afskaffet,
hvis Menneskeslægten gjør noget Fremskridt henimod
Retfærdighed og Fornuft; et mere menneskeligt og ikke
mindre helligt Baand vil da erstatte dette og vil formaae
at sikkre Børnenes Existens uden for bestandig
at lænkebinde Forældrenes Frihed. Men Mændene ere
altfor raae, og Kvinderne alfor feige til at forlange en % PRINTED AS IS: „alfor“
mere ædel Lov end den Jernlov, som behersker dem.
% --- p. 87 ---
\setcounter{footnote}{0}%
\opage{87}For Væsener uden Samvittighed og uden Dyd passe
tunge Lænker. De Forbedringer, som nogle ædle Aander
drømme om, er det umuligt at virkeliggjøre i dette
Aarhundrede; disse Aander glemme at de ere hundrede
Aar forud for deres Samtidige, og at førend man forandrer
Loven, maa man forandre Mennesket.“ --- Hvad
er det, som gjør, at det ikke vil gaae med Fourier's
Phalanger? Den forkvaklede Menneskenatur. „Menneskene
skulle ordnes efter deres forskjellige Attraaer i
Kredse paa 12—1800, der have deres egne Bygninger,
Phalanstères. Alle Kredse Jorden over skulle staae
i Forbindelse med hverandre, og Forbindelsestraadene
løbe sammen i Constantinopel. Enhver kan vælge hvad
Gjerning, han har Lyst til, og naar han bliver kjed af
den, kan han ombytte den med en anden. Ingen Attraa
skal der sættes Grændse for, ingen Selvfølelse
skal ydmyges, men Enhver have sin Deel af Udbyttet
i Forhold til Kapital, Arbeide og Duelighed.“ Kan
der tilbydes noget Bedre? Den livløse Natur er normal,
den gaaer endnu som fra evig Tid efter Newtons \emph{attraction matérielle}, men saa forkvaklet er den menneskelige
Natur, at den ikke længere lystrer sin egen
Naturlov: den vil jo ikke gaae efter Fourier's attraction
passionée.

Men den stakkels Menneskenatur er lovlig undskyldt.
Hvad kan det hjælpe at den har en Villie,
naar Villien er et Produkt af modsvarende Tilbøieligheder
(Egoisme og Altruisme), ligesom Vandet er et
Produkt af Ilt og Brint? Physik-Philosophernes store
Propheter, August Comte og Stuart Mill, have jo grundig
beviist, at de i „Samfunds-Statistiken“ og „Samfunds-Dynamiken“
hverken vide Besked med Naturnødvendig%
% --- p. 88 ---
\setcounter{footnote}{0}%
\opage{88}heden eller med den menneskelige Frihed, naar det kommer
til Stykket. At „de menneskelige Handlinger ere det
Resultat, som frembringes af Menneskenaturens almindelige
Forhold samt Menneskets egen særlige Karakteer“,
er uimodsigeligt. Men naar vi saa spørge, hvorledes
den særlige Karakteer danner sig, saa faae vi det
Svar, „at Karakteren igjen er en Følge af de naturlige
Omstændigheder, der opdrager den“, og der tilføies,
som Noget, der forstaaer sig af sig selv, at Individets
„egne bevidste Bestræbelser ogsaa høre med til Omstændighederne“.
Saa seer man paa de statistiske Tabeller,
og seer, hvorledes det Eiendommelige i alle
Frembringelser beroer paa Omstændighederne. Den
hollandske Jordbunds Fugtighed frembringer Græsset,
Græsset frembringer Køerne, Køerne frembringe Smørret
og Osten, Smørret og Osten frembringe Hollænderen, ---
paa Grund af Omstændighederne.

Hvordan skal man nu faae en saadan Verden forbedret?
Lægger man Eftertryk paa den menneskelige
Villies selvbevidste Anstrængelse, kun at denne Villie
ikke maa bilde sig ind, at den kan hæve sig over
Naturaarsagerne eller frigjøre sig fra Betingelsernes og
Omstændighedernes Indflydelse, saa forudsætter man
dog noget Særegent, man kalder Villie. Verdensforbedringen
skal foregaae ved Millioner Villiers forenede Anstrængelser.
Villien maa anstrænge sig: Naturen anstrænger
sig ikke. Verdensgaaden ligger da indesluttet
i det ene Spørgsmaal: hvad er Villien? Svarer man:
Villien er en sammensat Størrelse, et Produkt af Elementer
som „Dyd og Last, som Vitriol og Sukker“, saa
maae vi spørge: hvilken af de mange Elementer er det
da, der maa anstrænge sig? Er det kun et enkelt Ele%
% --- p. 89 ---
\setcounter{footnote}{0}%
\opage{89}ment, eller skulle maaskee alle Elementerne anstrænge
sig? Skulle de anstrænge sig før Produktet, med Produktet
eller efter Produktet? --- Man bliver vist
nødt til at opgive Verdensforbedringen og lade det
beroe ved, at glæde sig sympathetisk over „Naturens
Almagt“, og smile, naar man seer „den evige
Kemiker“ ved Analyser og Sammensætninger af Elementerne
„forfærdige Skjæbner, fremkalde Storhed eller
Forfald“ og forud bestemme hvert Folk“, og da ogsaa
hvert Individ i hvert Folk, „hvad Stort, det skal udrette
og hvilken Elendighed det skal døie“.

\emph{Haab} og \emph{Mismod}. At de nye Propheter gaae frem
med et straalende Haab, er forstaaeligt; de ville jo Frihed
og Fremskridt, og Frihed og Fremskridt er Menneskehedens
Løsen. Men at det desuagtet kan mørknes for dem,
og Skyer af Mismod brede sig over deres Aasyn, naar de
ret betænke, hvor forkvaklet Menneskenaturen er, og
hvor langsomt det gaaer med Verdens Forbedring, er
ogsaa forstaaeligt. Det er jo godt nok med al denne
Skæbnesympatie, denne Udsoning med Nødvendigheden,
denne Hvilen ved Naturens Bryst; men Sligt har ogsaa
sin betænkelige Side. „Den evige Kemiker“ med sine
Experimenter er ikke altid at stole paa, Naturens Moderbryst
er snart for varmt og snart for koldt. „Over for den
store og frugtbare Moder, hvem Intet forstyrrer, fordi
Intet er udenfor hende, udvikler sig“ --- siger Taine ---
„i Sjælen en Følelse, en dyb og hidtil ukjendt. Selve
dens Grundvold kommer for en Dag, de utallige Lag,
som Livet har afsat over den, alle Levninger af Lidenskaber
og Forhaabninger, alt det menneskelige Skarn,
der har ophobet sig paa Overfladen, svinder bort; den
bliver eenfold, den gjenfinder de ældste Tiders Instinkt,
% --- p. 90 ---
\setcounter{footnote}{0}%
\opage{90}de ubestemte, eensformige Ord, der fordum satte den i
Forbindelse med Guderne, de naturlige Guder, der leve
i Tingene; den føler, at alle de Ord, som den siden
har sagt eller hørt, kun ere indviklet Snak, kun Øieblikkets
Udbrud, kun Støi paa Gaden, og at hvis der
er et Minut, som er sundt og attraaværdigt i Livet,
saa er det det, i hvilket Sjælen glemmer den store
Myretues Summen og Plagerier, og, som de gamle Vise
sagde, opfatter Sphærernes Harmonie, det evige Universums
Hjertebanken.“

Men Naturen er ikke altid saa festlig stemt, og
Sjælen ikke altid beredt til en saadan Evigheds-Høitid.
Naar Grundkræfterne slaae sig løs og Elementerne
hamre, maa Individet krybe i Skjul, hvis det ikke vil
knuses. Selv i den ophøiede Ro, hvormed den tause
Natur seer ned paa Myretuens Travlhed med Store og
Smaa, er der en Ligegyldighed, som minder om det
overlegne Smiil, hvormed Olympens Guder fordum saae
ned paa de dødelige Slægter. „Hvilken Samling af
Ruiner, hvilken Kirkegaard er ikke Historien! Hvor
mange Menneskehjerter have ikke banket, af hvilke
der ikke er blevet andet Spor tilbage end en Form,
hugget i en Steen! Hvor ligegyldig smiler ikke den
fredelige Himmel, hvor grusom er ikke den lysende
Kuppels Skjønhed, der Dag efter Dag staaer udspændt
over Slægter som falde, ligesom en Baldakin staaer
over en Ligkiste.“

Hos blødere Naturer, hos en Alfred de Musset,
kan en saadan Stemning af Mismod i kritiske, smerterige
Øieblikke endog slaae over i den sorteste Fortvivlelse.
Under Masken \emph{Rolla} har Musset udtalt Fortvivlelsen
over, at det Gamle er faldet, og at det Nye
% --- p. 91 ---
\setcounter{footnote}{0}%
\opage{91}har ingen Grund. „Jeg troer ikke, Kristus, paa dit
hellige Ord; for seent er jeg kommen til en Verden,
der er overlevet. Af et Aarhundrede uden Haab opstaaer
et andet uden Frygt; vort Aarhundredes Kometer
have affolket Himlene. --- Dog, lad det være
tilladt det mest vantro Barn i dette Aarhundrede uden
Tro at kysse Støvet paa dit Kors og at græde, o Kristus,
paa denne kolde Jord, som levede af din Død og
som vil døe uden Dig. For hvem arbeide I, taabelige
Nedbrydere, naar I lemlæste Kristus paa hans \emph{Alter}\dots{}
I vilde danne Mennesket efter Eders Indbilding, I \emph{vilde}
forme en Verden. Velan, I have formet den. Eders
Verden er prægtig og Eders Menneske er fuldkomment,
Bjergene ere jævnede, Sletten er oplyst, I have forstandig
omhugget Livets Træ. Alt er vel feiet paa Eders Jernveie,
Alt er stort, Alt er godt, men man kvæles i Eders
Luft. I lade klinge prægtige Ord, de flagre langt bort
i de forpestede Vinde. --- Hykleriet er dødt, man troer
ei mere paa Præsterne; Dyden døer, man troer ei mere
paa Gud. Voltaire kan da sove tilfreds med det hæslige
Smiil flagrende over hans indtørrede Kinder. Hans
Aarhundrede, siger man, var for ungt til at kunne læse
ham; vort maa behage ham, og hans Mænd ere fødte;
thi den uhyre Bygning er falden over os, den, han med
sine brede Hænder undergravede Nat og Dag. Velan
da, siden det er saaledes: Leve Kjærligheden, som Beruselsen
ledsager, lad de brændende Kys gløde af
Spaniens Viin. Lad Svimmelaanden og de larmende
Fester føre os i Armene paa Glædens Genius. Velan
da, besynger Bakkus, Kjærligheden og Afsindet, drikker
paa Tiden, der svinder, til Døden, til Livet. Glemmer,
% --- p. 92 ---
\setcounter{footnote}{0}%
\opage{92}drikker, lad Friheden leve, besynger Guldet og Natten,
Vinen og Skjønheden!“

\emph{Nihilisme}. Man skjønner heraf, hvad Verden er, hvad
Livet er værd og hvad al menneskelig Digten og Tragten
har at betyde; det er Skuffelse, Blændværk og Illusion altsammen.
I Forfatteren af „Det Ubevidstes Philosophie“,
E. v. Hartmann, har den moderne Verdensanskuelse,
\emph{Monismen}, en Prophet, der er saa meget mere tilforladelig,
som han tillige er monistisk Philosoph. Der
er, ifølge v. Hartmann, tre Bevidsthedstrin, svarende
til tre Stadier paa Illusionernes Vei. --- Illusionens første
Stadium er betegnet derved, at Individet bilder sig
ind at kunne i den givne Verdenstilstand opnaae Lyksalighed
paa Jorden. At det er et Bedrag, er let at
indsee. Al Lyst er kun en Betingelse for Ulyst, al Nydelse
for Lidelse, al Glæde Betingelse for Sorg og
Smerte, og omvendt. Dersom Menneskene blot vilde
gjøre sig det klart, at Glæden forholder sig til den
overstandne Sorg som Smertens Ophør til Pinebænkens
Kvaler, vilde de snart blive kurerede for den om sig
gribende Nydelsessyge. Hvad man nutildags kalder
Pauperismens Spøgelse, er intet Andet end Bevidstheden
om, at Kampen med Nøden, Sorgerne og deres Lindring
ligger paa Smerte-Siden af Fornemmelsens Nulpunkt,
medens tidligere, da Massens Armod var ti Gange større,
denne Bevidsthed manglede, og Folket bar sin Fattigdom
som af Guds Naade. --- Paa Illusionens andet Trin
bilder Individet sig ind at skulle naae Lykken i et hiinsidigt
Liv efter Døden. Det er det kristelige Bedrag.
Jesu Verdensanskuelse er for naiv og barnlig til at adskille
Sjæl og Legeme, hvorfor ogsaa „Kjødets Opstandelse“
er med i den tredie Troesartikel. --- Illusionens
% --- p. 93 ---
\setcounter{footnote}{0}%
\opage{93}tredie og sidste Stadium er det, hvorpaa vi kjøre som
paa Jernbane. Indbildingen er, at Lykken ligger i en
stor politisk social Fremtid. Man venter sig Frelsen
af en forbedret Samfundsorden, af Fremskridt i Industri
og Videnskab: et uhyre Bedrag!

„Enden bliver, at den gamle, udlevede Menneskeslægt
er træt, træt til Døden, træt af Arbeiderne og
af Nydelserne, uden Arvinger og uden Kjærlighed, træt
af Alt. Da vil den komme, hiin ophøiede Melancholie,
som man sædvanlig finder hos Genier og aandelig høitstaaende
Oldinge; da vil den komme, den rene Himmelklarhed,
hiin kjølig guddommelige Ro, som vifter i
Spinozas Ethik, hvor Lidenskaberne sluges af Fornuftens
Afgrund, fordi de klart og tydelig ere fattede i
Ideer. Da er Illusionen død; da er Haabet udbrændt,
thi hvorpaa skulde man endnu haabe? Da erkjender
den dødtrætte Menneskehed kun Eet, der er saligt, ---
Hvile fra Sorg og Savn, Smerteløshed i det rene Intet.“

Men tør vi nu paastaae, at de nye Propheter alle
uden Undtagelse ere Nihilister? Det kommer an paa,
hvad man forstaaer ved Nihilisme. Efter den sædvanlige
Mening er Intet det rene Nul, det Tyndeste og
Tommeste og Kjedsommeligste af Alt: Noget er altid
bedre end slet Intet. Fra dette Synspunkt skulde Nihilisterne
altsaa være Mennesker, som Intet vilde, ikke
meente det alvorlig med Noget og intet Maal havde,
for hvis Opnaaelse de fandt det værdt at arbeide. Det
er en stor Misforstaaelse. Nihilisterne ere Folk med
Sands for det virkelige Livs Opgaver, ja med en næsten
graadig Appetit paa Virkeligheden. Det Intet, der
afgiver Slagordet, betyder kun, at Gud er Intet, at det
Overnaturlige, det Hiinsidige er Intet, men ikke, at Na%
% --- p. 94 ---
\setcounter{footnote}{0}%
\opage{94}turen med dens Kræfter og Love, ikke, at det nærværende
Liv med dets Arbeider, Midler og Maal er Intet.
Er Livet da virkelig bedre end Intet? Det kommer
an paa Omstændighederne. Saalænge Livskræfterne tilbyde
sig, bør man bruge dem. Bør? Forstaaer sig,
det er ingen ubetinget Pligt; Alt er relativt, det Ubetingede,
det Absolute er Intet. Enhver maa bruge
Kræfterne i Overensstemmelse med sin individuelle
Natur, sine Drifter og Tilbøieligheder, og være beredt
paa Følgerne. Derfor finde vi Ørkesløse og Arbeidsomme,
Letsindige og Alvorlige, Feige og Modige, Blødagtige
og Haardføre, Nydelsessyge og Handlekraftige,
Egoister, som kun sørge for sig selv, og Altruister, som
offre sig og virke for den hele Menneskehed, i brogede
Blandinger samlede under Nihilismens Flag. De indrømme
alle, at Livet kan være bedre end Intet, men
det tilstaae ogsaa alle, at det kan være værre end
Intet: kun de philosophiske Hoveder --- men de ere
jo næsten alle philosophiske Hoveder --- vide med sig
selv, at det Hele i Grunden er Intet.

Dette „i Grunden Intet“ er Hemmeligheden, den
Hemmelighed, der gjør det af med Dualismen. Dualismens
Skavank er en Tohed, og Toheden splitter sig
igjen til Mangfoldighed. Den virkelige Natur er befængt
med Dualisme: Opad staaer imod Nedad, Høire
imod Venstre, Magnetens positive Pol mod den negative,
Syren imod Basen, Hanblomsten imod Hunblomsten,
end ikke det flygtigste provisoriske Ægteskab kan fuldbyrdes
uden Dualisme. Monismen vil Eenheden; men
den forskjelsløse Eenhed af Alt er Intet. Saalænge
Noget adskiller sig fra Andet, saalænge det ene Blad i
Skoven er uligt det andet, den ene Bølge i Havet mod%
% --- p. 95 ---
\setcounter{footnote}{0}%
\opage{95}sat den anden, saalænge har man kun Fleerheden ikke
Eenheden. For at naae den fuldstændige Eenhed maae
vi tænke os al Ulighed, al Forskjellighed udslettet.
Naar al Forskjellighed falder bort, staaer den forskjelsløse
Eenhed tilbage som Grunden; men den forskjelsløse
Eenhed er i Grunden Intet: den gjennemførte Monisme
er Nihilisme. Mener man, at den philosophiske Nihilisme
er upraktisk, men at Monismen dog er den eneste
Verdensanskuelse, som praktisk forstaaet, passer for Nutiden,
den eneste, der er i Stand til at vække Ungdommens
Mod og pege paa Maalet, saa spørge vi: hvad er
det da for en Befrier, der staaer der --- henne ved Maalet?
Døden, den store Nihilist! Døden er praktisk, og
Døden er Intet. Døden udjævner Alt, den mætter Ethiken,
den fordøier Kritiken, den ophæver Modsigelsens
Grundsætning. I Nihilismen, grundig forstaaet, er der
ingen Forskjel mellem Godt og Ondt, mellem Sandhed
og Løgn. De gamle Livets Propheter ere døde, de nye
Dødens Propheter leve --- for at døe; saa mødes tilsidst
de Nye med de Gamle i Dødens rene Intet, den uendelige
Al-Eenhed: det er monistisk Forsoning.

% --- p. 96 ---
\setcounter{footnote}{0}%
\lecture{Sjette Forelæsning.}\opage{96}Der er trods den stærke Gjæring i Tiden, dog endnu
udentvivl et Fleertal af Medborgere, som, ihvorvel de
have Adskilligt at indvende mod det Bestaaende, med
Uvillie vende sig bort fra de nye Propheter; den
besindige Dom lyder saaledes: Naar vi høre disse
Menneskers ofte vilde og forargelige Udbrud af Gudsbespottelse,
naar vi see den overmodige Selvtillid,
hvormed de optræde, den Hensynsløshed, hvormed de
forvirre Begreberne om Pligt og Samvittighed og sønderrive
sædelige Baand, den Voldsomhed, hvormed de angribe
Samfundets Grundpiller, uden dog ved deres løse,
hverandre modsigende Projekter at sætte noget Bedre
i Stedet, saa maae vi, med al Agtelse for Enkeltes
Talent og velmeente Iver, holde os overbeviste om, at
de ere paa en Afvei. Ja, der er vel dem, hvis Forkastelsesdom
vil give sig et haardere Udtryk. Og dog
tør vi, hvad Sagen i det Hele angaaer, ingenlunde mene,
at Almeenheden er paa det Rene med sig selv. Der
røres ved ømme Sider, der tales om Noget, som vækker
Opmærksomhed; de nye Propheter have Ord, der klinge
i Nutidens Ører.

Det er ganske sandt, at Ordførere, der forstaae at
ægge Lidenskaberne og blænde Mængden ved lovende
% --- p. 97 ---
\setcounter{footnote}{0}%
\opage{97}Udsigter, maae have et stort Forspring for dem, der
føle sig forpligtede til at indskærpe Selvbegrændsning
og sindig Fremgang. De, der sige til den menige
Mand: Veien er banet, Tiden er kommen, slaae til!
faae vistnok flere Tilhængere end de, som tidlig og
sildig maae minde om at gjøre smaa Skridt, eftersom
Hindringerne ere mangfoldige og de hullede Veie lidet
fremkommelige. Hvorledes skulle de Førere, der kræve
Opdragelse, Tugt og Orden, kunne holde Skridt med
dem, hvis Feltraab er: Sammenhold, Dristighed, hurtig
Afgjørelse! Derfor maae vi ingenlunde bilde os ind, at
Stormen efterhaanden vil sagtnes og Gjæringen sætte
sig, naar man kun i Stilhed bliver ved med det Gamle
og forsigtig undgaaer at røre ved de egentlige Grundspørgsmaal.
„Bare holde igjen, holde igjen!“ Dette
„Holde igjen“ virker narkotisk paa Samfundets Nerver;
det Bestaaende tier sine Livsspørgsmaal ihjel, som en
dvask Moder ligger sine Børn ihjel. Naar man saa en
Gang som ved et Pistolskud vækkes op af sin aandelige
Døs, tager man paa, som gjaldt det at skyde
Graaspurve med Kanoner. Orthodoxien farer i Harnisk,
der ringes med den store Klokke: „Antikristen
er kommen, han gjør Tegn og forfører Folkene; Gog
og Magog staae for Døren. Frem med Sakramenterne,
bort med Fornuften; kun Præsten har Myndighed, kun
Præsten tør see den levende Satan stivt i det lumske
Øie!“

Vi maae for alting ikke glemme, at de geniale Kræfter
ogsaa stamme fra Aand, og at Aanden ikke lader sig
binde i udvortes Lænker; jo mere Lænkerne tynge, desto
virksommere ere disse Aander, desto ivrigere arbeide
de. Vi maae ikke glemme, at der paa Aandsdannelsens
% --- p. 98 ---
\setcounter{footnote}{0}%
\opage{98}Vei henstaae Spørgsmaal i Mængde, afgjørende Spørgsmaal,
som den frie Tankes Propheter have reist, men
som det Bestaaendes Forsvarere saa gjerne gaae uden
om? Ja, der er Mange, der mene, at det er af gode
Grunde. Vi maae ikke glemme, at der i Naturalismens
Leire findes Mænd, som have skarpt Øie for sjælelige
Tilstande, og tale den grundige Undersøgelses klare,
indtrængende Sprog. Naar en saadan Forsker samler
sig til et paamindende Ord, kan han have Adskilligt at
lægge sin Samtid paa Sinde; vi lade ham tale.

„Der er, som bekjendt, en betydelig Forskjel
imellem juridisk og psychologisk Sandhed, imellem
juridisk og psychologisk Ret. Naar En forsikkrer, at
han troer hvert Ord i den lutherske Kirkes symbolske
Bøger, saa har han, hvis der af hans andre Udtalelser
ellers ikke kan anføres Noget, som strider herimod,
unegtelig juridisk Ret til at kalde sig Lutheraner. Men
om han ogsaa har psychologisk Ret, er et andet
Spørgsmaal. Sæt at han, følgende Kulturens Strøm,
efterhaanden er kommet til at leve med heelt andre
Forestillinger om Himmel og Jord, om Naturligt og
Overnaturligt, om Guds og Menneskets Væsen, ja med
Forestillinger, som de symbolske Bøger udtrykkelig
fordømme, saa befinder han sig, om han end ikke er
blevet sig Afvigelserne klart bevidst, i en Tilstand af
psychologisk Usandhed; og hvis han med fuld Bevidsthed
om Afvigelserne, beraaber sig paa sin juridiske
Ret, saa sige vi ham op i Øinene, at han lever i en
psychologisk Løgn.“

„I Nutidens Lutheranere med Eders positive Aabenbaring,
Eders Bibel og symbolske Bøger, leve I virkelig
med Luthers Tro, i Luthers Aand efter Luthers Anviis%
% --- p. 99 ---
\setcounter{footnote}{0}%
\opage{99}ning? Dersom I have Luthers Tro paa Kristus, saa
maae I jo ogsaa have Luthers Tro paa Djævelen, thi
den er fuldt ud hjemlet i Skriften. At de gamle
Lutheranere for ramme Alvor troede paa det Overnaturlige,
beviste de ved deres levende Tro paa Djævlebesættelser.
Troe I endnu paa Djævlebesættelser, da
maae I ikke blot finde det naturligt, at de moderne
Propheter lee ad Eder, men, hvad mere er, I maae
bryde med hele den nyere Kultur, den nærværende
Virkelighed, hvori I dog leve og aande. I stamme og
famle og søge Udflugter, naar man gaaer Eder paa
Klingen med Djævelen, og mærke ikke, at I med det
Samme stamme og famle i Eders Tro paa Kristus.
Kristi Personlighed faaer jo netop sin Betydning ved
dens udtrykkelige Modsætning til den personlige Djævel,
dertil er, efter Skriftens Ord, Guds Søn aabenbaret,
at han skal forstyrre Djævelens Gjerning. I tale med
Salvelse om Synd og Naade, men om I ogsaa med
Luther og Skriften selv vide, hvad Naaden virker,
hvorledes den virker, og hvad Synd er, derom er det
tilladt at tvivle; vist er det, at mange af de Friheder,
I nu ansee for tilladelige, have de gamle Lutheranere
stemplet som Dødssynder. I staae med den ene Fod i
Nutiden, med den anden i Fortiden; I have to Rum i
Bevidstheden, I gaae fra Nutid til Fortid saa let som
man i et Huus gaaer fra det ene Værelse ind i det
andet, naar der kun er en Dør imellem, og Døren
staaer paa Klem“.

„I lade haant om Monismen, fordi den bliver
staaende ved vort Livs Naturgrund uden at indblande
et Overnaturligt; I opbygges ved Dualismen, at Kristus
strider imod Satan, saa at det nye Menneske, især om
% --- p. 100 ---
\setcounter{footnote}{0}%
\opage{100}Søndagen, er i Færd med at kvæle det gamle. Men
I ville dog ikke have den gamle Adam reent kvalt, thi I ere
humane. Eders Fædre vare haardhjertede, de pidskede
den Gamle og kunde uden Skaansel see den lede
Synder hængt. Men I ere ikke haardhjertede, I have
Medlidenhed med den gamle Adam, han er jo dog saa
ung endnu! Vel skal han straffes den Skjelm; men
naar det kommer til Execution, naar han spræller i
Snoren og gisper efter Luft, saa finde I Straffen barbarisk.
Og den gamle Eva --- endnu saa ung og skjøn!
hvor megen Uret, ja hvor mange Krænkelser har hun ikke
at hævne. Men naar hun seer sin Adams angrende Mine,
hans bedende Blik, og hører hans Suk, saa glemmer
hun Alt, saa tilgiver hun Alt. Hun bøier sig over den
Ydmyge, hun klapper hans Kind, hun kysser hans
Mund, hun vifter ham med Poesiens Vifte, hun svaler
hans Pande af Erotikens Eau de Cologne Flaske. See
saadan er Eders lutherske Dualisme!“

„Det maa dog være skrækkeligt at leve med
et spaltet Væsen, en Sjæl, hvis ene, syndige Deel
med Rette tilhører Djævelen, hvis anden og gjenfødte
Deel tilhører Kristus. I gaae med en tvedeelt
Sjæl og troe paa et evigt Liv. Djævelen med den
gamle Adam trækker til den ene Side, Kristus med
den Hellig-Aand til den anden. Hvis Kristus seirer,
saa blive I evig salige, men --- I glemme det ikke! ---
hvis Djævelen seirer, saa blive I evig fordømte. Hvor
kunne I leve paa saadanne Vilkaar? Jeg veed jo nok,
at der er dem iblandt Eder, som bilde sig ind, at her
slet ingen Fare er, da I jo have Daaben og Troen.
Men den Tro, der skal bringe Eenhed i en saadan
% --- p. 101 ---
\setcounter{footnote}{0}%
\opage{101}Adsplittelse, Fred og Forlig trods en saadan Spænding,
maa have mere end herkuliske Kræfter; det maa være
en mirakelgjørende Tro, en Tro, der kan flytte Bjerge.
Hvilke Mirakler have I da gjort, hvormange Bjerge har
Eders Tro flyttet? De virkelige Bjerge lade I kjønt
staae; saa maa det vel være Phantasiebjerge. I have
Lutherdom i Følelser, Rørelser og trøsterige Phantasier.
I stride tappert med den gamle Adam og vinde den
ene Phantasieseier mere glimrende end den anden.
Kommer Djævelen engang imellem bag paa Eder med
en Smule Politik, hvad gjør saa det? I have jo Klippen
i Troens Ord og i Daaben et Naadens Kinderpulver.“

„I have maaskee læst en Fortælling af Hoffmann,
der hedder „Elixire des Teufels“, eller, hvis I ikke.
have læst den, saa læser den endelig, at I ret maae
faae Indtryk af, hvad det vil sige, at faae sit Jeg, sit
Væsen, dualistisk spaltet i Phantasie. --- En ung, begavet
Munk ved Navn Medardus har i et Kloster fundet en
Flaske med Djævle-Elixir; han drikker af Flasken, og
strax bliver hans Væsen spaltet, saa at han kan vise
sig i mere end een Skikkelse. Indviklet i Elskovseventyr
og skrækkelige Forbrydelser, bliver det ham
reent umuligt at finde sig selv igjen: „Jeg er, siger
han, det, som jeg synes, og synes ikke det, som jeg er;
mig selv en uudgrundelig Gaade, ligger jeg i Splid
med mit eget Jeg.“ Efter at have været i Fængsel og være
blevet forhørt og frikjendt for et Mord, som han dog
virkelig har begaaet, løslades han, og nu kommer hans
Livs lykkeligste Dag, det er hans Bryllupsdag. Han
fortæller: „Der opstod en dump Larm paa Gaden, hule
Stemmer raabte imellem hverandre, og den drønende
% --- p. 102 ---
\setcounter{footnote}{0}%
\opage{102}Raslen af en tung, langsomt rullende Vogn lod sig
høre. Jeg ilede til Vinduet. Der stod netop foran
Palladset den af Rakkerknægtene trukne Vogn, paa
hvilken Munken sad baglænds; foran ham en Capuciner,
der høit og ivrigt bad sammen med ham. Han var
vanheldet af Dødsangstens Bleghed og det strittende
Skjæg; dog var den skrækkelige Dobbeltgjængers Træk
mig kun alt for kjendelige. Ret som Vognen, der et
Øjeblik var blevet standset af den sig sammentrængende
Folkemasse, atter rullede videre, kastede han det stive
forfærdelige Blik fra de funklende Øine op paa mig og
lo og hylede op: Brudgom, Brudgom --- kom op paa
Taget --- saa ville vi brydes med hinanden, og den,
som støder den anden ned, er Konge og tør drikke
Blod. Jeg raabte: Forfærdelige Menneske, hvad vil Du
mig? --- Aurelia omfavnede mig, rev mig fra Vinduet
og raabte: For Guds og den hellige Jomfrues Skyld\dots{}
De føre Medardus, min Broders Morder, til Døden\dots{}
Leonard --- Leonard! Da bleve Helvedaanderne i mit
Indre vaagne og reiste sig iveiret med den Magt, der
var dem skjænket over mig brødefulde, forrykte Synder.
--- Jeg greb Aurelia med Raseri, saa hun foer sammen:
Ha, ha, ha --- afsindige Kvinde --- jeg --- jeg, din
Elsker, din Brudgom er Medardus, er din Broders
Morder. Du, Munkens Brud, vil hyle Fordærvelse ned
over din Brudgom? Ho, ho, ho --- jeg er Konge ---
jeg drikker dit Blod.“ --- Han farer ned paa Gaden, befrier
Munken, hugger og stikker og slaaer om sig og
flygter ind i Skoven. „Kun den Tanke at flye som et
hidset Dyr stod fast i min Sjæl. Jeg reiste mig, men
neppe var jeg gaaet nogle Skridt, før et Menneske
% --- p. 103 ---
\setcounter{footnote}{0}%
\opage{103}sprang op paa min Ryg og slyngede sine Arme om
min Hals. Forgjæves søgte jeg at ryste ham af, jeg
kastede mig ned, jeg trykkede mig ind imod Træerne.
Alt forgjæves; Mennesket fniste og lo haanligt. Da
brød Maanen klart lysende gjennem de sorte Grantræer,
og Munkens dødblege, græsselige Ansigt, den
formentlige Medardus's, Dobbeltgjængerens Ansigt
stirrede paa mig med et græsseligt Blik som op fra
Vognen. Hihihi --- lille Broer --- Du maa bære mig,
nu er jeg altid hos Dig, jeg kan ikke løbe, jeg kommer
lige fra Galgen, de vilde radbrække mig.“

„Lutheranere, forstaae I denne Historie? --- Luther
forstod den, da han sad ene i Natten og kastede Blækhuset
mod Væggen paa Wartburg. Det er en lille
Historie om Dualismen.“

Efter at have hørt, hvor bittert og skarpt den
moderne Monismes Ordfører udtaler sig mod Dualismen,
ere vi komne til et Vendepunkt, hvor det egentlige
Stridsspørgsmaal: Aandsvillie eller Driftvillie? kan
træde tydelig frem, og det store Mellemværende vise
sig i fuld Belysning. Dualismen er ingen vilkaarlig
Opfindelse, et Indfald, man kan slaae ihjel med et
Pennestrøg. Dualismen har Livets og Realitetens
Vidnesbyrd for sig: Gamle og nye Propheter, hvilken
Dualisme! Saa længe al Aand endnu ikke er blevet til
Natur, saa længe den elementære Natur endnu ikke
er forflygtiget til Aand, saa længe der endnu er en
Rest af Modsætning mellem Aandsvillie og Driftvillie
tilbage, saa længe staaer Dualismen ved Magt. Dyret
med sin Driftvillie mærker ikke Dualismen; men
Mennesket med sin Aandsvillie faaer den tilgavns at
% --- p. 104 ---
\setcounter{footnote}{0}%
\opage{104}føle. Mennesket er dualistisk, Aben er monistisk;
skulde det radikale Middel til Dualismens Afskaffelse
maaskee bestaae i at reducere Mennesket til Abe?

Det eneste alvorlige Forsøg, den nye Tid har seet,
paa at udjævne det Dualistiske, er aabenbart Rationalismen,
ikke i sin ældre, nøgterne Forstandsform, men
i en alsidigere, aandrigere, med Ideer fortrolig Form
som \emph{Nyrationalisme}. Denne Bestræbelse har fundet
dygtige Førere, begavede Talsmænd i Tydskland, i
Schweiz, i Holland, i Sverrige. Hvad vil da Nyrationalismen?
„Den vil med Schleiermacher søge Religionens
inderste Væsen i \emph{Sindets Dyb}, i Menneskets Livsgrund,
i Aandens egentlige Midtpunkt, men ogsaa sætte
Erkjendelsens og Villiens Udfoldning i inderlig Forbindelse
med dette centrale Forhold. Den vil fra et ethisk
Standpunkt tilveiebringe en Rensning og Fornyelse af
en stor Deel Dogmer, henhørende til den anthropologisk-soteriologiske
Kreds, af Læren om den frie Villie, om
Synden, om Naaden og Istedtrædelsen o. s. v. Den
vil søge at opnaae en dybere Forbindelse mellem guddommelig
Afhængighed og menneskelig Frihed, at
komme til en sandere Forstaaelse af Vexelvirkningen
mellem den guddommelige og menneskelige Faktor i
Frelsens Foregang. Den vil ikke blot rense og inderliggjøre
Dogmatiken, omdanne dens udvortes juridiske og
magiske Forestillinger til sædelig religiøse; den vil
ogsaa uddybe Moralen, skille den ved dens forfængelige
Selvretfærdighed, dens overfladiske Pelagianisme, kort
sagt, ved dens Endelighed, idet den nedsænker Sædelighedens
Rødder i Religionens Dyb, Menneskets endelige
Villie i den guddommelige Villies Uendelighed, hans
Frihed i Bundenhed af det Guddommelige. Den vil
% --- p. 105 ---
\setcounter{footnote}{0}%
\opage{105}derved ophæve den skjæbnesvangre og fordærvelige
Adskillelse af det Religiøse og det Sædelige.“\footnote{Den nyeste Theologies Historie i kort Fremstilling af Dr. Karl Schwartz. Dansk Oversættelse. Kjøbenhavn 1869. S. 466---67.}

Det er i Sandhed en ædel Bestræbelse; men Spørgsmaalet
er, om den lader sig gjennemføre. At Rationalismen
ikke tager Hensyn til den gamle Kirkes Tarv, er indlysende.
Den vil hævde den videnskabelige Forsknings
Uafhængighed af Skriftens og Symbolernes Autoritet,
og heri maae vi give den Ret; men naar den saa
fordrer, at Aabenbaringstroen skal underkaste sig
Videnskabens Autoritet, saa undergraver den jo baade
Kirken og Menighedslivet. Den indrømmer ikke nogen
principiel Forskjel mellem Tro og Viden. Er det f. Ex.
afgjort, at de historiske Beretninger om Kristi Opstandelse
ingen Beviiskraft have i Videnskaben, saa tør
de heller ingen Gyldighed have i Kirken. Kan det
lutherske Nadverdogme, videnskabelig seet, ikke bestaae
for Kritikens Dom, strax maa Kirken opgive sit Dogme
og gjøre Sakramentets Brug afhængig af Videnskabens
Kjendelse. Mon en saadan Fordring er begrundet?
Ja, i en radikal Misforstaaelse af Troens Væsen.

Hvad er det vel, den religiøse Bevidsthed søger i
de historiske Beretninger om Kristi Opstandelse? Et
Vidnesbyrd om, at Apostlene virkelig have troet paa
denne Opstandelse, at Kirken er stiftet, og Menigheden
døbt paa denne Tro. Lad Kritiken bevise, at Beretningerne,
som vi have dem, ikke umiddelbart hidrøre
fra Øienvidner, at Uovereensstemmelser i det Enkelte
tyde paa, at de ere gaaede fra Mund til Mund, før de
% --- p. 106 ---
\setcounter{footnote}{0}%
\opage{106}bleve nedskrevne: siger Intet; Kirken kan indrømme
det, Menigheden blive vidende derom, uden i mindste
Maade at afficeres deraf. Vidnesbyrdet om Kristi Opstandelse
er et Vidnesbyrd af Troende om Troen til
Troende; Vidnesbyrdets Kreds er saa sluttet, at ingen
Kritik udenfra formaaer at gjennembryde den.

Kun een Indsigelse vilde have Gyldighed, men den
hentes ikke fra Bibelkritiken, den hentes fra Physiken,
eller rettere fra Physik-Philosophien, og lyder saaledes:
„Kristus er ikke opstanden, thi Opstandelsen er et
Mirakel, og Miraklet er en Urimelighed, en Modsigelse,
en Umulighed“. Vil nu Rationalismen bøie sig for
denne Videns Autoritet, for at tilfredsstille den moderne
Bevidsthed, saa faaer den at tage Conseqvenserne.

Hvad Videnskaben skal anerkjende, maa kunne
bevises i Videnskaben. Rationalismen taler om den
levende Gud, en guddommelig Kjærlighed, et levende
Forsyn, en hellig, guddommelig, retfærdig Villie, og
vil for Religionens og Sædelighedens Skyld ikke opgive
Sjælens Udødelighed; men, spørge vi, lader Sligt
sig begrunde videnskabeligt? Den moderne Viden kan
lige saa lidt finde sig til Rette med Sjælens Udødelighed
som med Kristi Opstandelse. Skal Guds Kjærlighed
filtreres i Videnskaben, saa bliver det i Sandhed en
kold Kjærlighed.

Rationalismen vil frigjøre den religiøse Bevidsthed
fra umiddelbar Afhængighed af de i Kirken overleverede
Læreformer. Den gængse mekaniske Inspirationslære
benævnes som „et til det Utaalelige plumpt Hylster,
hvori Kjernen er: de bibelske Bøgers enestaaende Værd“.
Det kirkelige Inspirationsdogme, hedder det, har „længe
været erkjendt som Noget, der satte i Forlegenhed,
% --- p. 107 ---
\setcounter{footnote}{0}%
\opage{107}som et hæmmende Onde. Det er ikke Andet end en
Overdrivelse og utilbørlig Forstærkelse af Sandheden,
tjener ikke nogen virkelig, men kun en indbildt Trang
hos Fromheden, nemlig den: at have et absolut færdigt
Udtryk for Sandheden.“

Vi indrømme det. Personlig Frigjørelse fra udvortes
Autoritet er af lige saa gjennemgribende Betydning
for Religionen som for Videnskaben. Det er
Rationalismens udødelige Fortjeneste at have kæmpet
for denne Frigjørelse i Kirken; men desto beklageligere
er det, at den har misforstaaet Frigjørelsens Væsen.
At ville, paa Skriftens og Kirkens Autoritet, tvinge den
religiøse Bevidsthed i Nutiden til en bogstavelig Antagelse
af de bibelske Forestillinger om det overnaturlige
Rige, og endnu mere, til en bogstavelig Fastholden
ved den orthodoxe Dogmatiks indtil de mindste Enkeltheder
udformede Supranaturalmekanisme, er ganske
vist et saa uhyre Misgreb, at man vanskeligt skal kunne
sige, om det er Mændene af den frie Tanke eller
Mændene af den bundne, der bidrage mest til at undergrave
Religionen. Men fordi Israels Folk engang har
levet i en ganske anden Forestillingskreds end vi,
følger saa deraf, at det har havt sit Aandsliv i lutter
Indbildninger? Fordi vore orthodoxe Fædre have dannet
sig Begreber om det Overnaturliges Indgriben i det
Naturlige, hvis objective Gyldighed vi ingenlunde kunne
anerkjende, følger saa deraf, at det Overnaturlige slet
ikke er, at Naturen selv ikke kan have sin Grund i et
Overnaturligt? Hvis Rationalisternes „levende Gud“
skal betyde det Samme, som den levende Natur, saa
indsees ikke, hvad Forskjel der er mellem Nyrationalisme
og Naturalisme. Men hvis den levende Gud som
% --- p. 108 ---
\setcounter{footnote}{0}%
\opage{108}Aand skal være forskjellig fra Naturen, saa maa han
jo være enten et overnaturligt eller et unaturligt Væsen.

Nyrationalismen har et væsentligt Fortrin deri, at
den med Sigte paa Foreningen af det Religiøse og
Sædelige hævder det Humane og vælger sit Midtpunkt
i Menneskets Samvittighed. Ud fra Samvittigheden, det
virksomt nærværende Indre, stræber den „at drage
Religionen fra dens abstracte, overjordiske, oversædelige
Høide midt ind i de sædelige Interesser og forsone
den med Nutidens hele Kulturliv.“ Mon det vil lykkes?
„Menighederne“, siger Karl Schwarz, „ere vaagnede af
deres Slummer, de lutherske Præsters Unatur, Utaalsomhed
og Herskesyge er stegen i saadan Grad, at den
har fremkaldt den stærkeste Modvirkning fra Folket.
Den dybere religiøse Trang hos de Dannede fordrer
bydende og fuldt bevidst en anden Tilfredsstillelse end
den tidligere tilbudte. Vi befinde os endnu i Forventningens
og Forhaabningens Øieblik. Men den nye Tid
nærmer sig, vi føle dens Aandepust vifte. Det, som
Alt peger hen imod, og som maa tilstræbes af alle
klare Aander, er Kirkens Befrielse fra det hidtilværende
Statsvæsens faste Omfavnelse, for at den ikke skal
blive dragen med i alle dets uheldsvangre Kriser.
Denne Befrielse er enstydig med en Opbyggelse fra
Folkebevidsthedens Grund, en fri Udformning af den
ægte, enfoldige, inderlige Kristendom, som den lever i
Folket. Kun i en af Menighedernes Midte fremgaaende
organisk Sammenslutning ville alle de nedtrykte og
skjulte Aandskræfter frigjøres, Nutidens hele Dannelse
føres tilbage i de udtørrede kirkelige Kanaler, og
Kirken selv blive \emph{Sandhedens} og \emph{Livets} \emph{Hjem}%
% --- p. 109 ---
\setcounter{footnote}{0}%
\opage{109}\emph{sted}, Forkynderske af den inderligste Samvittighedstro.“
--- Mon denne Spaadom vil gaae i Opfyldelse?

Skal et Menneske i sin Samvittighed finde Vidnesbyrd
om en evig Retfærdighed, en hellig, guddommelig
Kjærlighed, da maa han have en Troessamvittighed;
thi i en Videnssamvittighed findes Intet derom. Den,
der samvittighedsfuldt vil fjerne Alt, hvad der ikke
har Bekræftelse i Videnskaben, og frigjøre sin Sjæl fra
Forventninger, hvis Opfyldelse ikke lader sig bringe
til fuld Samklang med jordiske Betingelser og uforanderlige
Naturlove, vil i sin Samvittighed ikke finde
noget Spor enten til overjordisk Retfærdighed eller
guddommelig Kjærlighed. Skulde han en Gang i beaandet
Stemning lade sig henrive til at gribe efter det
Overjordiske, kan han ordentlig undsee sig for sig selv
bag efter og føle Samvittighedsbebreidelse; --- „det
var jo kun en Stemning, Hjertet har sin Ret, Livet sin
Poesi“. Det er forgjæves at minde ham om, at det
personlige Værd kun er Indbildning, hvis der i Villien
ikke er noget Evigt og Ubetinget: han slaaer det hen
som en subjectiv Idee og søger Livsgaaden løst i
Physiologien. Uvidenskabelige, halvvidenskabelige
Mennesker ansee hans Vantro for Trods. De forstaae
ham ikke; det er ikke Trods, det er Samvittighed,
videnskabelig Samvittighed. Det har maaske kostet
ham vaagne Nætter og bittre Taarer, inden han gav
Slip paa de trøsterige Illusioner; men --- „Sandheden
er objectiv, Videnskaben er objectiv!“ Mon en saadan
Kandidat i moderne Viden skulde føle Trang til at
lade sin Dannelse gaae tilbage „i de udtørrede kirkelige
Kanaler“ og blive Medlem af en nyrationalistisk
Menighed?

% --- p. 110 ---
\setcounter{footnote}{0}%
\opage{110}Med Troessamvittigheden er det en anden Sag;
den har en Evighedsgnist, et Ubetinget i Selvet. Men
netop herved kommer Modsætningen mellem Overnaturligt
og Naturligt frem i Bevidsthedens Allerinderste.
Idet Personligheden skal til at danne sig, træder Aandsvillien
op imod Driftvillien, og Mennesket finder sig
selv som et dualistisk Væsen. Det Smertelige ved
Dualismen bestaaer ikke deri, at der i Livsgrunden
selv er forskjelligartede Anlæg. Individets Naturliv afgiver
netop Grundlaget for dets Aandsliv. Nei, Spliden
opstaaer af, at Driftvillien i sin sandselige Umiddelbarhed
ikke vil underkaste sig Aandsvillien. Hvad vil
da Driftvillien? En øieblikkelig Tilfredsstillelse af
den naturlige Lyst eller, hvis den øieblikkelige Lyst
skal offres, dog en Tilfredsstillelse af en selvisk herskende
Lidenskab. Og hvad vil Aandsvillien? Herske,
og herske ubetinget. Det er en sand Plage for det
naturlige Menneske at faae Aandsvillien med, thi den
vil aldrig tillade Drifterne, end ikke de geniale, at
slaae sig løs og gaae paa egen Haand; den passer paa
dem som en Formynder, den sørger for deres Opdragelse
og holder dem under streng Tugt. Den
umiddelbare Følelse kan ordentlig finde det oprørende,
at de uskyldige Tilbøieligheder, Nydelsernes glade
Genier, disse Naturens legende Børn saa smaalig
holdes i Ave; men Samvittigheden, Troessamvittigheden,
giver Aandsvillien Ret; begge ere enige i, at der ubetinget
maa holdes fast paa det Ubetingede. Det er
ikke af Herskesyge heller ikke af Had til Drifterne eller
af Foragt for Individets naturlige Anlæg, at Aandsvillien
saaledes vil herske; nei, det er for Frihedens
Skyld, for Personlighedens Skyld, at der kan komme
% --- p. 111 ---
\setcounter{footnote}{0}%
\opage{111}Selvstændighed i Individets Karakteer, Eenhed i Menneskets
Væsen. Den sande Eenhed er Selveenhed: Udviklingens
Maal er det personlige Selv.

Ved at hæve sig over det blot Naturlige og forvandle
Drifterne til Midler viser Aandsvillien sit Slægtskab
med det Overnaturlige. Livsopgaven kan ikke
løses asketisk ved at underkue det Naturlige; men den
kan heller ikke løses monistisk ved at fornegte Aandsvillien,
for at kultivere Driftvillien. Er der slet intet
Overnaturligt i Menneskets Villie, saa maae vi lade
Naturen raade sig selv, og saa seer det betænkeligt
ud med det nyrationalistiske Sædelighedsprincip. Hvad
er vel i dybeste Forstand det Overnaturlige? Netop
det, hvori alt det Naturlige grunder. „Jo mere af det
Overnaturlige, der er i et Menneskes Villie, desto mere
personlig er denne Villie; jo mere personlig Villien er,
desto mere vil den formaae at gjennemarbeide Individualitetens
Naturgrund, udvikle dens Evner og styrke
dens Anlæg, først og sidst dens Evighedsanlæg, at det
Forkrænkelige kan døe bort, og det Uforkrænkelige
leve.“ Det er Troessamvittighedens Vidnesbyrd. ---
Mon et Menneske med en saadan „Samvittighedstro“
ikke skulde foretrække at blive i Luthers gamle Kirke
og --- med Forbehold af Tilegnelsens personlige Frihed
--- bøie sig i Ærefrygt for dens Dogmer og Sakramenter,
istedenfor at søge Trøst og Raad i den rationalistiske
Eenhed af Tro og Viden?

Tilværelsen er et Under --- alles Sein ist Wunder;
det bliver indrømmet saa vel fra Videnskabens som
fra Religionens Side, at her er en Gaade, som ikke
kan løses. Nutidens Antagelse, at den sandselige
Natur med dens Phænomener og Love er et blindt
% --- p. 112 ---
\setcounter{footnote}{0}%
\opage{112}Produkt af en blind Nødvendighed, der ikke forstaaer
sig selv, er lige saa eensidig og for den personlige
Selvforstaaelse utilfredsstillende, som Middelalderens
Forestilling om det Overnaturliges vilkaarlige Indgriben
i Tingenes Sammenhæng er oprørende for Nutidens
videnskabelige Bevidsthed. Den gaadefulde Vexelvirkning
af Overnaturligt og Naturligt --- det perennerende
Under --- er en for den menneskelige Forstand
uigjennemtrængelig Fordoblingseenhed af Naturens
Midler og Aandens Maal, en Fordoblingseenhed, som
fra Videns Standpunkt maa bedømmes efter Nødvendighedslovene,
fra Troens Standpunkt efter Frihedslovene.
Saa vist som Videnskaben paa Forskningens, Kritikens
og Kulturens Vegne har Ret til at erklære ethvert nok
saa dybsindigt Forsøg, som vil forklare den overnaturlige
Villies Indgriben ved at oversee, misforstaae eller
spidsfindig omgaae Naturlovene, for huult og phantastisk,
lige saa vist har Religionen med sin Aabenbaringstro,
i Personlighedens, Humanitetens og Frihedens Navn,
Ret til at protestere imod ethvert Overgreb, som under
Foregivende af at hæve Selvbevidstheden op paa et
høiere Trin, gjør Selvets Væsen til Skin og vitterlig
krænker Aandslovene. Indenfor disse Grændser er
Maalestokken holdt --- til en Karakteristik af \emph{gamle}
og \emph{nye Propheter}.

\end{document}
